%============================================================
%  main_sp.tex
%  Tensor Formats in Banach Spaces:
%  Minimal Subspaces, Geometry, and Applications to Open Quantum Systems
%
%  Antonio Falcó — Universidad Cardenal Herrera-CEU
%
%  TARGET: Ergebnisse der Mathematik und ihrer Grenzgebiete
%          (3. Folge), Springer
%
%  DOCUMENT CLASS: SNmono.cls  (Springer Nature monograph class,
%                  v6.00, 11-Nov-2025)
%
%  COMPILATION:
%    pdflatex main_sp && bibtex main_sp && pdflatex main_sp && pdflatex main_sp
%  or:
%    latexmk -pdf main_sp
%============================================================

%------------------------------------------------------------
%  DOCUMENT CLASS
%  SNmono v6.00 valid options:
%    envcountchap  — number theorems per chapter  (e.g. Theorem 1.1)
%    envcountsect  — number theorems per section  (e.g. Theorem 1.1.1)
%  Options 'sechang' and 'openright' are NOT supported in v6.00.
%------------------------------------------------------------
\documentclass[envcountsect]{SNmono}

%------------------------------------------------------------
%  PACKAGES
%  SNmono v6.00 loads internally:
%    amsmath, amsthm, graphicx, hyperref, xcolor, inputenc, fontenc
%  We add only what is missing.
%------------------------------------------------------------

% Symbol fonts — NOT loaded by SNmono
\usepackage{amssymb}             % \mathbb, \mathfrak, etc.
\usepackage{mathrsfs}            % \mathscr
\usepackage{mathtools}           % \coloneqq, \xrightarrow, etc.

% Diagrams and trees
\usepackage{tikz}
\usepackage[all,cmtip]{xy}
\usepackage{synttree}

% Tables and lists
\usepackage{booktabs}
\usepackage{enumitem}

% Index
\usepackage{makeidx}
\makeindex

%------------------------------------------------------------
%  THEOREM ENVIRONMENTS
%  SNmono already defines via \spnewtheorem:
%    theorem, lemma, corollary, proposition, definition,
%    example, exercise, remark, note, case, conjecture,
%    question, solution, property
%  We only declare the ones NOT provided by SNmono.
%  Syntax: \spnewtheorem{name}[shared-counter]{Caption}{\headfont}{\bodyfont}
%------------------------------------------------------------
\spnewtheorem{assumption}[theorem]{Assumption}{\bfseries}{\upshape}
\spnewtheorem{problem}[theorem]{Open Problem}{\bfseries}{\upshape}
\spnewtheorem{notation}[theorem]{Notation}{\itshape}{\upshape}

%------------------------------------------------------------
%  UNIFIED NOTATION
%------------------------------------------------------------
%============================================================
%  notation.tex
%  Unified notation and macros for the monograph
%  "Tensor Formats in Banach Spaces"
%
%  Antonio Falcó, CEU Universities
%
%  Usage:  %============================================================
%  notation.tex
%  Unified notation and macros for the monograph
%  "Tensor Formats in Banach Spaces"
%
%  Antonio Falcó, CEU Universities
%
%  Usage:  %============================================================
%  notation.tex
%  Unified notation and macros for the monograph
%  "Tensor Formats in Banach Spaces"
%
%  Antonio Falcó, CEU Universities
%
%  Usage:  \input{notation}  in main.tex preamble
%
%  PHILOSOPHY:
%  (a) Scalars / indices         : lowercase Roman or Greek
%  (b) Vectors in component spaces: lowercase bold  \bv, \bu, ...
%  (c) Tensors                   : uppercase bold   \bV, \bU, ...
%  (d) Tensor spaces (algebraic) : \VD, \Valpha, ...
%  (e) Topological completions   : decorated with norm subscript
%  (f) Trees, partitions         : calligraphic or Roman
%============================================================

% ---- Fields ----------------------------------------
\newcommand{\RR}{\mathbb{R}}          % real numbers
\newcommand{\CC}{\mathbb{C}}          % complex numbers
\newcommand{\KK}{\mathbb{K}}          % generic field (R or C)
\newcommand{\NN}{\mathbb{N}}          % positive integers
\newcommand{\NNz}{\mathbb{N}_0}       % non-negative integers  N cup {0}
\newcommand{\ZZ}{\mathbb{Z}}          % integers

% ---- Dimension index set ---------------------------
% D = {1,...,d}
\newcommand{\DD}{D}                   % dimension index set
\newcommand{\dD}{d}                   % cardinality |D|

% ---- Component spaces ------------------------------
% V_j  (j in D),  V_alpha  (alpha subset of D)
\newcommand{\Vj}{V_j}
\newcommand{\Valpha}{V_\alpha}
\newcommand{\Vbeta}{V_\beta}
\newcommand{\Vnu}{V_\nu}

% ---- Algebraic tensor spaces -----------------------
% Convention: _a\bigotimes = algebraic tensor product
% \VD   = _a otimes_{j in D} V_j
% \Vbigalpha = _a otimes_{j in alpha} V_j
\newcommand{\VD}{\mathbf{V}_\DD}
\newcommand{\VDalg}{{}_a\bigotimes_{\alpha \in \DD} V_\alpha}
\newcommand{\Vbigalpha}{\mathbf{V}_\alpha}
\newcommand{\Vbigbeta}{\mathbf{V}_\beta}
\newcommand{\VPk}{\mathbf{V}_{\mathcal{P}_k}}  % space at level-k partition

% Algebraic tensor product symbols
\newcommand{\atens}{\otimes_a}         % v \atens w
\newcommand{\Atens}{\bigotimes_a}      % \Atens_{j} V_j

% ---- Topological tensor spaces (completions) -------
% Completion of V_D under ||\cdot||_D:  \VDnorm
\newcommand{\VDnorm}{\mathbf{V}_{D,\|\cdot\|_D}}
\newcommand{\Valphanorm}{\mathbf{V}_{\alpha,\|\cdot\|_\alpha}}
% Topological tensor product: _{\|\cdot\|}\bigotimes
\newcommand{\tnorm}[1]{{}_{\|\cdot\|_{#1}}\bigotimes}

% ---- Tensor elements -------------------------------
\newcommand{\bv}{\mathbf{v}}
\newcommand{\bu}{\mathbf{u}}
\newcommand{\bw}{\mathbf{w}}
\newcommand{\bx}{\mathbf{x}}
\newcommand{\bA}{\mathbf{A}}              % bold A (e.g. HS operator as tensor)
\newcommand{\bF}{\mathbf{F}}           % vector field in ODE
\newcommand{\bU}{\mathbf{U}}

% ---- Norms -----------------------------------------
\newcommand{\norm}[2]{\left\|#1\right\|_{#2}}     % generic: \norm{v}{D}
\newcommand{\normD}[1]{\left\|#1\right\|_{D}}     % norm in V_D
\newcommand{\norma}[1]{\left\|#1\right\|_{\alpha}}
\newcommand{\normj}[1]{\left\|#1\right\|_{j}}
\newcommand{\opnorm}[3]{\left\|#1\right\|_{#2 \leftarrow #3}}  % operator norm

% ---- Dual spaces -----------------------------------
\newcommand{\dual}[1]{#1^*}            % topological dual  X^*
\newcommand{\algdual}[1]{#1'}          % algebraic dual  L(V,K)

% ---- Linear maps -----------------------------------
\newcommand{\cL}{\mathcal{L}}          % bounded linear maps  L(X,Y)
\newcommand{\GL}[1]{\mathrm{GL}(#1)}   % invertible automorphisms
\newcommand{\Id}[1]{\mathrm{id}_{#1}}  % identity on #1
\newcommand{\Proj}[2]{P_{#1 \oplus #2}}% projection onto #1 along #2

% ---- Grassmannian and split subspaces --------------
\newcommand{\Grass}[2]{\mathbb{G}_{#1}(#2)}  % r-dimensional subspaces
\newcommand{\GrassB}[1]{\mathbb{G}(#1)}      % all split subspaces
\newcommand{\AD}[1]{\mathcal{AD}(#1)}        % admissible Tucker ranks

% ---- Minimal subspaces -----------------------------
% U^min_alpha(v)  =  minimal subspace of v in V_alpha
\newcommand{\Umin}[2]{U^{\min}_{#1}(#2)}      % \Umin{\alpha}{\bv}
\newcommand{\Wmin}[2]{W^{\min}_{#1}(#2)}      % complement W
\newcommand{\alphabrank}{r_\alpha}             % alpha-rank (fixed)

% ---- Tucker format ---------------------------------
% M_r(V_D)  =  tensors in Tucker format with fixed rank r
\newcommand{\Tucker}[2]{\mathfrak{M}_{#1}(#2)}
\newcommand{\TuckerV}{\mathfrak{M}_{\fr}(\VD)}
% Tangent subspace at v
\newcommand{\ZD}[1]{Z^{(D)}(#1)}

% ---- Tree-based format -----------------------------
\newcommand{\TD}{T_D}                          % dimension partition tree
\newcommand{\TDtuck}{T_D^{\mathrm{Tucker}}}   % Tucker tree (depth 1)
\newcommand{\TDtt}{T_D^{\mathrm{TT}}}         % Tensor Train tree
\newcommand{\TDht}{T_D^{\mathrm{HT}}}         % Hierarchical Tucker tree

\newcommand{\Sons}[1]{S(#1)}                   % sons of vertex alpha
\newcommand{\Leaves}{\mathcal{L}(\TD)}         % leaves of T_D
\newcommand{\lev}[1]{\mathrm{level}(#1)}       % level of vertex
\newcommand{\dep}[1]{\mathrm{depth}(#1)}       % depth of tree

\newcommand{\Pk}{\mathcal{P}_k(\TD)}           % level-k partition
\newcommand{\PkT}[2]{\mathcal{P}_{#1}(#2)}    % general: \PkT{k}{T_D}

% FT_r(V_D, T_D)  =  tree-based tensors with fixed rank r
\newcommand{\FT}[3]{\mathcal{FT}_{#1}(#2,#3)}
\newcommand{\FTv}{\mathcal{FT}_{\fr}(\VD,\TD)}
\newcommand{\FTleq}[3]{\mathcal{FT}_{\leq #1}(#2,#3)}

% Tree-based rank tuples
\newcommand{\fr}{\mathfrak{r}}                 % fixed rank tuple
\newcommand{\fs}{\mathfrak{s}}                 % another rank tuple
\newcommand{\ADTD}{\mathcal{AD}(\VD,\TD)}     % admissible tuples

% Transfer tensor at vertex alpha
\newcommand{\Calpha}{C^{(\alpha)}}
\newcommand{\CD}{C^{(D)}}
\newcommand{\Cmu}{C^{(\mu)}}
% Matricisation map
\newcommand{\Mmat}[1]{\mathcal{M}_{#1}}

% ---- Banach manifold objects -----------------------
\newcommand{\chartmap}[1]{\xi_{#1}}            % chart map at v
\newcommand{\Tv}[2]{T_{#1}#2}                  % tangent space T_v M
\newcommand{\Emodel}[1]{\mathbf{E}_{#1}}       % model Banach space

% ---- Dirac-Frenkel objects -------------------------
\newcommand{\Pmetric}[1]{\mathcal{P}_{#1}}    % metric projection at v
\newcommand{\Fvf}{\mathbf{F}}                  % vector field

% ---- Norm conditions -------------------------------
\newcommand{\injnorm}{\|\cdot\|_\varepsilon}   % injective norm
\newcommand{\projnorm}{\|\cdot\|_\pi}          % projective norm
\newcommand{\condINJ}{\textup{(Inj)}}          % condition: >= injective
\newcommand{\condCROSS}{\textup{(Cross)}}      % condition: crossnorm

% ---- Miscellaneous ---------------------------------
\newcommand{\spann}{\mathrm{span}}
\newcommand{\rank}{\mathrm{rank}}
\newcommand{\depth}{\mathrm{depth}}
\newcommand{\level}{\mathrm{level}}
\newcommand{\dist}{\mathrm{dist}}
\newcommand{\cl}[1]{\overline{#1}}             % closure
\newcommand{\card}[1]{\#(#1)}                  % cardinality
\newcommand{\ac}{\alpha^c}                     % complement D\alpha
\newcommand{\pset}[1]{2^{#1}}                  % power set 2^D

% ---- Cross-reference shorthands --------------------
\newcommand{\refthm}[1]{Theorem~\ref{#1}}
\newcommand{\reflem}[1]{Lemma~\ref{#1}}
\newcommand{\refprop}[1]{Proposition~\ref{#1}}
\newcommand{\refdef}[1]{Definition~\ref{#1}}
\newcommand{\refsec}[1]{Section~\ref{#1}}
\newcommand{\refchap}[1]{Chapter~\ref{#1}}
\providecommand{\refeq}[1]{(\ref{#1})}

% Chart domain (added for Ch.8-9)
\newcommand{\chart}{\mathcal{U}}   % chart domain \chart(\bv) -> \mathcal{U}(\bv)

% Math operators (added for Ch.10)
\DeclareMathOperator*{\argmin}{arg\,min}
\DeclareMathOperator*{\argmax}{arg\,max}

% ============================================================
% Part V: Quantum States (Appendix E + Chapters 12--13)
% ============================================================

% ---- Operator ideals on Hilbert spaces -------------------
\newcommand{\cT}[1]{\mathcal{T}_{#1}}      % \cT{1}(H) = trace-class
                                             % \cT{2}(H) = Hilbert-Schmidt
\newcommand{\cTsa}[1]{\mathcal{T}_{#1}^{sa}}% \cTsa{1}(H) = self-adjoint trace-class
\newcommand{\cK}{\mathcal{K}}               % compact operators \cK(H)

% ---- Trace and singular values ---------------------------
\newcommand{\Tr}{\mathrm{Tr}}              % trace
\newcommand{\TrS}[1]{\mathrm{Tr}_{#1}}    % partial trace \TrS{2}(\rho)
\newcommand{\svals}[1]{s_{#1}}             % singular values s_k(A)
\newcommand{\supp}{\mathrm{supp}}          % support of an operator

% ---- Quantum state spaces --------------------------------
\newcommand{\Dstate}[1]{\mathcal{D}(#1)}   % all density operators
\newcommand{\Drank}[2]{\mathcal{D}_{#1}(#2)}% rank-r stratum \Drank{r}{H}
\newcommand{\Dpos}[1]{\mathcal{D}_+(#1)}   % positive core \Dpos{S}

% ---- Rank-one operators ----------------------------------
% u \roten v  =  the rank-one operator w \mapsto <v,w> u
\newcommand{\roten}{\otimes_{\overline{\phantom{v}}}}
% Usage: u \roten v  (read: u tensor conjugate-v)
% Displayed as: u \otimes \overline{v}  in the text

% ---- Hilbert tensor product ------------------------------
\newcommand{\Htens}{\otimes_H}             % Hilbert tensor product
\newcommand{\Hbar}[1]{\overline{#1}}       % conjugate Hilbert space

% ---- Projectors ------------------------------------------
\newcommand{\PS}[1]{P_{#1}}               % orthogonal projection onto S
\newcommand{\Prho}{\hat{P}_\rho}          % tangent projector at rho (Ch.13)

% ---- Lindblad / GKSL objects ----------------------------
\newcommand{\Lind}{\mathcal{L}}            % Lindblad generator
\newcommand{\Stiefel}[2]{\mathrm{St}(#1,#2)}% Stiefel manifold St(r,H)

% ============================================================
% Worked examples (Ch04, Ch10)
% ============================================================

% ---- Parametric PDE examples ----------------------------
\newcommand{\bmy}{\mathbf{y}}              % parameter vector y in Gamma
\newcommand{\bmx}{\mathbf{x}}              % spatial vector x (multi-dim)

% ---- Hamiltonian (avoids clash with Hilbert space H) -----
\newcommand{\Ham}{\mathcal{H}}             % Hamiltonian operator \Ham
\newcommand{\Hamone}{\mathcal{H}_1}        % single-particle Hamiltonian 1
\newcommand{\Hamtwo}{\mathcal{H}_2}        % single-particle Hamiltonian 2

% ============================================================
% Probability and statistics (Ch11)
% ============================================================
\newcommand{\EE}{\mathbb{E}}               % expectation operator
  in main.tex preamble
%
%  PHILOSOPHY:
%  (a) Scalars / indices         : lowercase Roman or Greek
%  (b) Vectors in component spaces: lowercase bold  \bv, \bu, ...
%  (c) Tensors                   : uppercase bold   \bV, \bU, ...
%  (d) Tensor spaces (algebraic) : \VD, \Valpha, ...
%  (e) Topological completions   : decorated with norm subscript
%  (f) Trees, partitions         : calligraphic or Roman
%============================================================

% ---- Fields ----------------------------------------
\newcommand{\RR}{\mathbb{R}}          % real numbers
\newcommand{\CC}{\mathbb{C}}          % complex numbers
\newcommand{\KK}{\mathbb{K}}          % generic field (R or C)
\newcommand{\NN}{\mathbb{N}}          % positive integers
\newcommand{\NNz}{\mathbb{N}_0}       % non-negative integers  N cup {0}
\newcommand{\ZZ}{\mathbb{Z}}          % integers

% ---- Dimension index set ---------------------------
% D = {1,...,d}
\newcommand{\DD}{D}                   % dimension index set
\newcommand{\dD}{d}                   % cardinality |D|

% ---- Component spaces ------------------------------
% V_j  (j in D),  V_alpha  (alpha subset of D)
\newcommand{\Vj}{V_j}
\newcommand{\Valpha}{V_\alpha}
\newcommand{\Vbeta}{V_\beta}
\newcommand{\Vnu}{V_\nu}

% ---- Algebraic tensor spaces -----------------------
% Convention: _a\bigotimes = algebraic tensor product
% \VD   = _a otimes_{j in D} V_j
% \Vbigalpha = _a otimes_{j in alpha} V_j
\newcommand{\VD}{\mathbf{V}_\DD}
\newcommand{\VDalg}{{}_a\bigotimes_{\alpha \in \DD} V_\alpha}
\newcommand{\Vbigalpha}{\mathbf{V}_\alpha}
\newcommand{\Vbigbeta}{\mathbf{V}_\beta}
\newcommand{\VPk}{\mathbf{V}_{\mathcal{P}_k}}  % space at level-k partition

% Algebraic tensor product symbols
\newcommand{\atens}{\otimes_a}         % v \atens w
\newcommand{\Atens}{\bigotimes_a}      % \Atens_{j} V_j

% ---- Topological tensor spaces (completions) -------
% Completion of V_D under ||\cdot||_D:  \VDnorm
\newcommand{\VDnorm}{\mathbf{V}_{D,\|\cdot\|_D}}
\newcommand{\Valphanorm}{\mathbf{V}_{\alpha,\|\cdot\|_\alpha}}
% Topological tensor product: _{\|\cdot\|}\bigotimes
\newcommand{\tnorm}[1]{{}_{\|\cdot\|_{#1}}\bigotimes}

% ---- Tensor elements -------------------------------
\newcommand{\bv}{\mathbf{v}}
\newcommand{\bu}{\mathbf{u}}
\newcommand{\bw}{\mathbf{w}}
\newcommand{\bx}{\mathbf{x}}
\newcommand{\bA}{\mathbf{A}}              % bold A (e.g. HS operator as tensor)
\newcommand{\bF}{\mathbf{F}}           % vector field in ODE
\newcommand{\bU}{\mathbf{U}}

% ---- Norms -----------------------------------------
\newcommand{\norm}[2]{\left\|#1\right\|_{#2}}     % generic: \norm{v}{D}
\newcommand{\normD}[1]{\left\|#1\right\|_{D}}     % norm in V_D
\newcommand{\norma}[1]{\left\|#1\right\|_{\alpha}}
\newcommand{\normj}[1]{\left\|#1\right\|_{j}}
\newcommand{\opnorm}[3]{\left\|#1\right\|_{#2 \leftarrow #3}}  % operator norm

% ---- Dual spaces -----------------------------------
\newcommand{\dual}[1]{#1^*}            % topological dual  X^*
\newcommand{\algdual}[1]{#1'}          % algebraic dual  L(V,K)

% ---- Linear maps -----------------------------------
\newcommand{\cL}{\mathcal{L}}          % bounded linear maps  L(X,Y)
\newcommand{\GL}[1]{\mathrm{GL}(#1)}   % invertible automorphisms
\newcommand{\Id}[1]{\mathrm{id}_{#1}}  % identity on #1
\newcommand{\Proj}[2]{P_{#1 \oplus #2}}% projection onto #1 along #2

% ---- Grassmannian and split subspaces --------------
\newcommand{\Grass}[2]{\mathbb{G}_{#1}(#2)}  % r-dimensional subspaces
\newcommand{\GrassB}[1]{\mathbb{G}(#1)}      % all split subspaces
\newcommand{\AD}[1]{\mathcal{AD}(#1)}        % admissible Tucker ranks

% ---- Minimal subspaces -----------------------------
% U^min_alpha(v)  =  minimal subspace of v in V_alpha
\newcommand{\Umin}[2]{U^{\min}_{#1}(#2)}      % \Umin{\alpha}{\bv}
\newcommand{\Wmin}[2]{W^{\min}_{#1}(#2)}      % complement W
\newcommand{\alphabrank}{r_\alpha}             % alpha-rank (fixed)

% ---- Tucker format ---------------------------------
% M_r(V_D)  =  tensors in Tucker format with fixed rank r
\newcommand{\Tucker}[2]{\mathfrak{M}_{#1}(#2)}
\newcommand{\TuckerV}{\mathfrak{M}_{\fr}(\VD)}
% Tangent subspace at v
\newcommand{\ZD}[1]{Z^{(D)}(#1)}

% ---- Tree-based format -----------------------------
\newcommand{\TD}{T_D}                          % dimension partition tree
\newcommand{\TDtuck}{T_D^{\mathrm{Tucker}}}   % Tucker tree (depth 1)
\newcommand{\TDtt}{T_D^{\mathrm{TT}}}         % Tensor Train tree
\newcommand{\TDht}{T_D^{\mathrm{HT}}}         % Hierarchical Tucker tree

\newcommand{\Sons}[1]{S(#1)}                   % sons of vertex alpha
\newcommand{\Leaves}{\mathcal{L}(\TD)}         % leaves of T_D
\newcommand{\lev}[1]{\mathrm{level}(#1)}       % level of vertex
\newcommand{\dep}[1]{\mathrm{depth}(#1)}       % depth of tree

\newcommand{\Pk}{\mathcal{P}_k(\TD)}           % level-k partition
\newcommand{\PkT}[2]{\mathcal{P}_{#1}(#2)}    % general: \PkT{k}{T_D}

% FT_r(V_D, T_D)  =  tree-based tensors with fixed rank r
\newcommand{\FT}[3]{\mathcal{FT}_{#1}(#2,#3)}
\newcommand{\FTv}{\mathcal{FT}_{\fr}(\VD,\TD)}
\newcommand{\FTleq}[3]{\mathcal{FT}_{\leq #1}(#2,#3)}

% Tree-based rank tuples
\newcommand{\fr}{\mathfrak{r}}                 % fixed rank tuple
\newcommand{\fs}{\mathfrak{s}}                 % another rank tuple
\newcommand{\ADTD}{\mathcal{AD}(\VD,\TD)}     % admissible tuples

% Transfer tensor at vertex alpha
\newcommand{\Calpha}{C^{(\alpha)}}
\newcommand{\CD}{C^{(D)}}
\newcommand{\Cmu}{C^{(\mu)}}
% Matricisation map
\newcommand{\Mmat}[1]{\mathcal{M}_{#1}}

% ---- Banach manifold objects -----------------------
\newcommand{\chartmap}[1]{\xi_{#1}}            % chart map at v
\newcommand{\Tv}[2]{T_{#1}#2}                  % tangent space T_v M
\newcommand{\Emodel}[1]{\mathbf{E}_{#1}}       % model Banach space

% ---- Dirac-Frenkel objects -------------------------
\newcommand{\Pmetric}[1]{\mathcal{P}_{#1}}    % metric projection at v
\newcommand{\Fvf}{\mathbf{F}}                  % vector field

% ---- Norm conditions -------------------------------
\newcommand{\injnorm}{\|\cdot\|_\varepsilon}   % injective norm
\newcommand{\projnorm}{\|\cdot\|_\pi}          % projective norm
\newcommand{\condINJ}{\textup{(Inj)}}          % condition: >= injective
\newcommand{\condCROSS}{\textup{(Cross)}}      % condition: crossnorm

% ---- Miscellaneous ---------------------------------
\newcommand{\spann}{\mathrm{span}}
\newcommand{\rank}{\mathrm{rank}}
\newcommand{\depth}{\mathrm{depth}}
\newcommand{\level}{\mathrm{level}}
\newcommand{\dist}{\mathrm{dist}}
\newcommand{\cl}[1]{\overline{#1}}             % closure
\newcommand{\card}[1]{\#(#1)}                  % cardinality
\newcommand{\ac}{\alpha^c}                     % complement D\alpha
\newcommand{\pset}[1]{2^{#1}}                  % power set 2^D

% ---- Cross-reference shorthands --------------------
\newcommand{\refthm}[1]{Theorem~\ref{#1}}
\newcommand{\reflem}[1]{Lemma~\ref{#1}}
\newcommand{\refprop}[1]{Proposition~\ref{#1}}
\newcommand{\refdef}[1]{Definition~\ref{#1}}
\newcommand{\refsec}[1]{Section~\ref{#1}}
\newcommand{\refchap}[1]{Chapter~\ref{#1}}
\providecommand{\refeq}[1]{(\ref{#1})}

% Chart domain (added for Ch.8-9)
\newcommand{\chart}{\mathcal{U}}   % chart domain \chart(\bv) -> \mathcal{U}(\bv)

% Math operators (added for Ch.10)
\DeclareMathOperator*{\argmin}{arg\,min}
\DeclareMathOperator*{\argmax}{arg\,max}

% ============================================================
% Part V: Quantum States (Appendix E + Chapters 12--13)
% ============================================================

% ---- Operator ideals on Hilbert spaces -------------------
\newcommand{\cT}[1]{\mathcal{T}_{#1}}      % \cT{1}(H) = trace-class
                                             % \cT{2}(H) = Hilbert-Schmidt
\newcommand{\cTsa}[1]{\mathcal{T}_{#1}^{sa}}% \cTsa{1}(H) = self-adjoint trace-class
\newcommand{\cK}{\mathcal{K}}               % compact operators \cK(H)

% ---- Trace and singular values ---------------------------
\newcommand{\Tr}{\mathrm{Tr}}              % trace
\newcommand{\TrS}[1]{\mathrm{Tr}_{#1}}    % partial trace \TrS{2}(\rho)
\newcommand{\svals}[1]{s_{#1}}             % singular values s_k(A)
\newcommand{\supp}{\mathrm{supp}}          % support of an operator

% ---- Quantum state spaces --------------------------------
\newcommand{\Dstate}[1]{\mathcal{D}(#1)}   % all density operators
\newcommand{\Drank}[2]{\mathcal{D}_{#1}(#2)}% rank-r stratum \Drank{r}{H}
\newcommand{\Dpos}[1]{\mathcal{D}_+(#1)}   % positive core \Dpos{S}

% ---- Rank-one operators ----------------------------------
% u \roten v  =  the rank-one operator w \mapsto <v,w> u
\newcommand{\roten}{\otimes_{\overline{\phantom{v}}}}
% Usage: u \roten v  (read: u tensor conjugate-v)
% Displayed as: u \otimes \overline{v}  in the text

% ---- Hilbert tensor product ------------------------------
\newcommand{\Htens}{\otimes_H}             % Hilbert tensor product
\newcommand{\Hbar}[1]{\overline{#1}}       % conjugate Hilbert space

% ---- Projectors ------------------------------------------
\newcommand{\PS}[1]{P_{#1}}               % orthogonal projection onto S
\newcommand{\Prho}{\hat{P}_\rho}          % tangent projector at rho (Ch.13)

% ---- Lindblad / GKSL objects ----------------------------
\newcommand{\Lind}{\mathcal{L}}            % Lindblad generator
\newcommand{\Stiefel}[2]{\mathrm{St}(#1,#2)}% Stiefel manifold St(r,H)

% ============================================================
% Worked examples (Ch04, Ch10)
% ============================================================

% ---- Parametric PDE examples ----------------------------
\newcommand{\bmy}{\mathbf{y}}              % parameter vector y in Gamma
\newcommand{\bmx}{\mathbf{x}}              % spatial vector x (multi-dim)

% ---- Hamiltonian (avoids clash with Hilbert space H) -----
\newcommand{\Ham}{\mathcal{H}}             % Hamiltonian operator \Ham
\newcommand{\Hamone}{\mathcal{H}_1}        % single-particle Hamiltonian 1
\newcommand{\Hamtwo}{\mathcal{H}_2}        % single-particle Hamiltonian 2

% ============================================================
% Probability and statistics (Ch11)
% ============================================================
\newcommand{\EE}{\mathbb{E}}               % expectation operator
  in main.tex preamble
%
%  PHILOSOPHY:
%  (a) Scalars / indices         : lowercase Roman or Greek
%  (b) Vectors in component spaces: lowercase bold  \bv, \bu, ...
%  (c) Tensors                   : uppercase bold   \bV, \bU, ...
%  (d) Tensor spaces (algebraic) : \VD, \Valpha, ...
%  (e) Topological completions   : decorated with norm subscript
%  (f) Trees, partitions         : calligraphic or Roman
%============================================================

% ---- Fields ----------------------------------------
\newcommand{\RR}{\mathbb{R}}          % real numbers
\newcommand{\CC}{\mathbb{C}}          % complex numbers
\newcommand{\KK}{\mathbb{K}}          % generic field (R or C)
\newcommand{\NN}{\mathbb{N}}          % positive integers
\newcommand{\NNz}{\mathbb{N}_0}       % non-negative integers  N cup {0}
\newcommand{\ZZ}{\mathbb{Z}}          % integers

% ---- Dimension index set ---------------------------
% D = {1,...,d}
\newcommand{\DD}{D}                   % dimension index set
\newcommand{\dD}{d}                   % cardinality |D|

% ---- Component spaces ------------------------------
% V_j  (j in D),  V_alpha  (alpha subset of D)
\newcommand{\Vj}{V_j}
\newcommand{\Valpha}{V_\alpha}
\newcommand{\Vbeta}{V_\beta}
\newcommand{\Vnu}{V_\nu}

% ---- Algebraic tensor spaces -----------------------
% Convention: _a\bigotimes = algebraic tensor product
% \VD   = _a otimes_{j in D} V_j
% \Vbigalpha = _a otimes_{j in alpha} V_j
\newcommand{\VD}{\mathbf{V}_\DD}
\newcommand{\VDalg}{{}_a\bigotimes_{\alpha \in \DD} V_\alpha}
\newcommand{\Vbigalpha}{\mathbf{V}_\alpha}
\newcommand{\Vbigbeta}{\mathbf{V}_\beta}
\newcommand{\VPk}{\mathbf{V}_{\mathcal{P}_k}}  % space at level-k partition

% Algebraic tensor product symbols
\newcommand{\atens}{\otimes_a}         % v \atens w
\newcommand{\Atens}{\bigotimes_a}      % \Atens_{j} V_j

% ---- Topological tensor spaces (completions) -------
% Completion of V_D under ||\cdot||_D:  \VDnorm
\newcommand{\VDnorm}{\mathbf{V}_{D,\|\cdot\|_D}}
\newcommand{\Valphanorm}{\mathbf{V}_{\alpha,\|\cdot\|_\alpha}}
% Topological tensor product: _{\|\cdot\|}\bigotimes
\newcommand{\tnorm}[1]{{}_{\|\cdot\|_{#1}}\bigotimes}

% ---- Tensor elements -------------------------------
\newcommand{\bv}{\mathbf{v}}
\newcommand{\bu}{\mathbf{u}}
\newcommand{\bw}{\mathbf{w}}
\newcommand{\bx}{\mathbf{x}}
\newcommand{\bA}{\mathbf{A}}              % bold A (e.g. HS operator as tensor)
\newcommand{\bF}{\mathbf{F}}           % vector field in ODE
\newcommand{\bU}{\mathbf{U}}

% ---- Norms -----------------------------------------
\newcommand{\norm}[2]{\left\|#1\right\|_{#2}}     % generic: \norm{v}{D}
\newcommand{\normD}[1]{\left\|#1\right\|_{D}}     % norm in V_D
\newcommand{\norma}[1]{\left\|#1\right\|_{\alpha}}
\newcommand{\normj}[1]{\left\|#1\right\|_{j}}
\newcommand{\opnorm}[3]{\left\|#1\right\|_{#2 \leftarrow #3}}  % operator norm

% ---- Dual spaces -----------------------------------
\newcommand{\dual}[1]{#1^*}            % topological dual  X^*
\newcommand{\algdual}[1]{#1'}          % algebraic dual  L(V,K)

% ---- Linear maps -----------------------------------
\newcommand{\cL}{\mathcal{L}}          % bounded linear maps  L(X,Y)
\newcommand{\GL}[1]{\mathrm{GL}(#1)}   % invertible automorphisms
\newcommand{\Id}[1]{\mathrm{id}_{#1}}  % identity on #1
\newcommand{\Proj}[2]{P_{#1 \oplus #2}}% projection onto #1 along #2

% ---- Grassmannian and split subspaces --------------
\newcommand{\Grass}[2]{\mathbb{G}_{#1}(#2)}  % r-dimensional subspaces
\newcommand{\GrassB}[1]{\mathbb{G}(#1)}      % all split subspaces
\newcommand{\AD}[1]{\mathcal{AD}(#1)}        % admissible Tucker ranks

% ---- Minimal subspaces -----------------------------
% U^min_alpha(v)  =  minimal subspace of v in V_alpha
\newcommand{\Umin}[2]{U^{\min}_{#1}(#2)}      % \Umin{\alpha}{\bv}
\newcommand{\Wmin}[2]{W^{\min}_{#1}(#2)}      % complement W
\newcommand{\alphabrank}{r_\alpha}             % alpha-rank (fixed)

% ---- Tucker format ---------------------------------
% M_r(V_D)  =  tensors in Tucker format with fixed rank r
\newcommand{\Tucker}[2]{\mathfrak{M}_{#1}(#2)}
\newcommand{\TuckerV}{\mathfrak{M}_{\fr}(\VD)}
% Tangent subspace at v
\newcommand{\ZD}[1]{Z^{(D)}(#1)}

% ---- Tree-based format -----------------------------
\newcommand{\TD}{T_D}                          % dimension partition tree
\newcommand{\TDtuck}{T_D^{\mathrm{Tucker}}}   % Tucker tree (depth 1)
\newcommand{\TDtt}{T_D^{\mathrm{TT}}}         % Tensor Train tree
\newcommand{\TDht}{T_D^{\mathrm{HT}}}         % Hierarchical Tucker tree

\newcommand{\Sons}[1]{S(#1)}                   % sons of vertex alpha
\newcommand{\Leaves}{\mathcal{L}(\TD)}         % leaves of T_D
\newcommand{\lev}[1]{\mathrm{level}(#1)}       % level of vertex
\newcommand{\dep}[1]{\mathrm{depth}(#1)}       % depth of tree

\newcommand{\Pk}{\mathcal{P}_k(\TD)}           % level-k partition
\newcommand{\PkT}[2]{\mathcal{P}_{#1}(#2)}    % general: \PkT{k}{T_D}

% FT_r(V_D, T_D)  =  tree-based tensors with fixed rank r
\newcommand{\FT}[3]{\mathcal{FT}_{#1}(#2,#3)}
\newcommand{\FTv}{\mathcal{FT}_{\fr}(\VD,\TD)}
\newcommand{\FTleq}[3]{\mathcal{FT}_{\leq #1}(#2,#3)}

% Tree-based rank tuples
\newcommand{\fr}{\mathfrak{r}}                 % fixed rank tuple
\newcommand{\fs}{\mathfrak{s}}                 % another rank tuple
\newcommand{\ADTD}{\mathcal{AD}(\VD,\TD)}     % admissible tuples

% Transfer tensor at vertex alpha
\newcommand{\Calpha}{C^{(\alpha)}}
\newcommand{\CD}{C^{(D)}}
\newcommand{\Cmu}{C^{(\mu)}}
% Matricisation map
\newcommand{\Mmat}[1]{\mathcal{M}_{#1}}

% ---- Banach manifold objects -----------------------
\newcommand{\chartmap}[1]{\xi_{#1}}            % chart map at v
\newcommand{\Tv}[2]{T_{#1}#2}                  % tangent space T_v M
\newcommand{\Emodel}[1]{\mathbf{E}_{#1}}       % model Banach space

% ---- Dirac-Frenkel objects -------------------------
\newcommand{\Pmetric}[1]{\mathcal{P}_{#1}}    % metric projection at v
\newcommand{\Fvf}{\mathbf{F}}                  % vector field

% ---- Norm conditions -------------------------------
\newcommand{\injnorm}{\|\cdot\|_\varepsilon}   % injective norm
\newcommand{\projnorm}{\|\cdot\|_\pi}          % projective norm
\newcommand{\condINJ}{\textup{(Inj)}}          % condition: >= injective
\newcommand{\condCROSS}{\textup{(Cross)}}      % condition: crossnorm

% ---- Miscellaneous ---------------------------------
\newcommand{\spann}{\mathrm{span}}
\newcommand{\rank}{\mathrm{rank}}
\newcommand{\depth}{\mathrm{depth}}
\newcommand{\level}{\mathrm{level}}
\newcommand{\dist}{\mathrm{dist}}
\newcommand{\cl}[1]{\overline{#1}}             % closure
\newcommand{\card}[1]{\#(#1)}                  % cardinality
\newcommand{\ac}{\alpha^c}                     % complement D\alpha
\newcommand{\pset}[1]{2^{#1}}                  % power set 2^D

% ---- Cross-reference shorthands --------------------
\newcommand{\refthm}[1]{Theorem~\ref{#1}}
\newcommand{\reflem}[1]{Lemma~\ref{#1}}
\newcommand{\refprop}[1]{Proposition~\ref{#1}}
\newcommand{\refdef}[1]{Definition~\ref{#1}}
\newcommand{\refsec}[1]{Section~\ref{#1}}
\newcommand{\refchap}[1]{Chapter~\ref{#1}}
\providecommand{\refeq}[1]{(\ref{#1})}

% Chart domain (added for Ch.8-9)
\newcommand{\chart}{\mathcal{U}}   % chart domain \chart(\bv) -> \mathcal{U}(\bv)

% Math operators (added for Ch.10)
\DeclareMathOperator*{\argmin}{arg\,min}
\DeclareMathOperator*{\argmax}{arg\,max}

% ============================================================
% Part V: Quantum States (Appendix E + Chapters 12--13)
% ============================================================

% ---- Operator ideals on Hilbert spaces -------------------
\newcommand{\cT}[1]{\mathcal{T}_{#1}}      % \cT{1}(H) = trace-class
                                             % \cT{2}(H) = Hilbert-Schmidt
\newcommand{\cTsa}[1]{\mathcal{T}_{#1}^{sa}}% \cTsa{1}(H) = self-adjoint trace-class
\newcommand{\cK}{\mathcal{K}}               % compact operators \cK(H)

% ---- Trace and singular values ---------------------------
\newcommand{\Tr}{\mathrm{Tr}}              % trace
\newcommand{\TrS}[1]{\mathrm{Tr}_{#1}}    % partial trace \TrS{2}(\rho)
\newcommand{\svals}[1]{s_{#1}}             % singular values s_k(A)
\newcommand{\supp}{\mathrm{supp}}          % support of an operator

% ---- Quantum state spaces --------------------------------
\newcommand{\Dstate}[1]{\mathcal{D}(#1)}   % all density operators
\newcommand{\Drank}[2]{\mathcal{D}_{#1}(#2)}% rank-r stratum \Drank{r}{H}
\newcommand{\Dpos}[1]{\mathcal{D}_+(#1)}   % positive core \Dpos{S}

% ---- Rank-one operators ----------------------------------
% u \roten v  =  the rank-one operator w \mapsto <v,w> u
\newcommand{\roten}{\otimes_{\overline{\phantom{v}}}}
% Usage: u \roten v  (read: u tensor conjugate-v)
% Displayed as: u \otimes \overline{v}  in the text

% ---- Hilbert tensor product ------------------------------
\newcommand{\Htens}{\otimes_H}             % Hilbert tensor product
\newcommand{\Hbar}[1]{\overline{#1}}       % conjugate Hilbert space

% ---- Projectors ------------------------------------------
\newcommand{\PS}[1]{P_{#1}}               % orthogonal projection onto S
\newcommand{\Prho}{\hat{P}_\rho}          % tangent projector at rho (Ch.13)

% ---- Lindblad / GKSL objects ----------------------------
\newcommand{\Lind}{\mathcal{L}}            % Lindblad generator
\newcommand{\Stiefel}[2]{\mathrm{St}(#1,#2)}% Stiefel manifold St(r,H)

% ============================================================
% Worked examples (Ch04, Ch10)
% ============================================================

% ---- Parametric PDE examples ----------------------------
\newcommand{\bmy}{\mathbf{y}}              % parameter vector y in Gamma
\newcommand{\bmx}{\mathbf{x}}              % spatial vector x (multi-dim)

% ---- Hamiltonian (avoids clash with Hilbert space H) -----
\newcommand{\Ham}{\mathcal{H}}             % Hamiltonian operator \Ham
\newcommand{\Hamone}{\mathcal{H}_1}        % single-particle Hamiltonian 1
\newcommand{\Hamtwo}{\mathcal{H}_2}        % single-particle Hamiltonian 2

% ============================================================
% Probability and statistics (Ch11)
% ============================================================
\newcommand{\EE}{\mathbb{E}}               % expectation operator


%------------------------------------------------------------
%  CUSTOM COMMANDS  (not in notation.tex)
%------------------------------------------------------------
\newcommand{\ie}{i.e.,\ }
\newcommand{\eg}{e.g.,\ }
\newcommand{\cf}{cf.\ }
\newcommand{\Norm}[1]{\left\|#1\right\|}
\newcommand{\Hitchcock}{\cite{Hitchcock1927}}

%------------------------------------------------------------
%  TOC DEPTH
%------------------------------------------------------------
\setcounter{tocdepth}{2}

%------------------------------------------------------------
%  AUTHOR / TITLE METADATA
%------------------------------------------------------------
\title{Tensor Formats in Banach Spaces}
\subtitle{Minimal Subspaces, Geometry, and Applications to Open Quantum Systems}
\author{Antonio Falcó}
\institute{%
  Antonio Falcó \at
  ESI International Chair@CEU-UCH,
  Departamento de Matemáticas, Física y Ciencias Tecnológicas,
  Universidad Cardenal Herrera-CEU, CEU Universities,
  San Bartolomé 55, 46115 Alfara del Patriarca (Valencia), Spain\\
  \email{afalco@uchceu.es}
}

%============================================================
%  DOCUMENT
%============================================================
\begin{document}

%------------------------------------------------------------
%  FRONT MATTER
%------------------------------------------------------------
\frontmatter

\maketitle

\tableofcontents

%------------------------------------------------------------
%  PREFACE
%------------------------------------------------------------
%============================================================
%  ch00_preface.tex  —  Preface
%  Tensor Formats in Banach Spaces
%  Antonio Falcó
%============================================================

\chapter*{Preface}
\addcontentsline{toc}{chapter}{Preface}
\label{chap:preface}

% TODO: Write preface (~8 pp.)
%
% Outline:
%  1. Motivation and context
%  2. The four source papers and how they fit together
%  3. What this monograph adds beyond the articles
%  4. Reader's guide
%  5. Acknowledgements


%------------------------------------------------------------
%  NOTATION CHAPTER
%------------------------------------------------------------
\chapter*{Notation and Conventions}
\addcontentsline{toc}{chapter}{Notation and Conventions}
\label{chap:notation}
See the file \texttt{notation.tex} for the complete list of macros
and symbols used throughout.

%------------------------------------------------------------
%  MAIN MATTER
%------------------------------------------------------------
\mainmatter

%%  PART I: ALGEBRAIC AND TOPOLOGICAL FOUNDATIONS
\part{Algebraic and Topological Foundations}
\label{part:foundations}

%============================================================
%  ch01_introduction.tex  —  Introduction and Motivation
%  TODO: Write content
%============================================================

\chapter{Introduction and Motivation}
\label{chap:intro}

\section*{To be written}

This chapter is a placeholder. Content will be added in subsequent
revisions.


%============================================================
%  ch02_banach.tex
%  Chapter 2: Banach Space Prerequisites
%
%  Tensor Formats in Banach Spaces
%  Antonio Falcó, CEU Universities
%
%  Estimated length: ~30 pp.
%============================================================

\chapter{Banach Space Prerequisites}
\label{chap:banach}

This chapter collects the functional-analytic and differential-geometric
background required throughout the monograph. The reader familiar with
Banach space theory and infinite-dimensional differential geometry may
proceed directly to Chapter~\ref{chap:algtensor}, referring back as needed.
All vector spaces are over the field $\KK$, which is either $\RR$ or $\CC$;
unless stated otherwise, $\KK = \RR$.

%------------------------------------------------------------
\section{Normed spaces and Banach spaces}
\label{sec:banach:basics}
%------------------------------------------------------------

We begin by fixing notation and recalling the central notions of normed and
Banach space theory that will be used throughout.

\subsection*{Basic definitions}

Let $X$ be a vector space over $\KK$. A \emph{norm} on $X$ is a map
$\|\cdot\|: X \to [0,\infty)$ satisfying:
\begin{enumerate}
\item[(N1)] $\|x\| = 0$ if and only if $x = 0$,
\item[(N2)] $\|\lambda x\| = |\lambda|\,\|x\|$ for all $\lambda \in \KK$,
    $x \in X$,
\item[(N3)] $\|x + y\| \leq \|x\| + \|y\|$ for all $x, y \in X$.
\end{enumerate}
The pair $(X, \|\cdot\|)$ is called a \emph{normed space}. A sequence
$(x_n)_{n \in \NN}$ in $X$ is a \emph{Cauchy sequence} if
$\|x_n - x_m\| \to 0$ as $n, m \to \infty$. A normed space in which every
Cauchy sequence converges is called a \emph{Banach space}.

Two norms $\|\cdot\|$ and $\|\cdot\|'$ on $X$ are \emph{equivalent} if
there exist constants $c, C > 0$ such that
$c\|x\| \leq \|x\|' \leq C\|x\|$ for all $x \in X$.
Equivalent norms define the same topology, and hence the same Cauchy
sequences, so they produce the same Banach space structure.

\begin{example}
\label{ex:banach:lp}
Let $I \subset \RR$ be a measurable set and $1 \leq p < \infty$. The
Lebesgue space $L^p(I)$, equipped with the norm
$\|f\|_{L^p} = \bigl(\int_I |f(x)|^p\,dx\bigr)^{1/p}$,
is a Banach space. For $p = \infty$ one uses the essential supremum norm.
The case $p = 2$ gives a Hilbert space with inner product
$\langle f, g\rangle = \int_I f(x)\overline{g(x)}\,dx$.
\end{example}

\begin{example}
\label{ex:banach:sobolev}
For a bounded open set $\Omega \subset \RR^n$, $1 \leq p < \infty$, and
$k \in \NNz$, the Sobolev space $W^{k,p}(\Omega)$ consists of all functions
$f \in L^p(\Omega)$ whose distributional derivatives $D^\alpha f$ of order
$|\alpha| \leq k$ belong to $L^p(\Omega)$. Equipped with the norm
\begin{equation*}
    \|f\|_{W^{k,p}} = \Bigl(\sum_{|\alpha| \leq k}
        \|D^\alpha f\|_{L^p}^p\Bigr)^{1/p},
\end{equation*}
it is a Banach space. For $p=2$ one writes $H^k(\Omega) = W^{k,2}(\Omega)$,
which is a Hilbert space. These spaces appear as natural component spaces in
high-dimensional tensor problems; see Example~\ref{ex:min:sobolev} and
Section~\ref{sec:topTBT:bestapprox}.
\end{example}

\subsection*{Continuous linear maps}

Let $X$ and $Y$ be normed spaces. A linear map $A: X \to Y$ is
\emph{continuous} if and only if it is \emph{bounded}, meaning
\begin{equation*}
    \opnorm{A}{Y}{X} \coloneqq \sup_{x \neq 0} \frac{\|Ax\|_Y}{\|x\|_X}
    = \sup_{\|x\|_X \leq 1} \|Ax\|_Y < \infty.
\end{equation*}
The quantity $\opnorm{A}{Y}{X}$ is the \emph{operator norm} of $A$.
The space of all continuous linear maps from $X$ to $Y$, denoted
$\cL(X, Y)$, is itself a normed space under the operator norm; it is a
Banach space whenever $Y$ is complete. We write $\cL(X) \coloneqq \cL(X,X)$.
The set of bounded linear automorphisms (bounded invertible maps with
bounded inverse) is
\begin{equation*}
    \GL{X} \coloneqq \{A \in \cL(X) : A \text{ invertible}\}.
\end{equation*}
This is an open subset of $\cL(X)$, and hence a Banach manifold modelled
on $\cL(X)$; see Example~\ref{ex:manifold:GL}.

\subsection*{Dual spaces and reflexivity}

The \emph{topological dual} of a normed space $(X, \|\cdot\|)$ is
$\dual{X} \coloneqq \cL(X, \KK)$, equipped with the dual norm
\begin{equation}
\label{eq:dualnorm}
    \|\varphi\|_{\dual{X}} \coloneqq \sup_{\|x\| \leq 1} |\varphi(x)|.
\end{equation}
The dual $\dual{X}$ is always a Banach space, regardless of whether $X$
is complete. The \emph{bidual} $X^{**} \coloneqq (\dual{X})^*$ contains
$X$ as a canonically embedded subspace via the isometric embedding
$\iota: X \to X^{**}$, $\iota(x)(\varphi) = \varphi(x)$. When $\iota$ is
surjective, $X$ is called \emph{reflexive}. All $L^p$ spaces with
$1 < p < \infty$ are reflexive; $L^1$ and $L^\infty$ are not.

Reflexivity plays an important role in the existence of best approximations.
Recall that a closed set $C$ in a normed space $X$ is \emph{proximinal} if
every point of $X$ has at least one nearest point in $C$. In a reflexive
Banach space, every closed convex set is proximinal
(see, e.g., \cite[p.\,61]{Hackbusch2019}). This fact will be used
extensively in Chapter~\ref{chap:minsubspaces} to prove the existence
of best low-rank tensor approximations.

\subsection*{Convexity and smoothness}

Two geometric properties of the norm are of particular importance for
the Dirac--Frenkel variational principle in Chapter~\ref{chap:DF}.

\begin{definition}
\label{def:strictly_convex}
A normed space $(X, \|\cdot\|)$ is \emph{strictly convex} (or
\emph{rotund}) if $\|x + y\|/2 < 1$ for all $x, y \in X$ with
$\|x\| = \|y\| = 1$ and $x \neq y$.
\end{definition}

\begin{definition}
\label{def:uniformly_convex}
A normed space $(X, \|\cdot\|)$ is \emph{uniformly convex} if for every
$\varepsilon > 0$ there exists $\delta > 0$ such that
$\|x + y\|/2 \leq 1 - \delta$ whenever $\|x\|, \|y\| \leq 1$ and
$\|x - y\| \geq \varepsilon$.
\end{definition}

Every uniformly convex Banach space is reflexive and strictly convex
(Clarkson's theorem). In particular, $L^p$ spaces with $1 < p < \infty$
are uniformly convex. Strict convexity guarantees uniqueness of metric
projections onto closed convex sets, a property needed for the
well-posedness of the Dirac--Frenkel equation.

\begin{definition}
\label{def:smooth}
A Banach space $X$ is \emph{smooth} if the limit
$\lim_{t \to 0} \bigl(\|x + ty\| - \|x\|\bigr)/t$
exists for all $x, y \in X$ with $\|x\| = 1$.
\end{definition}

Smoothness is the dual property to strict convexity in the following
sense: $X$ is smooth if and only if $\dual{X}$ is strictly convex, and
$X$ is strictly convex if and only if $\dual{X}$ is smooth.

%------------------------------------------------------------
\section{Complemented subspaces and projections}
\label{sec:banach:complemented}
%------------------------------------------------------------

The central geometric notion underlying all manifold structures in this
monograph is that of a complemented (or split) subspace. Unlike the
Hilbert space setting, complementability is a non-trivial condition in
general Banach spaces.

\subsection*{Projections}

A bounded linear map $P \in \cL(X)$ is a \emph{projection} if $P^2 = P$.
In that case $\mathrm{Im}\,P$ is a closed subspace, $\mathrm{Ker}\,P$ is
closed, and $X = \mathrm{Im}\,P \oplus \mathrm{Ker}\,P$ (as a direct sum
of closed subspaces). The complementary projection is $\Id{X} - P$.

\begin{definition}
\label{def:split_subspace}
Let $X$ be a Banach space. A closed subspace $U \subset X$ is a
\emph{split subspace} (or \emph{complemented subspace}) of $X$ if
there exists a closed subspace $W \subset X$ such that
$X = U \oplus W$.
The pair $(U, W)$ is then called a \emph{pair of complementary
subspaces}, and $W$ is a \emph{(topological) complement} of $U$ in $X$.
We write $\Proj{U}{W}$ for the unique continuous projection from $X$
onto $U$ along $W$.
\end{definition}

The following proposition characterises split subspaces in terms of
projections and provides three equivalent conditions that will each be
used at different points of the text.

\begin{proposition}[{\cite[Proposition~4.2]{FHN2015}}]
\label{prop:split_equiv}
Let $X$ be a Banach space and $U \subset X$ a closed subspace. The
following conditions are equivalent:
\begin{enumerate}
\item[\textup{(a)}] $U$ is a split subspace of $X$.
\item[\textup{(b)}] There exists $P \in \cL(X)$ with $P^2 = P$ and
    $\mathrm{Im}\,P = U$.
\item[\textup{(c)}] There exists $Q \in \cL(X)$ with $Q^2 = Q$ and
    $\mathrm{Ker}\,Q = U$.
\end{enumerate}
\end{proposition}

In condition~(b), the projection $P$ is none other than $\Proj{U}{W}$
for any choice of closed complement $W$ of $U$; in~(c), one has $Q =
\Proj{W}{U}$.

\begin{proposition}[{\cite[Theorem~4.5]{FHN2015}}]
\label{prop:finitedim_split}
Let $X$ be a Banach space. Every finite-dimensional subspace of $X$ is
a split subspace.
\end{proposition}

In a Hilbert space, every closed subspace is complemented (the orthogonal
complement provides the split). This is no longer true in general Banach
spaces: by a classical result of Lindenstrauss and Tzafriri, a Banach
space in which every closed subspace is complemented is necessarily
isomorphic to a Hilbert space. The failure of complementability in
non-Hilbert spaces is precisely what makes the immersion theorems of
Chapter~\ref{chap:tucker_manifold} non-trivial.

\subsection*{The nilpotent algebra associated to a splitting}

The following result, central to the construction of local charts in
Chapters~\ref{chap:tucker_manifold} and~\ref{chap:TB_manifold},
identifies a particularly well-behaved sub-algebra of $\cL(X)$
associated to any splitting $X = U \oplus W$.

\begin{proposition}
\label{prop:nilpotent_algebra}
Let $X$ be a Banach space and $U, W \in \GrassB{X}$ with $X = U \oplus W$.
Define
\begin{equation*}
    \cL_{(U,W)}(X) \coloneqq
    \bigl\{
        \Proj{W}{U} \circ S \circ \Proj{U}{W} : S \in \cL(X)
    \bigr\}
    \cong \cL(U, W).
\end{equation*}
Then:
\begin{enumerate}
\item[\textup{(i)}] The natural map
    $\cL(U, W) \to \cL_{(U,W)}(X)$,
    $L \mapsto \Proj{W}{U} \circ L \circ \Proj{U}{W}$,
    is an isometric isomorphism.
\item[\textup{(ii)}] $\cL_{(U,W)}(X)$ is a nilpotent sub-algebra of
    $\cL(X)$: for every $L, L' \in \cL_{(U,W)}(X)$ one has $L \circ L' = 0$.
\item[\textup{(iii)}] As a consequence, for every
    $L \in \cL_{(U,W)}(X)$:
    \begin{equation}
    \label{eq:exp_nilpotent}
        \exp(L) \coloneqq \sum_{n=0}^{\infty} \frac{L^n}{n!}
        = \Id{X} + L,
        \qquad
        \exp(-L) = \Id{X} - L = (\Id{X} + L)^{-1}.
    \end{equation}
\end{enumerate}
\end{proposition}

\begin{proof}
For the isometry in~(i), note that
$\|\Proj{W}{U} \circ L \circ \Proj{U}{W}\|_{X \leftarrow X}
= \|L\|_{W \leftarrow U}$ by the definition of the operator norm.
For~(ii): if $L = \Proj{W}{U} \circ S \circ \Proj{U}{W}$ and
$L' = \Proj{W}{U} \circ S' \circ \Proj{U}{W}$, then
\begin{equation*}
    L \circ L' =
    \Proj{W}{U} \circ S \circ (\Proj{U}{W} \circ \Proj{W}{U})
    \circ S' \circ \Proj{U}{W} = 0,
\end{equation*}
since $\Proj{U}{W} \circ \Proj{W}{U} = 0$. Statement~(iii) follows
immediately from~(ii) and the series definition of the exponential.
\end{proof}

The identity $\exp(L) = \Id{X} + L$ means that the exponential map on
the nilpotent algebra $\cL_{(U,W)}(X)$ is a polynomial of degree one,
and in particular a global analytic diffeomorphism. This is the key
that makes the local charts of the Grassmann--Banach manifold and the
Tucker manifold globally well-defined; see
Section~\ref{sec:banach:grassmann}.

\subsection*{Common complements and chart compatibility}

The following lemma, which underlies the compatibility of overlapping
charts in the Grassmann--Banach manifold, characterises when two split
subspaces share a common complement.

\begin{lemma}[{\cite[Lemma~2.1]{DixmierDouady1963}}]
\label{lem:common_complement}
Let $X$ be a Banach space and let $U, U', W \in \GrassB{X}$ with
$X = U \oplus W = U' \oplus W$. The following conditions are equivalent:
\begin{enumerate}
\item[\textup{(a)}] $U$ and $U'$ have $W$ as a common complement, i.e.,
    $X = U \oplus W = U' \oplus W$.
\item[\textup{(b)}] The restriction $\Proj{U}{W}\big|_{U'} : U' \to U$
    is a linear isomorphism (with bounded inverse).
\end{enumerate}
Moreover, in this case
$\bigl(\Proj{U}{W}\big|_{U'}\bigr)^{-1} = \Proj{U'}{W}\big|_U$.
\end{lemma}

%------------------------------------------------------------
\section{The Grassmann--Banach manifold}
\label{sec:banach:grassmann}
%------------------------------------------------------------

The Grassmannian of a Banach space is the natural parameter space for
subspaces, and it serves as the base manifold for the Tucker and
tree-based tensor manifolds.

\begin{definition}
\label{def:grassmannian}
Let $X$ be a Banach space. The \emph{Grassmann manifold} (or
\emph{Grassmannian}) of $X$ is
\begin{equation*}
    \GrassB{X} \coloneqq \{U \subset X : U \text{ is a split subspace of } X\}.
\end{equation*}
For $r \in \NNz$, the \emph{finite Grassmannian} is
\begin{equation*}
    \Grass{r}{X} \coloneqq
    \{U \in \GrassB{X} : \dim U = r\}.
\end{equation*}
\end{definition}

Both $X$ and $\{0\}$ belong to $\GrassB{X}$.
By Proposition~\ref{prop:finitedim_split}, every $r$-dimensional subspace
of $X$ belongs to $\GrassB{X}$, so $\Grass{r}{X}$ is well defined.

\begin{theorem}[Grassmann--Banach manifold, {\cite{Douady1966}}]
\label{thm:grassmann_manifold}
Let $X$ be a Banach space. The Grassmannian $\GrassB{X}$ is an analytic
Banach manifold, with atlas
$\bigl\{\GrassB{W, X},\, \Psi_{U \oplus W}\bigr\}_{U \in \GrassB{X}}$,
where for each splitting $X = U \oplus W$:
\begin{enumerate}
\item[\textup{(i)}] $\GrassB{W, X} \coloneqq
    \{V \in \GrassB{X} : X = V \oplus W\}$
    is the chart domain at $U$.
\item[\textup{(ii)}] The chart map
    $\Psi_{U \oplus W}: \GrassB{W, X} \to \cL(U, W)$
    is defined by
    \begin{equation}
    \label{eq:grassmann_chart}
        \Psi_{U \oplus W}(V) \coloneqq
        \Proj{W}{U}\big|_V \circ \bigl(\Proj{U}{W}\big|_V\bigr)^{-1},
    \end{equation}
    with inverse
    $\Psi_{U \oplus W}^{-1}(L) = G(L) \coloneqq
    \{(\Id{X} + L)(u) : u \in U\}$.
\end{enumerate}
The transition maps are analytic. Each chart domain
$\GrassB{W, X}$ is a Banach manifold modelled on $\cL(U, W)$.
\end{theorem}

The chart map $\Psi_{U \oplus W}$ sends a subspace $V \in \GrassB{W,X}$
to the linear map $L: U \to W$ whose graph $G(L)$ equals $V$. The
identity $\Psi_{U\oplus W}^{-1}(L) = \exp(L)(U)$, combined with
Proposition~\ref{prop:nilpotent_algebra}, shows that the inverse chart is
a degree-one polynomial map (hence analytic). Analyticity of the
transition maps follows by a direct computation using
Lemma~\ref{lem:common_complement}.

\begin{remark}
\label{rem:grassmann_not_modelled}
The full Grassmannian $\GrassB{X}$ is a Banach manifold but is
\emph{not} modelled on a single Banach space: if $U$ and $U'$ are
finite-dimensional with $\dim U \neq \dim U'$, then $\cL(U, W)$ and
$\cL(U', W')$ are non-isomorphic. This is not a deficiency but a
structural fact: $\GrassB{X}$ decomposes into connected components
$\GrassB{X} = \bigsqcup_{r \in \NNz} \Grass{r}{X}$, each of which
\emph{is} modelled on a single Banach space.
\end{remark}

\begin{proposition}
\label{prop:grassr_connected}
For each $r \in \NNz$, the finite Grassmannian $\Grass{r}{X}$ is a
connected component of $\GrassB{X}$, and it is an analytic Banach manifold
modelled on $\cL(U, W)$ for any $U \in \Grass{r}{X}$ and complement $W$
\textup{(}the isomorphism class of $\cL(U, W)$ does not depend on the
choice of $U$ or $W$\textup{)}.
\end{proposition}

\begin{remark}[Grassmannian of a normed space]
\label{rem:grassmann_normed}
The construction extends to normed (non-complete) spaces. Let
$(X, \|\cdot\|)$ be a normed space with completion $\bar X$. Define
$\Grass{r}{X}$ as the set of $r$-dimensional subspaces of $X$ that are
also split in $\bar X$. By Lemma~\ref{lem:normed_grassmann}, every chart
domain $\GrassB{W, \bar X}$ is contained in $\Grass{r}{X}$, so the atlas
of $\Grass{r}{\bar X}$ restricts to an analytic atlas on $\Grass{r}{X}$.
This variant is needed in Section~\ref{sec:tucker_manifold:charts}, where
the component spaces $V_\alpha$ are normed but not necessarily complete.
\end{remark}

The tangent space to $\GrassB{X}$ at $U$ is identified as follows.

\begin{proposition}
\label{prop:grassmann_tangent}
For $U \in \GrassB{X}$ with complement $W$, the tangent space
$T_U \GrassB{X} \cong \cL(U, W)$.
\end{proposition}

This identification is consistent with the finite-dimensional case: in
$\RR^n$, the tangent space to the Grassmannian $Gr(r, n)$ at a subspace
$U$ is the space of linear maps from $U$ to its orthogonal complement,
recovering the classical description.

%------------------------------------------------------------
\section{Banach manifolds}
\label{sec:banach:manifolds}
%------------------------------------------------------------

We recall the definition of a Banach manifold in its general form, since
the full Grassmannian $\GrassB{X}$ and the full tensor space
$\VDnorm$ are manifolds not modelled on a fixed Banach space.

\begin{definition}
\label{def:banach_manifold}
Let $\mathbb{M}$ be a set. An \emph{atlas of class $\mathcal{C}^p$}
(for $p \geq 0$ or $p = \infty$ or $p = \omega$, the analytic case)
on $\mathbb{M}$ is a family of \emph{charts}
$\{(M_i, u_i)\}_{i \in A}$ satisfying:
\begin{enumerate}
\item[\textup{(AT1)}] $\{M_i\}_{i\in A}$ covers $\mathbb{M}$.
\item[\textup{(AT2)}] For each $i \in A$, the chart $u_i: M_i \to U_i$
    is a bijection onto an open set $U_i$ in some Banach space $X_i$,
    and for any $i, j \in A$, the set $u_i(M_i \cap M_j)$ is open in $X_i$.
\item[\textup{(AT3)}] Whenever $M_i \cap M_j \neq \emptyset$,
    the \emph{transition map}
    $u_j \circ u_i^{-1} : u_i(M_i \cap M_j) \to u_j(M_i \cap M_j)$
    is a diffeomorphism of class $\mathcal{C}^p$.
\end{enumerate}
Two atlases are \emph{compatible} if their union is again an atlas. A
\emph{$\mathcal{C}^p$-Banach manifold} is a set $\mathbb{M}$ together
with a maximal (complete) atlas of class $\mathcal{C}^p$. If all
$X_i$ are isomorphic to a single Banach space $X$, we say $\mathbb{M}$
is \emph{modelled on $X$}.
\end{definition}

\begin{example}
\label{ex:manifold:banach_space}
Every Banach space $X$ is a $\mathcal{C}^\infty$-Banach manifold
modelled on itself, with the single chart $(X, \Id{X})$.
\end{example}

\begin{example}
\label{ex:manifold:GL}
If $X$ is a Banach space, then $\GL{X}$ is an open subset of $\cL(X)$,
and hence an analytic Banach manifold modelled on $\cL(X)$. The maps
$A \mapsto A^{-1}$ and $(A,B) \mapsto AB$ are analytic
\textup{(}see~\cite[2.7, 2.8]{Upmeier1985}\textup{)}.
\end{example}

\begin{definition}
\label{def:tangent_space}
Let $\mathbb{M}$ be a $\mathcal{C}^p$-Banach manifold with $p \geq 1$
and $m \in \mathbb{M}$. A \emph{tangent vector} at $m$ is an equivalence
class of triples $(U_i, u_i, v)$ where $(U_i, u_i)$ is a chart at $m$
and $v \in X_i$, under the equivalence
$(U_i, u_i, v) \sim (U_j, u_j, w)$ if $D(u_j \circ u_i^{-1})(u_i(m)) v = w$.
The \emph{tangent space} $T_m \mathbb{M}$ is the set of all tangent vectors
at $m$, which inherits the structure of a Banach space from any chart.
\end{definition}

\begin{definition}
\label{def:immersion}
A $\mathcal{C}^1$-map $f: \mathbb{M} \to \mathbb{N}$ between Banach
manifolds is an \emph{immersion} at $m$ if $Df(m): T_m\mathbb{M} \to
T_{f(m)}\mathbb{N}$ is injective with closed split image. It is an
\emph{immersion} if it is an immersion at every point. An
\emph{immersed submanifold} of $\mathbb{N}$ is the image of an injective
immersion.
\end{definition}

The closedness and split condition on the image of $Df(m)$ is automatic
in finite dimensions and in Hilbert spaces (where every closed subspace is
split), but is a genuine constraint in Banach spaces. This is the key
difficulty resolved in Section~\ref{sec:tucker_manifold:immersion}.

%------------------------------------------------------------
\section{Multilinear maps and Fréchet differentiability}
\label{sec:banach:multilinear}
%------------------------------------------------------------

Since the tensor product map is multilinear, the following general results
will be used repeatedly.

\begin{definition}
A map $M: X_1 \times \cdots \times X_d \to Y$ between normed spaces is
\emph{multilinear} if it is linear in each argument separately. It is
\emph{bounded} (equivalently, \emph{continuous}) if
\begin{equation*}
    \|M\| \coloneqq \sup_{\substack{x_j \in X_j \\ \|x_j\|_j \leq 1}}
    \|M(x_1, \ldots, x_d)\|_Y < \infty.
\end{equation*}
\end{definition}

The following proposition shows that continuity implies a much stronger
regularity.

\begin{proposition}[{\cite[Proposition~79]{HJ2012}}]
\label{prop:multilinear_smooth}
Let $X_1, \ldots, X_d$ and $Y$ be normed spaces and let
$M: X_1 \times \cdots \times X_d \to Y$ be a continuous multilinear map.
Then $M$ is $\mathcal{C}^\infty$-Fréchet differentiable, and
$D^k M = 0$ for all $k > d$.
\end{proposition}

The vanishing of derivatives beyond order $d$ reflects the polynomial
nature of multilinear maps. In the complex case, a stronger conclusion
holds.

\begin{proposition}[{\cite[Theorem~160]{HJ2012}}]
\label{prop:analytic_c1}
Let $X$, $Y$ be complex Banach spaces and $U \subset X$ open. A map
$f: U \to Y$ is analytic if and only if it is $\mathcal{C}^1$-Fréchet
differentiable.
\end{proposition}

\begin{corollary}
\label{cor:multilinear_analytic}
Every continuous multilinear map between complex Banach spaces is analytic.
\end{corollary}

In real Banach spaces, $\mathcal{C}^\infty$ and analytic are distinct
notions; the tensor product map is $\mathcal{C}^\infty$ but need not be
analytic over $\RR$. For the real case, $\mathcal{C}^\infty$-regularity
is sufficient for the manifold constructions in Chapters~\ref{chap:tucker_manifold}
and~\ref{chap:TB_manifold}.

\begin{remark}
\label{rem:regularity_manifold}
Proposition~\ref{prop:multilinear_smooth} is the reason why the Tucker
and tree-based tensor manifolds are $\mathcal{C}^\infty$-Banach manifolds
(and analytic over $\CC$): their chart transition maps are compositions of
continuous multilinear maps with bounded linear projections, all of which
are $\mathcal{C}^\infty$-Fréchet differentiable.
\end{remark}

\subsection*{The Grassmannian revisited: Lemma on normed spaces}

For later use in Section~\ref{sec:tucker_manifold:charts}, we record
the following lemma about Grassmannians of normed (non-complete) spaces.

\begin{lemma}
\label{lem:normed_grassmann}
Let $(X, \|\cdot\|)$ be a normed space with completion $\bar X$. Let
$U \in \Grass{r}{\bar X}$ with $U \subset X$, and suppose
$\bar X = U \oplus W$ for some $W \in \GrassB{\bar X}$. Then every
$V \in \GrassB{W, \bar X}$ satisfies $V \subset X$.
\end{lemma}

\begin{proof}
Observe that $X = U \oplus (W \cap X)$, where $W \cap X$ is dense in $W$.
Suppose for contradiction that there exists $V \in \GrassB{W, \bar X}$
with $V \not\subset X$. Since $V$ is finite-dimensional and $\bar X =
V \oplus W$, the subspace $V \cap X$ is a proper subspace of $V$. On
the other hand, $X = (V \cap X) \oplus (W \cap X)$ implies $\bar X =
(V \cap X) \oplus W$, so $V \cap X \in \GrassB{W, \bar X}$. But
$\dim(V \cap X) < \dim V = r$, contradicting the fact that every element
of $\GrassB{W, \bar X}$ has dimension $r$.
\end{proof}


%------------------------------------------------------------
\section{Historical notes}
\label{sec:banach:history}
%------------------------------------------------------------

The theory of Banach spaces was developed primarily in the 1920s and
1930s, with Stefan Banach's monograph~\cite{Banach1932} providing the
first systematic treatment. The duality theory and reflexivity go back
to Banach and Alaoglu; the theorem that uniform convexity implies
reflexivity is due to Clarkson (1936) and Milman (1938).

\paragraph{Complemented subspaces.}
The study of complemented subspaces has its roots in the work of
Banach himself, who showed that $c_0$ has uncomplemented subspaces.
The decisive result --- that a Banach space in which every closed
subspace is complemented must be isomorphic to a Hilbert space ---
was proved by Lindenstrauss and Tzafriri in 1971. This explains why
complementability is a genuine and non-trivial constraint in the
Banach tensor setting of this monograph.

The characterisation of split subspaces via bounded projections
(Proposition~\ref{prop:split_equiv}) and the finite-dimensionality
criterion (Proposition~\ref{prop:finitedim_split}) are classical; the
formulation used here follows~\cite{FHN2015}. The nilpotent algebra
associated to a splitting (Proposition~\ref{prop:nilpotent_algebra})
and its use to identify the chart inverse as a degree-one polynomial
map is a key observation of~\cite{FHN2019}.

The lemma on common complements (Lemma~\ref{lem:common_complement})
is a version of results going back to Dixmier and Douady~\cite{DixmierDouady1963},
who studied continuous fields of Hilbert spaces. The formulation for
general Banach spaces and its application to the transition maps of the
Tucker manifold appears in~\cite{FHN2019}.

\paragraph{Grassmann--Banach manifold.}
The infinite-dimensional Grassmannian was constructed as an analytic
Banach manifold by Douady~\cite{Douady1966} in the context of moduli
spaces of compact analytic subspaces. The atlas described in
Theorem~\ref{thm:grassmann_manifold} goes back to this work; the
specific formulation in terms of the chart map~\eqref{eq:grassmann_chart}
and the identification of the tangent space with $\cL(U,W)$ follows
the presentation of~\cite{Upmeier1985}. The finite-dimensional
Grassmannian $Gr(r,n)$ is a smooth compact manifold of dimension
$r(n-r)$, a fact well known since Grassmann's original work in the
1840s; the infinite-dimensional version requires the split condition
as a substitute for compactness arguments.

\paragraph{Multilinear maps and Fréchet differentiability.}
The fact that continuous multilinear maps are $\mathcal{C}^\infty$
(Proposition~\ref{prop:multilinear_smooth}) is a standard result of
infinite-dimensional calculus; see~\cite{HJ2012} for a careful
treatment. The equivalence between $\mathcal{C}^1$ and analyticity
in complex Banach spaces (Proposition~\ref{prop:analytic_c1}) is the
Gâteaux--Cauchy theorem; see~\cite{HJ2012}. These two propositions are
the main tools that allow one to conclude that the Tucker and tree-based
tensor manifolds are $\mathcal{C}^\infty$ (real case) or analytic
(complex case) without further calculation: their atlas transition maps
are compositions of continuous multilinear maps and bounded linear
projections, all of which satisfy the hypotheses of
Propositions~\ref{prop:multilinear_smooth} and~\ref{prop:analytic_c1}.

%============================================================
%  ch03_algtensor.tex
%  Chapter 3: Algebraic Tensor Spaces
%
%  Tensor Formats in Banach Spaces
%  Antonio Falcó, CEU Universities
%
%  Estimated length: ~25 pp.
%============================================================

\chapter{Algebraic Tensor Spaces}
\label{chap:algtensor}

This chapter develops the purely algebraic foundations of tensor spaces.
We introduce the algebraic tensor product via its universal property,
establish the key notions of elementary tensor, matricisation, and multilinear
rank, and define the Tucker format and the set $\Tucker{\fr}{\VD}$ of tensors
of fixed Tucker rank. No topology is assumed in this chapter; that enters
in Chapter~\ref{chap:topTBT}.

Throughout, $D = \{1, \ldots, d\}$ with $d \geq 2$ is the \emph{dimension
index set}, and $V_\alpha$ ($\alpha \in D$) are vector spaces over $\KK$.
Concerning the foundational construction of tensor products we follow
Greub~\cite{Greub1981}.

%------------------------------------------------------------
\section{The algebraic tensor product}
\label{sec:algtensor:def}
%------------------------------------------------------------

\subsection*{Universal property and construction}

\begin{definition}
\label{def:algtensor}
The \emph{algebraic tensor product} of the spaces $V_\alpha$, $\alpha \in D$,
is a pair $(\VD, \bigotimes)$ consisting of a vector space $\VD$ and a
multilinear map
\begin{equation}
\label{eq:tensor_product_map}
    \bigotimes : \prod_{\alpha \in D} V_\alpha \longrightarrow \VD,
    \qquad
    (v_\alpha)_{\alpha \in D} \mapsto \bigotimes_{\alpha \in D} v_\alpha,
\end{equation}
satisfying the following \emph{universal property}: for every vector space
$Z$ and every multilinear map $M: \prod_{\alpha \in D} V_\alpha \to Z$,
there exists a unique linear map $\widehat{M}: \VD \to Z$ such that
$M = \widehat{M} \circ \bigotimes$.

We write $\VD = {}_a\bigotimes_{\alpha \in D} V_\alpha$ and refer to elements
of the form $\bigotimes_{\alpha \in D} v_\alpha$ (with $v_\alpha \in V_\alpha$)
as \emph{elementary tensors}. Every element of $\VD$ is a \emph{finite}
linear combination of elementary tensors.
\end{definition}

The algebraic tensor product exists and is unique up to isomorphism; see
\cite{Greub1981} for the standard construction. The prefix `$a$' in
${}_a\bigotimes$ emphasises the algebraic (non-topological) nature.
The underlying field is $\KK = \RR$ throughout, but all results hold
equally for $\KK = \CC$.

\begin{remark}
\label{rem:tensor_not_basis}
The set of elementary tensors is \emph{not} a basis of $\VD$; it
generates $\VD$ but with many linear dependencies. For example,
$\lambda v \otimes w = v \otimes \lambda w$ for all $\lambda \in \KK$.
If $\{e_i\}$ and $\{f_j\}$ are bases of $V_1$ and $V_2$ respectively,
then $\{e_i \otimes f_j\}$ is a basis of $V_1 \atens V_2$.
\end{remark}

\subsection*{Bilinear identification and the $\alpha$-unfolding}

For each $\alpha \in D$, the algebraic tensor space can be rewritten as
a two-factor product by grouping all indices except $\alpha$:
\begin{equation}
\label{eq:bilinear_split}
    \VD \cong V_\alpha \atens \Vbigalpha^{[\alpha]},
    \qquad
    \Vbigalpha^{[\alpha]} \coloneqq {}_a\bigotimes_{\beta \in D \setminus \{\alpha\}} V_\beta.
\end{equation}
Concretely, there is a canonical linear isomorphism
$\Phi_\alpha: \VD \xrightarrow{\;\sim\;} V_\alpha \atens \Vbigalpha^{[\alpha]}$
and throughout this monograph we identify $\bv \in \VD$ with
$\Phi_\alpha(\bv) \in V_\alpha \atens \Vbigalpha^{[\alpha]}$ without
further comment. An elementary tensor splits as
$\bigotimes_{\beta \in D} v_\beta \leftrightarrow v_\alpha \otimes
\bigotimes_{\beta \neq \alpha} v_\beta$.

\subsection*{Tensor products of linear maps}

Given linear maps $A_\alpha: V_\alpha \to W_\alpha$ for $\alpha \in D$,
their tensor product is the unique linear map
\begin{equation}
\label{eq:tensor_of_maps}
    \bigotimes_{\alpha \in D} A_\alpha :
    \VD \longrightarrow \mathbf{W}_D
    \coloneqq {}_a\bigotimes_{\alpha \in D} W_\alpha,
\end{equation}
determined on elementary tensors by
$\bigl(\bigotimes_{\alpha \in D} A_\alpha\bigr)
\bigl(\bigotimes_{\alpha \in D} v_\alpha\bigr)
= \bigotimes_{\alpha \in D} A_\alpha v_\alpha$.
This is the unique linear extension guaranteed by the universal property.
We denote the algebraic dual of $V$ by $\algdual{V} = L(V, \KK)$
and the space of all linear maps from $V$ to $W$ by $L(V, W)$.

\subsection*{Partitions and iterated tensor products}

\begin{definition}
\label{def:partition_tensor}
Let $\mathcal{P}$ be a partition of $D$ (i.e., a collection of
non-empty, pairwise disjoint subsets of $D$ whose union is $D$).
For each $\alpha \in \mathcal{P}$, set
$\Vbigalpha \coloneqq {}_a\bigotimes_{j \in \alpha} V_j$
(with the convention $\mathbf{V}_{\{j\}} = V_j$). Then the
\emph{iterated tensor product} associated to $\mathcal{P}$ is
\begin{equation}
\label{eq:iterated_tensor}
    {}_a\bigotimes_{\alpha \in \mathcal{P}} \Vbigalpha
    \cong \VD,
\end{equation}
where the isomorphism is canonical and consistent with the
identifications of elementary tensors.
\end{definition}

This associativity-of-tensor-products fact is fundamental: it allows us
to regroup the factors of $\VD$ according to any hierarchical structure,
which is precisely what the dimension partition trees of
Chapter~\ref{chap:trees} will exploit.

\begin{example}
\label{ex:iterated_tensor}
For $D = \{1,2,3,4\}$ and the partition
$\mathcal{P} = \{\{1,2\}, \{3,4\}\}$:
\begin{equation*}
    (V_1 \atens V_2) \atens (V_3 \atens V_4) \cong
    V_1 \atens V_2 \atens V_3 \atens V_4.
\end{equation*}
The partition $\mathcal{P}' = \{\{1\},\{2\},\{3\},\{4\}\}$ (trivial)
recovers the original four-fold product. The non-trivial partition
$\mathcal{P}$ is the structure underlying the Tucker format
(depth-2 tree), while more refined partitions correspond to hierarchical
formats; see Chapter~\ref{chap:trees}.
\end{example}

%------------------------------------------------------------
\section{Matricisation and multilinear rank}
\label{sec:algtensor:rank}
%------------------------------------------------------------

\subsection*{Matricisation in finite dimensions}

When the component spaces are finite-dimensional, the tensor $\bv \in \VD$
can be represented as a multi-dimensional array, and the identification
\eqref{eq:bilinear_split} induces a matrix.

\begin{definition}
\label{def:matricisation}
Let $D = \{1, \ldots, d\}$, $\dim V_\alpha = n_\alpha < \infty$, and
fix bases of each $V_\alpha$. Identify $\bv \in \VD$ with its
coordinate array $C \in \RR^{n_1 \times \cdots \times n_d}$.
For each $\alpha \in D$, the \emph{$\alpha$-matricisation} (or
\emph{$\alpha$-unfolding}) of $\bv$ is the matrix
\begin{equation*}
    \Mmat{\alpha}(\bv) \in \RR^{n_\alpha \times \prod_{\beta \neq \alpha} n_\beta}
\end{equation*}
obtained by reordering the coordinate array $C$ so that the $\alpha$-index
labels the rows and all other indices, in lexicographic order, label the
columns. The rank of this matrix,
$\rank \Mmat{\alpha}(\bv) \in \{0, 1, \ldots, n_\alpha\}$,
is the \emph{$\alpha$-rank} of $\bv$.
\end{definition}

The notion of $\alpha$-rank was introduced by Hitchcock~\cite{Hitchcock1927}
in 1927 under the name ``rank on the $j$th index''. In the
infinite-dimensional case, the correct replacement of
$\rank \Mmat{\alpha}(\bv)$ is the dimension of the minimal subspace
$\Umin{\alpha}{\bv}$, developed fully in Chapter~\ref{chap:minsubspaces}.

\subsection*{Explicit example: three-mode matricisation}

\begin{example}[Matricisation for $d=3$]
\label{ex:matricisation_d3}
Let $n_1=n_2=n_3=2$. Consider $\bv \in \RR^2 \otimes \RR^2 \otimes \RR^2$
with coordinate array
$C_{111}=C_{122}=C_{212}=C_{221}=1$, all other entries zero;
equivalently,
\begin{equation*}
    \bv = e_1{\otimes}e_1{\otimes}e_1
        + e_1{\otimes}e_2{\otimes}e_2
        + e_2{\otimes}e_1{\otimes}e_2
        + e_2{\otimes}e_2{\otimes}e_1.
\end{equation*}
The three matricisations, with columns ordered lexicographically by the
remaining indices, are:
\begin{align*}
    \Mmat{1}(\bv) &=
    \begin{pmatrix} 1 & 0 & 0 & 1 \\ 0 & 1 & 1 & 0 \end{pmatrix}
    \in \RR^{2 \times 4}, \\
    \Mmat{2}(\bv) &=
    \begin{pmatrix} 1 & 0 & 0 & 1 \\ 0 & 1 & 1 & 0 \end{pmatrix}
    \in \RR^{2 \times 4}, \\
    \Mmat{3}(\bv) &=
    \begin{pmatrix} 1 & 0 & 0 & 1 \\ 0 & 1 & 1 & 0 \end{pmatrix}
    \in \RR^{2 \times 4}.
\end{align*}
All three have rank~2, so the Tucker rank of $\bv$ is $(2,2,2)$ and
$\Umin{\alpha}{\bv} = \RR^2 = V_\alpha$ for each $\alpha$.
\end{example}

\subsection*{CP rank vs.\ Tucker rank}

\begin{example}[CP rank strictly exceeds Tucker rank]
\label{ex:cp_vs_tucker}
Consider the tensor $\bv$ of Example~\ref{ex:matricisation_d3}. Its
Tucker rank is $(2,2,2)$, yet its CP rank (the minimum number of
elementary summands) equals~4: the representation above has 4 terms,
and a 3-term representation can be ruled out by direct computation.

By contrast, the tensor
$\bw = e_1{\otimes}e_1{\otimes}e_1 + e_2{\otimes}e_2{\otimes}e_2
\in \RR^2 \otimes \RR^2 \otimes \RR^2$
has Tucker rank $(2,2,2)$ (each matricisation has rank~2) but CP rank~2.

This illustrates a fundamental distinction: the Tucker rank
$(r_1,\ldots,r_d)$ describes the subspace structure
($\Umin{\alpha}{\bv}$ has dimension $r_\alpha$), while the CP rank
counts minimal summands. One always has $r_\alpha \leq r_{\mathrm{CP}}(\bv)$,
but the inequality is strict in general for $d \geq 3$. Computing
Tucker rank is tractable (SVD of the matricisations); computing CP rank
is NP-hard~(H\aa stad, 1990). This computational advantage is a key
motivation for the Tucker-format framework of this monograph.
\end{example}

\begin{definition}
\label{def:fullrank_core}
For a rank tuple $\fr = (r_\alpha)_{\alpha \in D} \in \NN^d$, define
\begin{equation}
\label{eq:Rstar}
    \RR^{\times_{\alpha \in D} r_\alpha}_*
    \coloneqq
    \Bigl\{
        C \in \RR^{\times_{\alpha \in D} r_\alpha} :
        \rank \Mmat{\alpha}(C) = r_\alpha \;\text{ for all } \alpha \in D
    \Bigr\}.
\end{equation}
Elements of $\RR^{\times_{\alpha \in D} r_\alpha}_*$ are called
\emph{full-rank core tensors} of shape $\fr$.
\end{definition}

\begin{remark}
\label{rem:Rstar_open}
Since $\rank \Mmat{\alpha}(C) = r_\alpha$ if and only if
$\Mmat{\alpha}(C)\Mmat{\alpha}(C)^T \in \GL{\RR^{r_\alpha}}$,
and the determinant is a continuous function,
$\RR^{\times_{\alpha \in D} r_\alpha}_*$ is an open subset of
$\RR^{\times_{\alpha \in D} r_\alpha}$, hence a smooth finite-dimensional
manifold in its own right. In the scalar case $r_\alpha = 1$ for all
$\alpha$, we recover $\RR_* = \RR \setminus \{0\} = \GL{\RR}$.
\end{remark}

%------------------------------------------------------------
\section{The Tucker format}
\label{sec:algtensor:tucker}
%------------------------------------------------------------

\subsection*{Minimal subspaces: first appearance}

We introduce minimal subspaces here in their algebraic form; the full
theory (existence, uniqueness, lattice properties, best approximation)
is developed in Chapter~\ref{chap:minsubspaces}.

\begin{definition}
\label{def:minimal_subspace_algebraic}
Let $\bv \in \VD = {}_a\bigotimes_{\alpha \in D} V_\alpha$ and
$\alpha \in D$. Under the identification $\VD \cong V_\alpha \atens
\Vbigalpha^{[\alpha]}$, the \emph{minimal subspace} of $\bv$ in
direction $\alpha$ is the unique subspace $\Umin{\alpha}{\bv} \subset V_\alpha$
of smallest dimension such that
$\bv \in \Umin{\alpha}{\bv} \atens \Vbigalpha^{[\alpha]}$.
The \emph{$\alpha$-rank} of $\bv$ is $r_\alpha(\bv) \coloneqq \dim \Umin{\alpha}{\bv}$.
\end{definition}

The key properties — existence, uniqueness, and the symmetry
$\dim \Umin{\alpha}{\bv} = \dim \Umin{D \setminus \{\alpha\}}{\bv}$ —
are established in Chapter~\ref{chap:minsubspaces}. Here we proceed
taking these properties for granted.

\subsection*{The Tucker format and the set $\Tucker{\fr}{\VD}$}

\begin{definition}
\label{def:tucker_format}
Let $\fr = (r_\alpha)_{\alpha \in D} \in \NNz^d$ and let $U_\alpha \in
\Grass{r_\alpha}{V_\alpha}$ for each $\alpha \in D$. A tensor
$\bv \in \VD$ is said to be in \emph{Tucker format} with subspace
tuple $(U_\alpha)_{\alpha \in D}$ if
\begin{equation*}
    \bv \in {}_a\bigotimes_{\alpha \in D} U_\alpha.
\end{equation*}
The \emph{Tucker rank} of $\bv$ is the tuple
$\fr(\bv) = (r_\alpha(\bv))_{\alpha \in D}$,
and the set of tensors with fixed Tucker rank $\fr$ is
\begin{equation}
\label{eq:Tucker_set}
    \Tucker{\fr}{\VD} \coloneqq
    \bigl\{
        \bv \in \VD : r_\alpha(\bv) = r_\alpha \text{ for all } \alpha \in D
    \bigr\}.
\end{equation}
\end{definition}

The Tucker format thus represents $\bv$ as an element of a subspace
$U_1 \otimes \cdots \otimes U_d$ spanned by $r_\alpha$-dimensional
subspaces $U_\alpha \subset V_\alpha$. Unlike the subspaces, the
representation is not unique; the set $\Tucker{\fr}{\VD}$ collects all
tensors for which the \emph{minimal} such subspaces have dimension
exactly $r_\alpha$.

\subsection*{Admissible Tucker ranks}

\begin{definition}
\label{def:admissible_tucker}
A tuple $\fr = (r_\alpha)_{\alpha \in D} \in \NNz^d$ is an
\emph{admissible Tucker rank} for $\VD$ if there exists $\bv \in \VD$
with $r_\alpha(\bv) = r_\alpha$ for all $\alpha \in D$. The set of all
admissible Tucker ranks is
\begin{equation}
\label{eq:AD_tucker}
    \AD{\VD} \coloneqq
    \bigl\{
        (r_\alpha(\bv))_{\alpha \in D} \in \NNz^d : \bv \in \VD
    \bigr\}.
\end{equation}
\end{definition}

The zero tuple $\mathbf{0} = (0,\ldots,0)$ is always admissible
(corresponding to $\bv = 0$). The tuple $\mathbf{1} = (1,\ldots,1)$
is admissible if and only if $\dim V_\alpha \geq 1$ for all $\alpha$.
Since the $r_\alpha$ are bounded above by $\dim V_\alpha$ (in finite
dimensions), $\AD{\VD}$ is a finite set when all $V_\alpha$ are
finite-dimensional.

The following decomposition is one of the foundational structural facts
of the whole theory.

\begin{proposition}
\label{prop:tucker_decomposition}
The algebraic tensor space decomposes as a disjoint union:
\begin{equation}
\label{eq:tucker_disjoint}
    \VD = \bigsqcup_{\fr \in \AD{\VD}} \Tucker{\fr}{\VD}.
\end{equation}
\end{proposition}

\begin{proof}
Since every $\bv \in \VD$ has a well-defined Tucker rank tuple
(Chapter~\ref{chap:minsubspaces}), the sets $\Tucker{\fr}{\VD}$ cover
$\VD$; they are disjoint by definition of the rank.
\end{proof}

The decomposition \eqref{eq:tucker_disjoint} is the algebraic precursor
of the manifold decomposition established in
Chapter~\ref{chap:tucker_manifold}: each component $\Tucker{\fr}{\VD}$
will be shown to carry the structure of an infinite-dimensional
$\mathcal{C}^\infty$-Banach manifold.

\subsection*{Explicit representation of Tucker tensors}

The following lemma, which combines the definition of $\Tucker{\fr}{\VD}$
with the properties of minimal subspaces, gives three equivalent
characterisations of Tucker tensors. It is the algebraic heart of the
Tucker manifold construction.

\begin{lemma}
\label{lem:tucker_characterisation}
Let $\VD = {}_a\bigotimes_{\alpha \in D} V_\alpha$ and $\fr \in \AD{\VD}$
with $\fr \neq \mathbf{0}$. For $\bv \in \VD$, the following conditions
are equivalent.
\begin{enumerate}
\item[\textup{(a)}] $\bv \in \Tucker{\fr}{\VD}$.
\item[\textup{(b)}] \textup{(Tucker representation)} For each $\alpha \in D$ there exists a
    basis $\{u^{(\alpha)}_{i_\alpha} : 1 \leq i_\alpha \leq r_\alpha\}$
    of $\Umin{\alpha}{\bv}$ and a unique full-rank core tensor
    $\CD \in \RR^{\times_{\alpha \in D} r_\alpha}_*$,
    once the bases are fixed, such that
    \begin{equation}
    \label{eq:tucker_core}
        \bv = \sum_{\substack{1 \leq i_\alpha \leq r_\alpha \\ \alpha \in D}}
        \CD_{(i_\alpha)_{\alpha \in D}}
        \bigotimes_{\alpha \in D} u^{(\alpha)}_{i_\alpha}.
    \end{equation}
\item[\textup{(c)}] \textup{(Bilinear form)} For each $\alpha \in D$ there exist
    linearly independent vectors
    $\{u^{(\alpha)}_{i_\alpha}\}_{i_\alpha=1}^{r_\alpha}$ in $V_\alpha$
    and linearly independent vectors
    $\{\bU^{(\alpha)}_{i_\alpha}\}_{i_\alpha=1}^{r_\alpha}$
    in $\Vbigalpha^{[\alpha]}$ such that
    \begin{equation}
    \label{eq:tucker_bilinear}
        \bv = \sum_{i_\alpha=1}^{r_\alpha}
        u^{(\alpha)}_{i_\alpha} \otimes \bU^{(\alpha)}_{i_\alpha}.
    \end{equation}
    Moreover, when \eqref{eq:tucker_core} holds, the vectors
    $\bU^{(\alpha)}_{i_\alpha}$ are given by
    \begin{equation}
    \label{eq:tucker_Ualpha}
        \bU^{(\alpha)}_{i_\alpha} =
        \sum_{\substack{1 \leq i_\beta \leq r_\beta \\ \beta \in D \setminus \{\alpha\}}}
        \CD_{i_\alpha, (i_\beta)_{\beta \neq \alpha}}
        \bigotimes_{\beta \neq \alpha} u^{(\beta)}_{i_\beta}.
    \end{equation}
\end{enumerate}
\end{lemma}

\begin{proof}
\textbf{(a) $\Leftrightarrow$ (c).} If (a) holds, then by
Proposition~\ref{prop:min_d2} (Chapter~\ref{chap:minsubspaces}),
$\bv \in \Umin{\alpha}{\bv} \atens \Umin{D\setminus\{\alpha\}}{\bv}$
with $\dim \Umin{\alpha}{\bv} = \dim \Umin{D\setminus\{\alpha\}}{\bv} = r_\alpha$,
which gives~(c). Conversely, (c) implies $\dim \Umin{\alpha}{\bv} = r_\alpha$
for all $\alpha$, hence~(a).

\textbf{(b) $\Rightarrow$ (c).} Immediate.

\textbf{(c) $\Rightarrow$ (b).} From (c), one has
$\Umin{\alpha}{\bv} = \spann\{u^{(\alpha)}_{i_\alpha}\}$ for each $\alpha$,
and since $\bv \in {}_a\bigotimes_{\alpha \in D} \Umin{\alpha}{\bv}$,
there exists $\CD \in \RR^{\times_{\alpha \in D} r_\alpha}$ satisfying
\eqref{eq:tucker_core}. From \eqref{eq:tucker_Ualpha} and the linear
independence of the $\bU^{(\alpha)}_{i_\alpha}$, one deduces
$\rank \Mmat{\alpha}(\CD) = r_\alpha$ for all $\alpha$, so
$\CD \in \RR^{\times_{\alpha \in D} r_\alpha}_*$.
\end{proof}

\begin{remark}
\label{rem:tucker_identification}
From Lemma~\ref{lem:tucker_characterisation} and
Remark~\ref{rem:Rstar_open}, once bases of the minimal subspaces
$\Umin{\alpha}{\bv}$ are fixed ($\alpha \in D$), there is a natural
identification
\begin{equation}
\label{eq:tucker_iso}
    \Tucker{\fr}\Bigl({}_a\bigotimes_{\alpha \in D} \Umin{\alpha}{\bv}\Bigr)
    \xrightarrow{\;\sim\;}
    \RR^{\times_{\alpha \in D} r_\alpha}_*,
    \qquad
    \bu \mapsto E^{(D)},
\end{equation}
where $E^{(D)}$ is the coordinate array of $\bu$ in the chosen bases.
This identification is a diffeomorphism (both directions are polynomial),
and it is the local model for the Tucker manifold chart; see
Section~\ref{sec:tucker_manifold:charts}.
\end{remark}

\subsection*{The Tucker format as a flat case of tree-based formats}

The Tucker format corresponds to the \emph{shallowest} possible
hierarchical structure: the dimension partition tree $\TDtuck$ of
depth~1, whose root is $D$ and whose children are the singletons
$\{1\}, \ldots, \{d\}$:
\begin{equation*}
    \TDtuck = \bigl\{D, \{1\}, \{2\}, \ldots, \{d\}\bigr\},
    \qquad
    \Sons{D} = \bigl\{\{1\},\ldots,\{d\}\bigr\} = \Leaves[\TDtuck].
\end{equation*}
In Chapter~\ref{chap:algTBT} we will define tree-based formats for
arbitrary dimension partition trees $\TD$ and show that
$\Tucker{\fr}{\VD} = \FT{\fr}{\VD}{\TDtuck}$.

%------------------------------------------------------------
\section{The Hilbert--Schmidt norm as a tensor product norm}
\label{sec:algtensor:hs}
%------------------------------------------------------------

The Hilbert--Schmidt operators provide the most natural bridge between
the abstract tensor framework of this chapter and the concrete
infinite-dimensional setting of Part~\ref{part:quantum}. This section
shows that the Hilbert--Schmidt norm is a crossnorm on $H \otimes H^*$,
computes it explicitly for functions in $L^2$, and connects the
Hilbert--Schmidt rank with the Tucker rank.

\subsection*{The Hilbert tensor product}

Let $H$ be a separable complex Hilbert space with inner product
$\langle\cdot,\cdot\rangle$ and dual $H^*$ (Riesz isomorphism
$J: H \to H^*, J(v) = \langle v, \cdot\rangle$). The \emph{Hilbert
tensor norm} on $H \otimes_a H^*$ is defined for elementary tensors by
\begin{equation}
\label{eq:hilbert_tensor_norm}
    \|u \otimes \varphi\|_H \coloneqq \|u\|_H \|\varphi\|_{H^*},
\end{equation}
and extended to all finite sums by
\begin{equation}
\label{eq:hilbert_tensor_norm_sum}
    \Bigl\|\sum_k u_k \otimes \varphi_k\Bigr\|_H^2
    \coloneqq \sum_{j,k} \varphi_j(u_k)\,\overline{\varphi_k(u_j)}
    = \Tr\bigl(([\varphi_j(u_k)]_{jk})^* [\varphi_j(u_k)]_{jk}\bigr).
\end{equation}
The completion of $(H \otimes_a H^*, \|\cdot\|_H)$ is the space
$\cT{2}(H)$ of Hilbert--Schmidt operators, with
$\cT{2}(H) \cong H \otimes_H H^*$ isometrically
(Theorem~\ref{thm:app:hs_tensor} in Appendix~\ref{app:quantum}).

\begin{proposition}[Hilbert tensor norm is a crossnorm satisfying \condINJ]
\label{prop:hilbert_tensor_crossnorm}
\label{rem:grothendieck_hilbert}
The Hilbert tensor norm $\|\cdot\|_H$ on $H \otimes_a H^*$ satisfies:
\begin{enumerate}
\item[\textup{(i)}] \condCROSS: $\|u \otimes J(v)\|_H = \|u\|_H\|v\|_H$
    for all $u, v \in H$.
\item[\textup{(ii)}] \condINJ: $\|\cdot\|_\varepsilon \leq \|\cdot\|_H$.
\item[\textup{(iii)}] $\|\cdot\|_H \leq \|\cdot\|_\pi$ (dominated by
    the projective norm).
\item[\textup{(iv)}] The Grothendieck inequality gives
    $\|\cdot\|_\pi \leq K_G \|\cdot\|_\varepsilon$ on $H \otimes_a H^*$,
    so all three norms are equivalent with constant $K_G \approx 1.78$.
\end{enumerate}
\end{proposition}

\begin{proof}
Part~(i) is immediate from~\eqref{eq:hilbert_tensor_norm}.
Part~(ii): for any $\phi \in H^*$ with $\|\phi\|_{H^*} \leq 1$ and
$\psi \in (H^*)^*$ with $\|\psi\|_{(H^*)^*} \leq 1$, identifying
$\psi = J^{-1}(\psi) \in H$ (using $J$), one has
$|(\phi \otimes \psi)(\bv)| = |\phi((\text{Id} \otimes \psi)(\bv))|
\leq \|\bv\|_H$ by the Cauchy--Schwarz inequality in $\cT{2}(H)$,
giving $\|\bv\|_\varepsilon \leq \|\bv\|_H$.
Part~(iii) is standard: $\|\cdot\|_H \leq \|\cdot\|_\pi$ for any
reasonable crossnorm. Part~(iv) is Grothendieck's theorem
(Remark~\ref{rem:grothendieck_hilbert}).
\end{proof}

\subsection*{Hilbert--Schmidt rank equals Tucker rank}

\begin{proposition}[Hilbert--Schmidt rank and Tucker rank]
\label{prop:app:hs_rank}
Let $A \in \cT{2}(H)$ and identify $A$ with a tensor $\bA \in H \otimes_H H^*$.
Then $r_1(\bA) = r_2(\bA) = \mathrm{rank}(A)$, and the SVD
$A = \sum_{k=1}^r s_k v_k \otimes J(u_k)$ is the Tucker decomposition
with minimal subspaces $\Umin{1}{\bA} = \overline{\mathrm{Im}(A^*)}
\subset H$ and $\Umin{2}{\bA} = J(\overline{\mathrm{Im}(A)}) \subset H^*$.
\end{proposition}

\begin{proof}
Under the identification $\cT{2}(H) \cong H \otimes_H H^*$, the
contraction map $T_\varphi: H \otimes_H H^* \to H^*$,
$T_\varphi(\bA) = (\varphi \otimes \Id{H^*})(\bA) = \varphi(A(\cdot))$
has range $J(\mathrm{Im}(A))$ when $\varphi$ ranges over $H^*$. Hence
$\Umin{2}{\bA} = J(\overline{\mathrm{Im}(A)})$, and similarly
$\Umin{1}{\bA} = \overline{\mathrm{Im}(A^*)}$.
\end{proof}

\subsection*{Explicit computation in $L^2$}

\begin{example}[Hilbert--Schmidt norm of a separable kernel]
\label{ex:hs_l2}
Let $H = L^2([0,1])$ and $H^* \cong H$ via $J$. Consider the rank-$r$
kernel $K(x,y) = \sum_{k=1}^r \lambda_k \phi_k(x)\psi_k(y)$, viewed as
an element of $L^2([0,1]^2) \cong L^2([0,1]) \otimes_H L^2([0,1])$.
The associated Hilbert--Schmidt operator $T_K: L^2 \to L^2$,
$T_K f = \int_0^1 K(\cdot, y) f(y)\,dy$, has:
\begin{itemize}
\item \emph{Hilbert--Schmidt norm:}
    $\|T_K\|_{\cT{2}}^2 = \|K\|_{L^2([0,1]^2)}^2
    = \int_0^1 \int_0^1 |K(x,y)|^2\,dx\,dy$.
\item \emph{Tucker rank:} $r_1(K) = r_2(K) = r$
    (if $\{\phi_k\}$, $\{\psi_k\}$ are linearly independent).
\item \emph{Minimal subspaces:} $\Umin{1}{K} = \spann\{\phi_1,\ldots,\phi_r\}$
    and $\Umin{2}{K} = \spann\{\psi_1,\ldots,\psi_r\}$ in $L^2([0,1])$.
\end{itemize}
For the specific case $K(x,y) = xy$ (rank~1 kernel):
$\|K\|_{\cT{2}}^2 = \int_0^1 x^2\,dx \int_0^1 y^2\,dy = \tfrac{1}{9}$,
so $\|K\|_{\cT{2}} = \tfrac{1}{3}$. The Tucker rank is $(1,1)$ and
the Tucker decomposition is $K = \phi \otimes J(\psi)$ with
$\phi(x) = x$, $\psi(y) = y$.
\end{example}

\begin{remark}[The Hilbert tensor norm in the real case]
\label{rem:hs_real}
When $\KK = \RR$ and $H = \RR^n$, the Hilbert--Schmidt operators
on $\RR^n$ are exactly the $n \times n$ matrices with the Frobenius norm
$\|A\|_F = (\sum_{ij} A_{ij}^2)^{1/2}$. The identification
$\cT{2}(\RR^n) \cong \RR^n \otimes_H \RR^n$ (using $H^* \cong H$
since $J$ is the identity in the real self-dual case) gives
$\|A\|_F = \|A\|_H$, and the Tucker rank of $A$ viewed as a tensor in
$\RR^n \otimes \RR^n$ equals $\mathrm{rank}(A)$ as a matrix. This is
the finite-dimensional instance of Proposition~\ref{prop:app:hs_rank}.
\end{remark}

%------------------------------------------------------------
\section{The tensor product map and its differential}
\label{sec:algtensor:tensmap}
%------------------------------------------------------------

When the component spaces are normed, the tensor product map inherits
both topological and differential properties that underpin the manifold
constructions in Part~\ref{part:geometry}.

\begin{definition}
\label{def:tensor_product_continuous}
Let $(V_\alpha, \|\cdot\|_\alpha)$ be normed spaces and let $\|\cdot\|_D$
be a norm on $\VD$. The \emph{tensor product map}
\begin{equation}
\label{eq:tensor_product_normed}
    \bigotimes :
    \prod_{\alpha \in D}(V_\alpha, \|\cdot\|_\alpha) \longrightarrow
    (\VD, \|\cdot\|_D),
    \qquad
    (v_\alpha)_{\alpha \in D} \mapsto \bigotimes_{\alpha \in D} v_\alpha,
\end{equation}
(where the product is equipped with the $\max$ norm
$\|(v_\alpha)\|_\times = \max_\alpha \|v_\alpha\|_\alpha$) is said to be
\emph{continuous} if $\|\bigotimes_\alpha v_\alpha\|_D \leq C \prod_\alpha
\|v_\alpha\|_\alpha$ for some constant $C > 0$.
\end{definition}

Continuity of the tensor product map is equivalent to
$\|\bigotimes\| < \infty$; it holds whenever $\|\cdot\|_D$ is a
crossnorm (Definition~\ref{def:crossnorm}) on $\VD$.

\begin{proposition}
\label{prop:tensor_smooth}
Let $(V_\alpha, \|\cdot\|_\alpha)$ be normed spaces and let $\|\cdot\|_D$
be a norm on $\VD$ such that the tensor product map
\eqref{eq:tensor_product_normed} is continuous. Then:
\begin{enumerate}
\item[\textup{(i)}] The map \eqref{eq:tensor_product_normed} is
    $\mathcal{C}^\infty$-Fréchet differentiable, with differential
    \begin{equation}
    \label{eq:diff_tensor_product}
        D\bigotimes(v_1,\ldots,v_d)(w_1,\ldots,w_d)
        = \sum_{\alpha \in D}
          v_1 \otimes \cdots \otimes v_{\alpha-1} \otimes
          w_\alpha \otimes v_{\alpha+1} \otimes \cdots \otimes v_d.
    \end{equation}
\item[\textup{(ii)}] The tensor product of bounded linear maps,
    \begin{equation*}
        \bigotimes :
        \prod_{\alpha \in D} \cL(U_\alpha, V_\alpha)
        \longrightarrow
        \cL\Bigl({}_a\bigotimes_{\alpha \in D} U_\alpha,\, \VD\Bigr),
        \quad
        (A_\alpha) \mapsto \bigotimes_{\alpha \in D} A_\alpha,
    \end{equation*}
    is also continuous and $\mathcal{C}^\infty$-Fréchet differentiable
    for any finite-dimensional subspaces $U_\alpha \subset V_\alpha$.
\end{enumerate}
\end{proposition}

\begin{proof}
Part~(i) follows directly from Proposition~\ref{prop:multilinear_smooth},
since $\bigotimes$ is a continuous multilinear map. The formula for the
differential follows by the Leibniz rule for multilinear maps. For~(ii),
one checks boundedness of the map using the continuity of
\eqref{eq:tensor_product_normed}:
\begin{equation*}
    \|\bigotimes_\alpha A_\alpha\|_{\VD \leftarrow \bigotimes U_\alpha}
    \leq C \prod_\alpha \|A_\alpha\|_{V_\alpha \leftarrow U_\alpha},
\end{equation*}
where $C$ is the constant in Definition~\ref{def:tensor_product_continuous},
and the finite-dimensionality of $U_\alpha$ ensures
$L(U_\alpha, V_\alpha) = \cL(U_\alpha, V_\alpha)$.
Fréchet differentiability then follows again from
Proposition~\ref{prop:multilinear_smooth}.
\end{proof}

\begin{remark}
\label{rem:regularity_Tucker}
Proposition~\ref{prop:tensor_smooth} is the reason why the Tucker manifold
$\Tucker{\fr}{\VD}$ is $\mathcal{C}^\infty$ (not merely continuous): all
coordinate change maps between charts are built from compositions of
projections and the tensor product map, all of which are
$\mathcal{C}^\infty$-Fréchet differentiable. See
Section~\ref{sec:tucker_manifold:cinfty} for the precise argument.
\end{remark}

\subsection*{Banach tensor spaces}

When a norm is fixed on $\VD$, its completion is the natural ambient space.

\begin{definition}
\label{def:banach_tensor_space}
Let $(V_\alpha, \|\cdot\|_\alpha)$ be normed spaces and $\|\cdot\|_D$ a
norm on $\VD = {}_a\bigotimes_{\alpha \in D} V_\alpha$. The
\emph{Banach tensor space} associated to this data is the completion
\begin{equation*}
    \VDnorm
    \coloneqq
    {}_{\|\cdot\|_D}\bigotimes_{\alpha \in D} V_\alpha
    \coloneqq
    \overline{\VD}^{\|\cdot\|_D}.
\end{equation*}
If $\VDnorm$ is a Hilbert space, it is called a \emph{Hilbert tensor space}.
\end{definition}

\begin{example}
\label{ex:Lp_tensor}
For $I_\alpha \subset \RR$ ($\alpha \in D$) and $1 \leq p < \infty$,
the Lebesgue spaces $L^p(I_\alpha)$ satisfy
\begin{equation*}
    L^p\Bigl(\prod_{\alpha \in D} I_\alpha\Bigr)
    = {}_{\|\cdot\|_{0,p}}\bigotimes_{\alpha \in D} L^p(I_\alpha),
\end{equation*}
where $\|\cdot\|_{0,p}$ denotes the standard $L^p$ norm on the product
domain. For $p=2$ this is a Hilbert tensor space.
\end{example}

\begin{example}
\label{ex:sobolev_tensor}
For Sobolev spaces $H^{N,p}(I_\alpha)$ with $1 < p < \infty$ and
$\mathbf{I} = \prod_\alpha I_\alpha$:
\begin{equation*}
    H^{N,p}(\mathbf{I})
    = {}_{\|\cdot\|_{N,p}}\bigotimes_{\alpha \in D} H^{N,p}(I_\alpha).
\end{equation*}
This is reflexive and separable. For $p=2$ it is a Hilbert tensor space.
These spaces arise naturally in the numerical treatment of high-dimensional
PDEs; see Section~\ref{sec:min:bestapprox} and
Section~\ref{sec:topTBT:bestapprox}.
\end{example}

\begin{remark}
\label{rem:algebraic_dense}
The algebraic tensor space $\VD$ is dense in $\VDnorm$ by construction.
The Tucker manifold $\Tucker{\fr}{\VD}$ lives inside the algebraic space,
but its closure (as a subset of $\VDnorm$) and its topology as a manifold
are determined by the norm $\|\cdot\|_D$ and, in particular, by whether
$\|\cdot\|_D$ is a crossnorm. This interplay between algebraic and
topological structure is a recurring theme throughout the monograph.
\end{remark}


%============================================================
%  ch04_minsubspaces.tex
%  Chapter 4: Minimal Subspaces in Tensor Representations
%
%  Tensor Formats in Banach Spaces
%  Antonio Falcó, CEU Universities
%
%  Based on: Falcó & Hackbusch, Found. Comput. Math. 12 (2012)
%  Estimated length: ~40 pp.
%============================================================

\chapter{Minimal Subspaces in Tensor Representations}
\label{chap:minsubspaces}

In Chapter~\ref{chap:algtensor} we introduced the Tucker format and
defined the minimal subspace $\Umin{\alpha}{\bv}$ informally, taking
its existence and properties for granted. This chapter provides the
complete theory: existence and uniqueness (Section~\ref{sec:min:existence}),
the lattice structure and inclusion under refinement of partitions
(Section~\ref{sec:min:lattice}), admissible ranks and the disjoint
decomposition of $\VD$ (Section~\ref{sec:min:admissible}), and the
existence of best approximations in Banach tensor spaces
(Section~\ref{sec:min:bestapprox}). The results in this chapter are
taken from~\cite{FH2012}, expanded with complete proofs and additional
examples.

%------------------------------------------------------------
\section{Existence and uniqueness of minimal subspaces}
\label{sec:min:existence}
%------------------------------------------------------------

\subsection*{The case $d = 2$: Schmidt decomposition}

The case of two factors $\VD = V_1 \atens V_2$ is classical and serves
as the inductive base. The key observation is a lattice property of
tensor subspaces.

\begin{lemma}[Intersection lemma]
\label{lem:intersection_d2}
Let $X_i, Y_i$ be subspaces of $V_i$ for $i = 1, 2$. Then
\begin{equation}
\label{eq:intersection_d2}
    (X_1 \atens X_2) \cap (Y_1 \atens Y_2)
    = (X_1 \cap Y_1) \atens (X_2 \cap Y_2).
\end{equation}
\end{lemma}

\begin{proof}
The inclusion $(X_1 \cap Y_1) \atens (X_2 \cap Y_2) \subseteq
(X_1 \atens X_2) \cap (Y_1 \atens Y_2)$ is immediate. For the
reverse inclusion, let $\bv \in (X_1 \atens X_2) \cap (Y_1 \atens Y_2)$
with representations
\begin{equation*}
    \bv = \sum_{\nu=1}^{n_x} x^{(1)}_\nu \otimes x^{(2)}_\nu
       = \sum_{\nu=1}^{n_y} y^{(1)}_\nu \otimes y^{(2)}_\nu,
    \quad
    x^{(i)}_\nu \in X_i,\; y^{(i)}_\nu \in Y_i.
\end{equation*}
By adding terms with zero coefficient, assume $n_x = n_y = n$ and
that $\{x^{(1)}_\nu\}$ and $\{y^{(1)}_\nu\}$ are each linearly
independent (passing to a minimal representation). The two
representations yield the same matrix $M_1(\bv)$ under any identification
with a bilinear form, so $\spann\{x^{(1)}_\nu\} = \spann\{y^{(1)}_\nu\}
\subset X_1 \cap Y_1$. A symmetric argument for the second factor gives
$\spann\{x^{(2)}_\nu\} \subset X_2 \cap Y_2$, so
$\bv \in (X_1 \cap Y_1) \atens (X_2 \cap Y_2)$.
\end{proof}

\begin{proposition}[Minimal subspace: case $d=2$]
\label{prop:min_d2}
Let $\bv \in V_1 \atens V_2$. There exist unique subspaces
$U_i \subset V_i$ of finite dimension such that:
\begin{enumerate}
\item[\textup{(i)}] $\bv \in U_1 \atens U_2$.
\item[\textup{(ii)}] If $\bv \in X_1 \atens X_2$ for any subspaces $X_i$,
    then $U_i \subset X_i$.
\end{enumerate}
Moreover, $\dim U_1 = \dim U_2$.
\end{proposition}

\begin{proof}
\textbf{Existence.} Fix any representation $\bv = \sum_\nu v^{(1)}_\nu
\otimes v^{(2)}_\nu$ and set $U_i = \spann\{v^{(i)}_\nu\}$. By passing
to a representation with linearly independent $v^{(1)}_\nu$, one can also
take $v^{(2)}_\nu$ linearly independent (otherwise row reduce).
Then $\dim U_1 = \dim U_2$ (both equal the rank of the bilinear form).

\textbf{Minimality.} If $\bv \in X_1 \atens X_2$, then by the
intersection lemma, $\bv \in (U_1 \cap X_1) \atens (U_2 \cap X_2)$,
which forces $U_i \subset X_i$ by minimality of $U_i$.

\textbf{Uniqueness.} If $U_i$ and $U_i'$ both satisfy (i)--(ii),
minimality of $U_i$ applied to $U_i'$ gives $U_i \subset U_i'$,
and vice versa.
\end{proof}

\subsection*{The general case: existence and uniqueness}

\begin{theorem}[Existence and uniqueness of minimal subspaces]
\label{thm:min_existence}
Let $D = \{1, \ldots, d\}$, $\VD = {}_a\bigotimes_{\alpha \in D}
V_\alpha$, and $\bv \in \VD$. For each $\alpha \in D$, under the
identification $\VD \cong V_\alpha \atens \Vbigalpha^{[\alpha]}$,
there exists a unique subspace $\Umin{\alpha}{\bv} \subset V_\alpha$
of finite dimension $r_\alpha < \infty$ such that:
\begin{enumerate}
\item[\textup{(a)}] $\bv \in \Umin{\alpha}{\bv} \atens
    \Vbigalpha^{[\alpha]}$.
\item[\textup{(b)}] If $\bv \in U_\alpha \atens \Vbigalpha^{[\alpha]}$
    for any subspace $U_\alpha \subset V_\alpha$, then
    $\Umin{\alpha}{\bv} \subset U_\alpha$.
\end{enumerate}
Furthermore, there exists a unique subspace
$\Umin{D\setminus\{\alpha\}}{\bv} \subset \Vbigalpha^{[\alpha]}$
of dimension $r_\alpha$ such that
$\bv \in \Umin{\alpha}{\bv} \atens \Umin{D\setminus\{\alpha\}}{\bv}$.
\end{theorem}

\begin{proof}
Apply Proposition~\ref{prop:min_d2} to the two-factor product
$V_\alpha \atens \Vbigalpha^{[\alpha]}$. The subspace $\Umin{\alpha}{\bv}$
is the minimal subspace in the $V_\alpha$-factor. Uniqueness follows
from the same argument as in $d=2$.
\end{proof}

\begin{definition}
\label{def:minimal_subspace}
The subspace $\Umin{\alpha}{\bv}$ of Theorem~\ref{thm:min_existence}
is called the \emph{minimal subspace} of $\bv$ in direction $\alpha$.
The integer $r_\alpha(\bv) \coloneqq \dim \Umin{\alpha}{\bv}$ is the
\emph{$\alpha$-rank} of $\bv$.
\end{definition}

\subsection*{Characterisation via contractions}

The following characterisation, which is central to both the proof of
Chapter~\ref{chap:tucker_manifold} and the tree-based theory of
Chapter~\ref{chap:algTBT}, expresses $\Umin{\alpha}{\bv}$ as the
range of a family of contractions.

\begin{theorem}[Contraction characterisation]
\label{thm:min_contraction}
Let $\bv \in \VD$ and $\alpha \in D$. Then
\begin{equation}
\label{eq:min_contraction}
    \Umin{\alpha}{\bv}
    = \spann\bigl\{
        (\Id{V_\alpha} \otimes \varphi^{[\alpha]})(\bv) :
        \varphi^{[\alpha]} \in (\Vbigalpha^{[\alpha]})^*
      \bigr\},
\end{equation}
where $(\Id{V_\alpha} \otimes \varphi^{[\alpha]})(\bv) \in V_\alpha$
is the contraction of $\bv$ with the functional $\varphi^{[\alpha]}$
in all directions except $\alpha$.
\end{theorem}

\begin{proof}
Write $\bv = \sum_\nu u_\nu \otimes \bU_\nu$ with $u_\nu \in V_\alpha$
and $\bU_\nu \in \Vbigalpha^{[\alpha]}$. For $\varphi \in
(\Vbigalpha^{[\alpha]})^*$:
$(\Id{V_\alpha} \otimes \varphi)(\bv) = \sum_\nu \varphi(\bU_\nu) u_\nu
\in \spann\{u_\nu\} = \Umin{\alpha}{\bv}$. Conversely, by taking
$\varphi = \varphi_\nu$ dual to a basis of $\Umin{D\setminus\{\alpha\}}{\bv}$,
one recovers a basis of $\Umin{\alpha}{\bv}$, so equality holds.
\end{proof}

\begin{remark}
\label{rem:min_replacement}
In Theorem~\ref{thm:min_contraction}, one can replace
$(\Vbigalpha^{[\alpha]})^*$ by $(\Umin{D\setminus\{\alpha\}}{\bv})^*$
without changing the span, since $\bv \in V_\alpha \atens
\Umin{D\setminus\{\alpha\}}{\bv}$. This gives the tighter formula
\begin{equation}
\label{eq:min_contraction_minimal}
    \Umin{\alpha}{\bv}
    = \spann\bigl\{
        (\Id{V_\alpha} \otimes \varphi^{[\alpha]})(\bv) :
        \varphi^{[\alpha]} \in (\Umin{D\setminus\{\alpha\}}{\bv})^*
      \bigr\},
\end{equation}
which is often more useful in practice.
\end{remark}

%------------------------------------------------------------
\section{Lattice properties and inclusion under partitions}
\label{sec:min:lattice}
%------------------------------------------------------------

\subsection*{Monotonicity}

\begin{proposition}[Monotonicity of minimal subspaces]
\label{prop:min_monotone}
Let $\bv \in \VD$ and $U_\alpha \in \Grass{r_\alpha}{V_\alpha}$
for each $\alpha \in D$. Then:
\begin{equation*}
    \bv \in {}_a\bigotimes_{\alpha \in D} U_\alpha
    \quad \Longrightarrow \quad
    \Umin{\alpha}{\bv} \subset U_\alpha \quad \forall\, \alpha \in D.
\end{equation*}
\end{proposition}

\begin{proof}
Under $\bv \in {}_a\bigotimes_{\alpha \in D} U_\alpha$, the identification
$\VD \cong V_\alpha \atens \Vbigalpha^{[\alpha]}$ gives
$\bv \in U_\alpha \atens \Vbigalpha^{[\alpha]}$,
so the minimality of $\Umin{\alpha}{\bv}$ (Theorem~\ref{thm:min_existence}(b))
yields $\Umin{\alpha}{\bv} \subset U_\alpha$.
\end{proof}

\subsection*{Inclusion under partition refinement}

The following proposition is the key algebraic fact underlying the
tree-based format: it shows that minimal subspaces at a coarser
level of a tree are contained in tensor products of minimal subspaces
at finer levels.

\begin{proposition}[Inclusion under refinement]
\label{prop:min_inclusion_partition}
Let $\bv \in \VD$ and let $\alpha \subset D$ with $\card{\alpha} \geq 2$.
Let $\mathcal{P}_\alpha$ be a non-trivial partition of $\alpha$. Then
\begin{equation}
\label{eq:min_inclusion}
    \Umin{\alpha}{\bv} \subset
    {}_a\bigotimes_{\beta \in \mathcal{P}_\alpha} \Umin{\beta}{\bv}.
\end{equation}
\end{proposition}

\begin{proof}
From the canonical identification $\VD \cong \Vbigalpha \atens
\Vbigalpha^{[\alpha]}$ and Theorem~\ref{thm:min_contraction},
\begin{equation*}
    \Umin{\alpha}{\bv} =
    \spann\bigl\{
        (\Id{\Vbigalpha} \otimes \varphi^{[\alpha]})(\bv)
        : \varphi^{[\alpha]} \in (\Umin{\ac}{\bv})^*
    \bigr\}.
\end{equation*}
For each such element $w_\alpha = (\Id{\Vbigalpha} \otimes
\varphi^{[\alpha]})(\bv) \in \Vbigalpha$, we apply
Proposition~\ref{prop:min_monotone} to the tensor space
$\Vbigalpha \cong {}_a\bigotimes_{\beta \in \mathcal{P}_\alpha} \Vbigbeta$.
Since $w_\alpha \in \Umin{\alpha}{\bv}$ and
$\Vbigalpha = {}_a\bigotimes_{\beta \in \mathcal{P}_\alpha} \Vbigbeta$,
one has $w_\alpha \in {}_a\bigotimes_{\beta \in \mathcal{P}_\alpha}
\Umin{\beta}{w_\alpha}$. Moreover, from the minimality of
$\Umin{\beta}{\bv}$ and the relation between contractions at different
levels, $\Umin{\beta}{w_\alpha} \subset \Umin{\beta}{\bv}$ for each
$\beta \in \mathcal{P}_\alpha$. Hence
$w_\alpha \in {}_a\bigotimes_{\beta \in \mathcal{P}_\alpha} \Umin{\beta}{\bv}$,
and taking the span over all such $w_\alpha$ gives~\eqref{eq:min_inclusion}.
\end{proof}

\begin{remark}
\label{rem:min_inclusion_tree}
Proposition~\ref{prop:min_inclusion_partition} applied iteratively at
every internal vertex $\alpha \in \TD$ of a dimension partition tree
(Chapter~\ref{chap:trees}) gives the fundamental compatibility of
minimal subspaces with the tree structure:
\begin{equation}
\label{eq:min_inclusion_tree}
    \Umin{\alpha}{\bv} \subset
    {}_a\bigotimes_{\beta \in \Sons{\alpha}} \Umin{\beta}{\bv}
    \quad \forall\, \alpha \in \TD \setminus \Leaves.
\end{equation}
This is the algebraic foundation of the tree-based format in
Chapter~\ref{chap:algTBT}.
\end{remark}

\subsection*{New characterisation of minimal subspaces}

The following result, proved in~\cite{FHN2021} and used in the construction
of tree-based tensor formats, gives a characterisation of $\Umin{\beta}{\bv}$
via contractions of $\Umin{\alpha}{\bv}$ for $\alpha \supset \beta$.

\begin{proposition}
\label{prop:min_char_partition}
Let $\bv \in \VD$, $\alpha \subset D$ with $\card{\alpha} \geq 2$,
and $\mathcal{P}_\alpha$ a non-trivial partition of $\alpha$. Suppose
that $\Vbigalpha$ and $\Vbigbeta$ for $\beta \in \mathcal{P}_\alpha$
are normed spaces. Then for each $\beta \in \mathcal{P}_\alpha$:
\begin{equation}
\label{eq:min_char_partition}
    \Umin{\beta}{\bv}
    = \spann\bigl\{
        (\Id{\Vbigbeta} \otimes \varphi^{(\alpha\setminus\beta)})
        (\bv_\alpha) :
        \bv_\alpha \in \Umin{\alpha}{\bv},\;
        \varphi^{(\alpha\setminus\beta)} \in
        ({}_a\bigotimes_{\gamma \in \mathcal{P}_\alpha \setminus \{\beta\}}
        \Vbigalpha[\gamma])^*
    \bigr\}.
\end{equation}
\end{proposition}

\begin{proof}
Applying Theorem~\ref{thm:min_contraction} twice — first for the
full index set $D$ and then for $\alpha$ — one obtains two nested
contraction characterisations. The formula~\eqref{eq:min_char_partition}
follows by composing these and using that $\bv \in \Vbigalpha \atens
\Umin{\ac}{\bv}$ allows the replacement of the ambient dual by that
of $\Umin{\ac}{\bv}$.
\end{proof}

%------------------------------------------------------------
\section{Admissible ranks and the disjoint decomposition}
\label{sec:min:admissible}
%------------------------------------------------------------

\subsection*{Admissible Tucker ranks}

Recall from Chapter~\ref{chap:algtensor} (Definition~\ref{def:admissible_tucker})
that the set of admissible Tucker ranks is
\begin{equation*}
    \AD{\VD} = \bigl\{
        (r_\alpha(\bv))_{\alpha \in D} \in \NNz^d : \bv \in \VD
    \bigr\}.
\end{equation*}
The following proposition characterises which tuples are admissible.

\begin{proposition}
\label{prop:admissible_characterisation}
A tuple $\fr = (r_\alpha)_{\alpha \in D} \in \NNz^d$ belongs to
$\AD{\VD}$ if and only if there exist subspaces
$U_\alpha \in \Grass{r_\alpha}{V_\alpha}$ and a tensor
$\bv \in \Tucker{\fr}{{}_a\bigotimes_{\alpha \in D} U_\alpha}$,
i.e., a tensor whose $\alpha$-unfolding has rank exactly $r_\alpha$
for all $\alpha \in D$.
\end{proposition}

\begin{proposition}
\label{prop:admissible_bounds}
For $\fr \in \AD{\VD}$:
\begin{enumerate}
\item[\textup{(i)}] $r_\alpha \leq \dim V_\alpha$ for all $\alpha \in D$.
\item[\textup{(ii)}] $r_\alpha = r_{\alpha^c} \coloneqq
    \dim \Umin{D\setminus\{\alpha\}}{\bv}$ for all $\alpha \in D$
    \textup{(}symmetry of marginal ranks\textup{)}.
\item[\textup{(iii)}] $r_\alpha \leq \prod_{\beta \neq \alpha} r_\beta$
    for all $\alpha \in D$.
\end{enumerate}
\end{proposition}

\begin{proof}
Statements~(i) and~(iii) follow from the representation formula
\eqref{eq:tucker_core} and rank properties of multi-index arrays.
Statement~(ii) follows from Proposition~\ref{prop:min_d2} applied
to $V_\alpha \atens \Vbigalpha^{[\alpha]}$.
\end{proof}

\subsection*{The disjoint decomposition}

Theorem~\ref{thm:min_existence} immediately gives the following
structural result.

\begin{theorem}
\label{thm:tucker_decomposition_full}
The algebraic tensor space $\VD$ admits the disjoint decomposition
\begin{equation}
\label{eq:tucker_disjoint_full}
    \VD = \bigsqcup_{\fr \in \AD{\VD}} \Tucker{\fr}{\VD}.
\end{equation}
Each $\Tucker{\fr}{\VD}$ is non-empty for $\fr \in \AD{\VD}$.
The zero tensor $\mathbf{0}$ belongs to $\Tucker{\mathbf{0}}{\VD}$
and each elementary tensor $\bigotimes_\alpha v_\alpha$ (with all
$v_\alpha \neq 0$) belongs to $\Tucker{\mathbf{1}}{\VD}$.
\end{theorem}

\begin{remark}[Topology of the decomposition]
\label{rem:tucker_topology}
While $\Tucker{\fr}{\VD}$ is algebraically well-defined, its
topological structure is non-trivial. The key result —
that each $\Tucker{\fr}{\VD}$ is a $\mathcal{C}^\infty$-Banach
manifold immersed in the ambient Banach tensor space $\VDnorm$ —
is established in Chapter~\ref{chap:tucker_manifold}. The immersion
requires the condition \condINJ\ on the norm $\|\cdot\|_D$
(Definition~\ref{def:crossnorm}).
\end{remark}

\subsection*{Stratification by rank}

\begin{definition}
\label{def:bounded_tucker_rank}
For $\fr \in \AD{\VD}$, the set of tensors with Tucker rank
\emph{bounded by} $\fr$ is
\begin{equation}
\label{eq:tucker_bounded}
    \Tucker{\leq\fr}{\VD}
    \coloneqq
    \bigsqcup_{\substack{\mathfrak{s} \in \AD{\VD} \\ \mathfrak{s} \leq \fr}}
    \Tucker{\mathfrak{s}}{\VD}
    =
    \bigl\{
        \bv \in \VD : r_\alpha(\bv) \leq r_\alpha \;\forall\, \alpha \in D
    \bigr\},
\end{equation}
where $\fs \leq \fr$ means $s_\alpha \leq r_\alpha$ for all $\alpha \in D$.
\end{definition}

The set $\Tucker{\leq\fr}{\VD}$ is not a linear subspace (it is not
closed under addition) but it is closed under scalar multiplication.
The key property for applications — that best approximations from
$\Tucker{\leq\fr}{\VD}$ exist in Banach tensor spaces — is proved in
Section~\ref{sec:min:bestapprox}.

%------------------------------------------------------------
\section{Best approximations in Banach tensor spaces}
\label{sec:min:bestapprox}
%------------------------------------------------------------

In numerical applications, one needs to approximate a given tensor
$\bv \in \VDnorm$ by a tensor of bounded rank. The existence of such
a best approximation is non-trivial: the set $\Tucker{\leq\fr}{\VD}$
is not convex, so general theorems about proximinal sets do not apply.
The proof uses the weak closedness of $\Tucker{\leq\fr}{\VD}$.

\subsection*{Crossnorms and the injective norm condition}

\begin{definition}
\label{def:crossnorm}
Let $(V_\alpha, \|\cdot\|_\alpha)$ be normed spaces and $\|\cdot\|_D$
a norm on $\VD$. We say $\|\cdot\|_D$ is a \emph{crossnorm} if
\begin{equation}
\label{eq:crossnorm}
    \Bigl\|\bigotimes_{\alpha \in D} v_\alpha\Bigr\|_D
    = \prod_{\alpha \in D} \|v_\alpha\|_\alpha
    \quad \forall\, v_\alpha \in V_\alpha.
\end{equation}
The \emph{injective norm} $\injnorm$ on $\VD$ is defined by
\begin{equation*}
    \|\bv\|_\varepsilon
    \coloneqq \sup\bigl\{
        |(\bigotimes_{\alpha \in D} \varphi_\alpha)(\bv)| :
        \varphi_\alpha \in \dual{V_\alpha},\;
        \|\varphi_\alpha\|_{\dual{V_\alpha}} \leq 1
    \bigr\}.
\end{equation*}
We say $\|\cdot\|_D$ satisfies condition \condINJ\ if
$\|\bv\|_D \geq \|\bv\|_\varepsilon$ for all $\bv \in \VD$.
\end{definition}

A crossnorm satisfies \condINJ\ if and only if it is a
\emph{reasonable crossnorm} in the sense that its dual norm is also
a crossnorm. The injective norm is the smallest reasonable crossnorm
(see \cite{Hackbusch2019}, Chapter~4).

\begin{remark}
\label{rem:crossnorm_examples}
The $L^p$ norm on $L^p(\prod_\alpha I_\alpha)$ is a reasonable
crossnorm for $1 \leq p \leq \infty$
(see \cite{Hackbusch2019}, Example~4.72).
The Sobolev norm $\|\cdot\|_{(0,1),p}$ on
$H^{1,p}(I_1) \atens H^{1,p}(I_2)$ is a crossnorm
(Example~4.42 in \cite{Hackbusch2019}), hence satisfies \condINJ.
\end{remark}

\subsection*{Weak closedness of bounded-rank sets}

\begin{theorem}[Weak closedness, {\cite[Theorem~5.1]{FH2012}}]
\label{thm:weak_closed}
Let $(V_\alpha, \|\cdot\|_\alpha)$ be normed spaces and let
$\|\cdot\|_D$ be a norm on $\VD$ satisfying condition \condINJ.
Then the set $\Tucker{\leq\fr}{\VD}$ is weakly closed in $\VDnorm$.
\end{theorem}

\begin{proof}[Sketch]
Let $(\bv_n) \subset \Tucker{\leq\fr}{\VD}$ converge weakly to
$\bv \in \VDnorm$. We must show $\bv \in \Tucker{\leq\fr}{\VD}$,
i.e., $r_\alpha(\bv) \leq r_\alpha$ for all $\alpha \in D$.

Fix $\alpha \in D$ and consider the $\alpha$-unfolding map
$\Mmat{\alpha}: \VDnorm \to \cL(\dual{(V_\alpha)}, \Vbigalpha^{[\alpha]})$
defined by $\Mmat{\alpha}(\bv)(\varphi) = (\varphi \otimes \Id{})(\bv)$.
Since $r_\alpha(\bv_n) \leq r_\alpha$, we have
$\rank \Mmat{\alpha}(\bv_n) \leq r_\alpha$ for all $n$. The condition
\condINJ\ implies that $\Mmat{\alpha}$ is continuous for the weak
topology on $\VDnorm$ and the weak operator topology on
$\cL(\dual{(V_\alpha)}, \Vbigalpha^{[\alpha]})$. Since the set of
operators of rank $\leq r_\alpha$ is closed in the weak operator topology
(it is a closed variety), the limit $\Mmat{\alpha}(\bv)$ has rank
$\leq r_\alpha$, i.e., $r_\alpha(\bv) \leq r_\alpha$.
\end{proof}

\subsection*{Existence of best approximations}

\begin{theorem}[Best approximation, {\cite[Theorem~5.3]{FH2012}}]
\label{thm:best_approx_tucker}
Let $(V_\alpha, \|\cdot\|_\alpha)$ be normed spaces, let $\|\cdot\|_D$
be a norm on $\VD$ satisfying \condINJ, and assume that $\VDnorm$ is
reflexive. Then for every $\bu \in \VDnorm$ and $\fr \in \AD{\VD}$,
there exists $\bv_* \in \Tucker{\leq\fr}{\VDnorm}$ such that
\begin{equation}
\label{eq:best_approx}
    \|\bu - \bv_*\|_D = \inf_{\bv \in \Tucker{\leq\fr}{\VDnorm}}
    \|\bu - \bv\|_D.
\end{equation}
\end{theorem}

\begin{proof}
Since $\VDnorm$ is reflexive, any bounded sequence has a weakly
convergent subsequence (Eberlein--Šmulian theorem). Let
$(\bv_n) \subset \Tucker{\leq\fr}{\VDnorm}$ be a minimising sequence
for the infimum in~\eqref{eq:best_approx}. The sequence $(\bv_n)$ is
bounded (the infimum is attained or approached by a bounded set), so
by reflexivity there is a weakly convergent subsequence
$\bv_{n_k} \rightharpoonup \bv_*$. By
Theorem~\ref{thm:weak_closed}, $\bv_* \in \Tucker{\leq\fr}{\VDnorm}$.
Since the norm is weakly lower semicontinuous,
$\|\bu - \bv_*\|_D \leq \liminf_k \|\bu - \bv_{n_k}\|_D = \inf$,
so $\bv_*$ is a best approximation.
\end{proof}

\begin{corollary}
\label{cor:best_approx_lp}
For $I_\alpha \subset \RR$ ($\alpha \in D$) and $1 < p < \infty$,
the Banach tensor space $L^p(\prod_\alpha I_\alpha)$ is reflexive and
its norm satisfies \condINJ\ (Example~\ref{rem:crossnorm_examples}).
Hence every $f \in L^p(\prod_\alpha I_\alpha)$ has a best
Tucker approximation of any prescribed rank.
\end{corollary}

\begin{corollary}
\label{cor:best_approx_sobolev}
For $1 < p < \infty$, $N \geq 1$, and bounded intervals $I_\alpha$:
\begin{enumerate}
\item[\textup{(i)}] $H^{N,p}(\prod_\alpha I_\alpha)$ with the standard
    Sobolev norm is reflexive and satisfies \condINJ.
\item[\textup{(ii)}] The space $H^{1,p}(I_1) \otimes_{\|\cdot\|_{(0,1),p}}
    H^{1,p}(I_2)$ with the partial-derivative crossnorm satisfies \condINJ.
\end{enumerate}
In both cases, best Tucker approximations of any rank exist.
\end{corollary}

\begin{example}
\label{ex:min:sobolev}
Consider $d=2$, $V_1 = H^1(0,1)$, $V_2 = H^1(0,1)$, and the
ambient space $V_1 \otimes_{H^1} V_2$ with the norm
$\|f\|_{H^1}^2 = \|f\|_{L^2}^2 + \|\partial_{x_2} f\|_{L^2}^2$.
This Banach tensor space is reflexive and satisfies \condINJ, so every
$f \in V_1 \otimes_{H^1} V_2$ has a best approximation of
rank $\leq r$ for any $r \geq 1$. This setting arises in parametric PDEs
where $x_1$ is the spatial variable and $x_2$ is the stochastic parameter.
\end{example}

\subsection*{Extension to intersections}

In many applications (including the tree-based format), the ambient
space is an intersection of Banach tensor spaces, each with its own norm.

\begin{theorem}[Best approximation in intersections, {\cite[Theorem~5.6]{FH2012}}]
\label{thm:best_approx_intersection}
Let $\VDnorm^{(1)}, \ldots, \VDnorm^{(k)}$ be reflexive Banach tensor
spaces over the same algebraic tensor space $\VD$, each with a norm
satisfying \condINJ. Equip the intersection
$Z = \VDnorm^{(1)} \cap \cdots \cap \VDnorm^{(k)}$ with the norm
$\|\bv\|_Z = \max_{i} \|\bv\|_i$. If $Z$ is reflexive (which holds,
e.g., when $k=2$ and both norms are reflexive), then every $\bu \in Z$
has a best approximation in $\Tucker{\leq\fr}{Z}$.
\end{theorem}

%------------------------------------------------------------
\section{Worked examples}
\label{sec:min:examples}
%------------------------------------------------------------

This section illustrates the abstract theory through two concrete
examples that arise naturally in scientific computing.

\subsection*{Example 1: A Gaussian kernel in $L^2([0,1]^2)$}

\begin{example}[Gaussian kernel]
\label{ex:min:gaussian}
Let $V_1 = V_2 = L^2([0,1])$ and consider the function
\begin{equation}
\label{eq:gaussian_kernel}
    f(x_1, x_2) \coloneqq e^{-\gamma(x_1 - x_2)^2},
    \quad (x_1, x_2) \in [0,1]^2,\quad \gamma > 0.
\end{equation}
We identify $f$ with a Hilbert--Schmidt operator
$T_f: L^2([0,1]) \to L^2([0,1])$ via
$(T_f g)(x_1) = \int_0^1 f(x_1, x_2) g(x_2)\,dx_2$.

\paragraph{Minimal subspaces.}
Since $T_f$ is a compact self-adjoint positive operator
(the kernel is symmetric and positive definite), its spectral
decomposition is
\begin{equation*}
    T_f = \sum_{k=1}^\infty \lambda_k\, \phi_k \otimes \phi_k,
    \quad \lambda_1 \geq \lambda_2 \geq \cdots > 0, \quad
    \sum_k \lambda_k < \infty,
\end{equation*}
with $\{\phi_k\}_{k \geq 1}$ an orthonormal basis of $L^2([0,1])$
consisting of eigenfunctions of $T_f$. In the tensor product notation,
\begin{equation*}
    f = \sum_{k=1}^\infty \lambda_k\, \phi_k \otimes \phi_k
    \in L^2([0,1]) \otimes_H L^2([0,1]),
\end{equation*}
so $r_1(f) = r_2(f) = \infty$ (the kernel has infinite rank). For any
$r \in \NN$, however, the minimal subspaces of the rank-$r$ truncation
\begin{equation*}
    f_r \coloneqq \sum_{k=1}^r \lambda_k\, \phi_k \otimes \phi_k
\end{equation*}
are $\Umin{1}{f_r} = \Umin{2}{f_r} = \spann\{\phi_1, \ldots, \phi_r\}$,
and $f_r \in \Tucker{(r,r)}{L^2([0,1])^{\otimes 2}}$.

\paragraph{Best approximation.}
By the Eckart--Young theorem (the $d=2$ case of
Theorem~\ref{thm:best_approx_tucker}), $f_r$ is the
best approximation of $f$ of Tucker rank $(r,r)$ in the Hilbert tensor
norm:
\begin{equation*}
    \|f - f_r\|_{L^2([0,1]^2)}^2
    = \sum_{k > r} \lambda_k^2
    = \inf_{\bg \in \Tucker{\leq(r,r)}{L^2}} \|f - \bg\|^2.
\end{equation*}

\paragraph{Decay of eigenvalues.}
For $\gamma > 0$, the eigenvalues of the Gaussian kernel
decay super-algebraically: $\lambda_k \sim C e^{-c\sqrt{k}}$
for constants $C, c > 0$ depending on $\gamma$. This rapid decay
means that $f_r$ approximates $f$ to any prescribed accuracy with
$r$ growing only polylogarithmically with the tolerance. More
precisely, for tolerance $\varepsilon > 0$:
\begin{equation}
\label{eq:gaussian_rank}
    \|f - f_r\|_{L^2}^2 \leq \varepsilon^2
    \quad \text{requires } r = O\bigl(\log(1/\varepsilon)^2\bigr).
\end{equation}
This is the mathematical reason why low-rank approximations work
exceptionally well for functions with smooth integral kernels: the
minimal subspace dimension needed for $\varepsilon$-accuracy grows only
logarithmically with $1/\varepsilon$.
\end{example}

\subsection*{Example 2: Parametric partial differential equations}

The main motivation for the tensor framework of this monograph is the
numerical treatment of high-dimensional PDEs. The following example
shows how the minimal subspace dimension encodes the sensitivity of the
solution with respect to parameters.

\begin{example}[Parametric elliptic PDE]
\label{ex:min:parametric_pde}
Let $D_x = (0,1) \subset \RR$ and $\Gamma = [-1,1]^s \subset \RR^s$
be a parameter space ($s \geq 1$). Consider the parametric family of
elliptic problems: find $u(\cdot, \bm{y}): D_x \to \RR$ such that
\begin{equation}
\label{eq:parametric_pde}
    -\frac{\partial}{\partial x}
    \Bigl[a(x, \bm{y})\,\frac{\partial u}{\partial x}(x, \bm{y})\Bigr]
    = f(x),
    \quad x \in D_x,\; \bm{y} = (y_1, \ldots, y_s) \in \Gamma,
\end{equation}
with homogeneous Dirichlet boundary conditions $u(0,\bm{y}) =
u(1,\bm{y}) = 0$. The diffusion coefficient has the affine form
\begin{equation*}
    a(x, \bm{y}) = a_0(x) + \sum_{j=1}^s y_j a_j(x),
\end{equation*}
where $a_0, a_1, \ldots, a_s \in L^\infty(D_x)$ with
$a(x, \bm{y}) \geq a_{\min} > 0$ uniformly.

\paragraph{Tensor product setting.}
The solution map $\bm{y} \mapsto u(\cdot, \bm{y})$ defines an element
\begin{equation*}
    \bm{u} \in \mathcal{V} \coloneqq
    H_0^1(D_x) \otimes L^2(\Gamma, \pi),
\end{equation*}
where $\pi$ is the uniform measure on $\Gamma$ and the tensor product
norm satisfies \condINJ\ (Corollary~\ref{cor:best_approx_sobolev}).

\paragraph{Minimal subspaces and parameter sensitivity.}
The $\alpha$-rank $r_\alpha(\bm{u})$ for $\alpha = \{x\}$ (the physical
variable) equals the effective rank of the solution manifold
$\{u(\cdot, \bm{y}) : \bm{y} \in \Gamma\} \subset H_0^1(D_x)$:
\begin{equation*}
    \Umin{x}{\bm{u}}
    = \overline{\spann\{u(\cdot, \bm{y}) : \bm{y} \in \Gamma\}}
    \subset H_0^1(D_x).
\end{equation*}
If $a_0$ dominates (i.e., the perturbations $y_j a_j$ are small),
the solution $u(\cdot, \bm{y})$ varies little with $\bm{y}$, and
$r_x(\bm{u})$ is small. More precisely, the Kolmogorov $n$-width
\begin{equation*}
    d_n\bigl(\{u(\cdot,\bm{y}) : \bm{y}\in\Gamma\},\,H_0^1(D_x)\bigr)
    \coloneqq \inf_{\substack{W \subset H_0^1(D_x) \\ \dim W = n}}
    \sup_{\bm{y} \in \Gamma}
    \inf_{w \in W} \|u(\cdot,\bm{y}) - w\|_{H_0^1}
\end{equation*}
decays at a rate determined by the smoothness of $a$ in the parameters
$\bm{y}$. When $a$ is analytic in $\bm{y}$ (which holds when the
$a_j$ are bounded away from zero), the $n$-width decays exponentially:
$d_n \leq C e^{-cn^{1/s}}$ for constants $C, c > 0$
(see~\cite{Nouy2017}).

\paragraph{Best approximation from Tucker format.}
The best Tucker-$(r,r)$ approximation of $\bm{u}$ in $\mathcal{V}$
is
\begin{equation*}
    \bm{u}_r = \sum_{k=1}^r \lambda_k\, v_k^{(x)} \otimes v_k^{(\bm{y})},
    \quad
    \|\bm{u} - \bm{u}_r\|_{\mathcal{V}}^2 = \sum_{k > r} \lambda_k^2,
\end{equation*}
where $\lambda_k$ are the singular values of the solution operator
$\bm{y} \mapsto u(\cdot, \bm{y})$ as a Hilbert--Schmidt operator. The
existence of this best approximation is guaranteed by
Corollary~\ref{cor:best_approx_sobolev}. The spatial basis functions
$v_k^{(x)} \in H_0^1(D_x)$ are the principal components of the
solution manifold, and $v_k^{(\bm{y})} \in L^2(\Gamma)$ encode the
parameter dependence. This decomposition is the starting point for
the \emph{proper orthogonal decomposition} (POD) and reduced basis
methods used in uncertainty quantification; see~\cite{Nouy2017}.
\end{example}

%------------------------------------------------------------
\section{Historical notes and comparison with the literature}
\label{sec:min:history}
%------------------------------------------------------------

The notion of minimal subspace was introduced in~\cite{FH2012}.
It generalises to infinite dimensions the classical concept of matrix
rank: for $d=2$ and finite-dimensional $V_\alpha$, the $\alpha$-rank
$r_\alpha(\bv) = \rank \Mmat{\alpha}(\bv)$ coincides with the classical
matrix rank of the $\alpha$-unfolding, introduced by
Hitchcock~\cite{Hitchcock1927} in 1927. The higher-order version for
finite-dimensional spaces was studied by Tucker~\cite{Tucker1966},
De Lathauwer et al.~\cite{DeLathauwer2000}, and Kolda--Bader~\cite{KoldaBader2009}.

The infinite-dimensional theory of~\cite{FH2012} is motivated by the
numerical treatment of high-dimensional PDEs, where $V_\alpha = L^2(\RR^3)$
or Sobolev spaces, and the ambient tensor space is infinite-dimensional.
The existence of best approximations (Theorem~\ref{thm:best_approx_tucker})
removes a fundamental obstacle in the convergence analysis of alternating
least-squares and related iterative methods.

The weak closedness argument of Theorem~\ref{thm:weak_closed} is
modelled on Uschmajew's result for the case of the spectral norm
on finite-dimensional spaces~\cite{Uschmajew2012}. The innovation
of~\cite{FH2012} is the extension to general reflexive Banach spaces
under the crossnorm condition \condINJ, which is the natural condition
for tensor product norms.

The connection between minimal subspaces and the geometry of Tucker
manifolds is developed in Chapter~\ref{chap:tucker_manifold}, where
the rank map $\varrho: \VD \to \prod_\alpha \Grass{r_\alpha}{V_\alpha}$
defined by $\varrho(\bv) = (\Umin{\alpha}{\bv})_{\alpha \in D}$ plays
a central role in the construction of local charts.



%%  PART II: TREE-BASED TENSOR FORMATS
\part{Tree-Based Tensor Formats}
\label{part:tree}

%============================================================
%  ch05_trees.tex  —  Dimension Partition Trees
%  TODO: Write content
%============================================================

\chapter{Dimension Partition Trees}
\label{chap:trees}

\section*{To be written}

This chapter is a placeholder. Content will be added in subsequent
revisions.


%============================================================
%  ch06_algTBT.tex  —  Algebraic Tree-Based Tensor Formats
%  TODO: Write content
%============================================================

\chapter{Algebraic Tree-Based Tensor Formats}
\label{chap:algTBT}

\section*{To be written}

This chapter is a placeholder. Content will be added in subsequent
revisions.


%============================================================
%  ch07_topTBT.tex  —  Topological Tree-Based Tensor Spaces
%  TODO: Write content
%============================================================

\chapter{Topological Tree-Based Tensor Spaces}
\label{chap:topTBT}

\section*{To be written}

This chapter is a placeholder. Content will be added in subsequent
revisions.



%%  PART III: DIFFERENTIAL GEOMETRY OF TENSOR FORMATS
\part{Differential Geometry of Tensor Formats}
\label{part:geometry}

%============================================================
%  ch08_tucker_manifold.tex
%  Chapter 8: The Tucker Manifold
%
%  Tensor Formats in Banach Spaces
%  Antonio Falcó, CEU Universities
%
%  Based on: FHN2019 (Secs. 3-4)
%  Estimated length: ~45 pp.
%============================================================

\chapter{The Tucker Manifold}
\label{chap:tucker_manifold}

This chapter establishes the central geometric result of the monograph:
the set $\Tucker{\fr}{\VD}$ of tensors with fixed Tucker rank carries
the structure of a $\mathcal{C}^\infty$-Banach manifold, and is an
immersed submanifold of the ambient Banach tensor space $\VDnorm$.
Section~\ref{sec:tucker_manifold:decomp} introduces the rank map and
motivates the manifold structure.
Section~\ref{sec:tucker_manifold:charts} constructs the local charts
explicitly. Section~\ref{sec:tucker_manifold:cinfty} proves the
$\mathcal{C}^\infty$ atlas and the fibre bundle structure
(Theorem~\ref{thm:Tucker_Banach_Manifold}).
Section~\ref{sec:tucker_manifold:tangent} identifies the tangent space.
Section~\ref{sec:tucker_manifold:immersion} proves the immersion theorem
(Theorem~\ref{thm:tucker_immersion}).
Throughout, $D = \{1,\ldots,d\}$, $d \geq 2$, and $(V_\alpha,
\|\cdot\|_\alpha)$ are normed spaces.

%------------------------------------------------------------
\section{The rank map and the decomposition of $\VD$}
\label{sec:tucker_manifold:decomp}
%------------------------------------------------------------

Recall from Theorem~\ref{thm:tucker_decomposition_full} that
\begin{equation*}
    \VD = \bigsqcup_{\fr \in \AD{\VD}} \Tucker{\fr}{\VD}.
\end{equation*}
The map
\begin{equation}
\label{eq:rank_map}
    \varrho : \VD \longrightarrow
    \bigcup_{\fr \in \AD{\VD}}
    \prod_{\alpha \in D} \Grass{r_\alpha}{V_\alpha},
    \qquad
    \bv \mapsto \bigl(\Umin{\alpha}{\bv}\bigr)_{\alpha \in D},
\end{equation}
sends each tensor to its tuple of minimal subspaces. For fixed
$\fr \in \AD{\VD}$ with $\fr \neq \mathbf{0}$, the restriction
\begin{equation}
\label{eq:rank_map_restricted}
    \varrho_\fr \coloneqq \varrho\big|_{\Tucker{\fr}{\VD}} :
    \Tucker{\fr}{\VD} \longrightarrow
    \prod_{\alpha \in D} \Grass{r_\alpha}{V_\alpha}
\end{equation}
is surjective: for any tuple $(U_\alpha)_{\alpha \in D}$ with
$U_\alpha \in \Grass{r_\alpha}{V_\alpha}$, any full-rank tensor in
${}_a\bigotimes_\alpha U_\alpha$ belongs to $\Tucker{\fr}{\VD}$ and
maps to $(U_\alpha)$.

The product of Grassmannians $\prod_{\alpha \in D}
\Grass{r_\alpha}{V_\alpha}$ is a Banach manifold (Theorem~\ref{thm:grassmann_manifold}
applied factor by factor), and the rank map $\varrho_\fr$ is its
natural base for the bundle structure of $\Tucker{\fr}{\VD}$.

%------------------------------------------------------------
\section{Local charts}
\label{sec:tucker_manifold:charts}
%------------------------------------------------------------

Fix $\fr \in \AD{\VD}$ with $\fr \neq \mathbf{0}$ and
$\bv \in \Tucker{\fr}{\VD}$. Let $V_{\alpha,\|\cdot\|_\alpha}$
denote the completion of $V_\alpha$ and choose for each $\alpha \in D$
a closed complement $\Wmin{\alpha}{\bv}$ of $\Umin{\alpha}{\bv}$ in
$V_{\alpha,\|\cdot\|_\alpha}$:
\begin{equation*}
    V_{\alpha,\|\cdot\|_\alpha}
    = \Umin{\alpha}{\bv} \oplus \Wmin{\alpha}{\bv}.
\end{equation*}
Such a complement exists by Proposition~\ref{prop:finitedim_split},
since $\Umin{\alpha}{\bv}$ is finite-dimensional.

\subsection*{Chart domain}

Define the \emph{chart domain at $\bv$} by
\begin{equation}
\label{eq:chart_domain}
    \chart(\bv) \coloneqq
    \bigl\{
        \bw \in \Tucker{\fr}{\VD} :
        \Umin{\alpha}{\bw} \in
        \GrassB{\Wmin{\alpha}{\bv}, V_\alpha},\;
        \forall \alpha \in D
    \bigr\},
\end{equation}
where $\GrassB{W, X} = \{U \in \GrassB{X} : X = U \oplus W\}$ is
the chart domain of the Grassmannian at $\Umin{\alpha}{\bv}$
(see Section~\ref{sec:banach:grassmann}).

\subsection*{Chart map}

Fix bases $\{u^{(\alpha)}_{i_\alpha}\}_{i_\alpha=1}^{r_\alpha}$ of
$\Umin{\alpha}{\bv}$ for each $\alpha \in D$. By
Lemma~\ref{lem:tucker_characterisation}, there is a unique
$\CD(\bv) \in \RR^{\times r_\alpha}_*$ such that $\bv$ has the
Tucker representation~\eqref{eq:tucker_core}.
For each $\bw \in \chart(\bv)$, define the \emph{chart map}
\begin{equation}
\label{eq:chart_map}
    \chartmap{\bv} : \chart(\bv)
    \longrightarrow
    \prod_{\alpha \in D} \cL\bigl(\Umin{\alpha}{\bv},
        \Wmin{\alpha}{\bv}\bigr)
    \times \RR^{\times r_\alpha}_*,
\end{equation}
by sending $\bw = (\bigotimes_\alpha \exp(L_\alpha))(\bu(E^{(D)}))$
to the pair $((L_\alpha)_{\alpha \in D},\, E^{(D)})$, where
$(L_\alpha, E^{(D)})$ are uniquely determined by
Lemma~\ref{lem:element} below.

\begin{lemma}[Local representation]
\label{lem:element}
Under the conditions above, for $\bw \in \Tucker{\fr}{\VD}$
the following are equivalent:
\begin{enumerate}
\item[\textup{(a)}] $\bw \in \chart(\bv)$.
\item[\textup{(b)}] There exists a unique pair
    $\bigl((L_\alpha)_{\alpha \in D},\, E^{(D)}\bigr)
    \in \prod_\alpha \cL\bigl(\Umin{\alpha}{\bv}, \Wmin{\alpha}{\bv}\bigr)
    \times \RR^{\times r_\alpha}_*$
    such that
    \begin{equation}
    \label{eq:local_rep}
        \bw
        = \Bigl(\bigotimes_{\alpha \in D} \exp(L_\alpha)\Bigr)
          \bigl(\bu(E^{(D)})\bigr),
    \end{equation}
    where $\bu(E^{(D)}) \in \Tucker{\fr}\bigl({}_a\bigotimes_\alpha
    \Umin{\alpha}{\bv}\bigr)$ is the tensor with core $E^{(D)}$ in
    the fixed bases $\{u^{(\alpha)}_{i_\alpha}\}$.
\end{enumerate}
\end{lemma}

\begin{proof}
We prove (a)$\Rightarrow$(b) and (b)$\Rightarrow$(a).

\medskip
\noindent\textbf{(a) $\Rightarrow$ (b): Existence and uniqueness.}

Suppose $\bw \in \chart(\bv)$, so $\Umin{\alpha}{\bw} \in
\GrassB{\Wmin{\alpha}{\bv}, V_\alpha}$ for every $\alpha \in D$.

\textit{Step 1: Recovering $L_\alpha$.}
Define $L_\alpha$ as the Grassmannian chart coordinate of
$\Umin{\alpha}{\bw}$ at $\Umin{\alpha}{\bv}$:
\begin{equation*}
    L_\alpha
    \coloneqq \Psi_{\Umin{\alpha}{\bv} \oplus \Wmin{\alpha}{\bv}}
    \bigl(\Umin{\alpha}{\bw}\bigr)
    = \Proj{\Wmin{\alpha}{\bv}}{\Umin{\alpha}{\bv}}
      \big|_{\Umin{\alpha}{\bw}}
      \circ
      \bigl(\Proj{\Umin{\alpha}{\bv}}{\Wmin{\alpha}{\bv}}
      \big|_{\Umin{\alpha}{\bw}}\bigr)^{-1}
    \in \cL\bigl(\Umin{\alpha}{\bv}, \Wmin{\alpha}{\bv}\bigr).
\end{equation*}
By Proposition~\ref{prop:nilpotent_algebra},
$\exp(L_\alpha) = \Id{V_\alpha} + L_\alpha$, and
\begin{equation}
\label{eq:exp_maps_U_to_Uw}
    \exp(L_\alpha)\bigl(\Umin{\alpha}{\bv}\bigr)
    = G(L_\alpha)
    = \Umin{\alpha}{\bw},
\end{equation}
since $G(L_\alpha)$ is the Grassmann chart inverse of $L_\alpha$
(Theorem~\ref{thm:grassmann_manifold}).

\textit{Step 2: Recovering $E^{(D)}$.}
Set $\widetilde{\bw} \coloneqq
\bigl(\bigotimes_{\alpha \in D} \exp(-L_\alpha)\bigr)(\bw)$.
Since $\exp(-L_\alpha) = \Id{} - L_\alpha = (\exp(L_\alpha))^{-1}$
maps $\Umin{\alpha}{\bw}$ back to $\Umin{\alpha}{\bv}$
by~\eqref{eq:exp_maps_U_to_Uw}, we get
$\Umin{\alpha}{\widetilde\bw} \subset \Umin{\alpha}{\bv}$ for each
$\alpha$. Proposition~\ref{prop:min_monotone} gives
$\widetilde\bw \in {}_a\bigotimes_\alpha \Umin{\alpha}{\bv}$,
and since $\exp(-L_\alpha)$ preserves ranks,
$\widetilde\bw \in \Tucker{\fr}({}_a\bigotimes_\alpha \Umin{\alpha}{\bv})$.
Its core array in the fixed bases is the unique $E^{(D)} \in
\RR^{\times r_\alpha}_*$ such that $\bu(E^{(D)}) = \widetilde\bw$.
Applying $\bigotimes_\alpha \exp(L_\alpha)$ recovers~\eqref{eq:local_rep}.

\textit{Uniqueness.}
If $(L_\alpha, E^{(D)})$ and $(L'_\alpha, E'^{(D)})$ both represent $\bw$,
applying $\bigotimes_\alpha \exp(-L'_\alpha)$ and Step~2 gives
$E^{(D)} = E'^{(D)}$, and then $L_\alpha = L'_\alpha$ for each $\alpha$.

\medskip
\noindent\textbf{(b) $\Rightarrow$ (a).}

For $(L_\alpha, E^{(D)})$ with $E^{(D)} \in \RR^{\times r_\alpha}_*$,
let $\bw$ be defined by~\eqref{eq:local_rep}. By~\eqref{eq:exp_maps_U_to_Uw},
$\Umin{\alpha}{\bw} = G(L_\alpha) \in \GrassB{\Wmin{\alpha}{\bv}, V_\alpha}$,
so $\bw \in \chart(\bv)$.
\end{proof}

Here $\exp(L_\alpha) = \Id{V_\alpha} + L_\alpha$ by the nilpotency
of $L_\alpha \in \cL_{(\Umin{\alpha}{\bv},\Wmin{\alpha}{\bv})}(V_\alpha)$
(Proposition~\ref{prop:nilpotent_algebra}), so the representation is
polynomial and $\mathcal{C}^\infty$.

\begin{remark}
\label{rem:chart_geometric}
Formula~\eqref{eq:local_rep} has a clean geometric interpretation:
$\bw$ is obtained from the ``reference'' tensor $\bu(E^{(D)})$
(which lives in the fixed subspaces $\Umin{\alpha}{\bv}$) by applying
the ``rotation'' $\bigotimes_\alpha \exp(L_\alpha)$ which moves each
subspace $\Umin{\alpha}{\bv}$ to $\Umin{\alpha}{\bw} = G(L_\alpha)
= \spann\{(\Id{} + L_\alpha)(u^{(\alpha)}_{i_\alpha})\}$.
\end{remark}

%------------------------------------------------------------
\section{The $\mathcal{C}^\infty$-Banach manifold structure}
\label{sec:tucker_manifold:cinfty}
%------------------------------------------------------------

\begin{lemma}[Smooth transitions]
\label{lem:premanifold}
Under the continuity assumption on the tensor product map
\eqref{eq:tensor_product_normed}, for $\bv, \bv' \in
\Tucker{\fr}{\VD}$ with $\chart(\bv) \cap \chart(\bv') \neq \emptyset$,
the transition map
\begin{equation*}
    \chartmap{\bv'} \circ \chartmap{\bv}^{-1} :
    \chartmap{\bv}\bigl(\chart(\bv) \cap \chart(\bv')\bigr)
    \longrightarrow
    \chartmap{\bv'}\bigl(\chart(\bv) \cap \chart(\bv')\bigr)
\end{equation*}
is $\mathcal{C}^\infty$-Fréchet differentiable \textup{(}analytic if
the $V_\alpha$ are complex Banach spaces\textup{)}.
\end{lemma}

\begin{proof}
Let $\bw \in \chart(\bv) \cap \chart(\bv')$ with chart coordinates
$\chartmap{\bv}(\bw) = ((L_\alpha), E^{(D)})$ and
$\chartmap{\bv'}(\bw) = ((L'_\alpha), E'^{(D)})$.
We show $(L_\alpha, E^{(D)}) \mapsto (L'_\alpha, E'^{(D)})$ is
$\mathcal{C}^\infty$ (analytic over $\CC$).

\medskip
\noindent\textbf{Step 1: Smoothness of the Grassmannian component.}

For each $\alpha \in D$, $\Umin{\alpha}{\bw} = G(L_\alpha)$
(the inverse Grassmann chart of $L_\alpha$ at $\Umin{\alpha}{\bv}$).
The map
\begin{equation*}
    L_\alpha \mapsto L'_\alpha
    = \Psi_{\Umin{\alpha}{\bv'} \oplus \Wmin{\alpha}{\bv'}}(G(L_\alpha))
\end{equation*}
is the transition map of $\Grass{r_\alpha}{V_\alpha}$ between the
charts at $\Umin{\alpha}{\bv}$ and $\Umin{\alpha}{\bv'}$. Since
$\Grass{r_\alpha}{V_\alpha}$ is an analytic Banach manifold
(Theorem~\ref{thm:grassmann_manifold}), its transition maps are
analytic, hence $L_\alpha \mapsto L'_\alpha$ is analytic.
More explicitly: $G(L_\alpha) = \{(\Id{}+L_\alpha)(u) : u \in \Umin{\alpha}{\bv}\}$
is affine-linear in $L_\alpha$; the projections in the chart formula
\eqref{eq:grassmann_chart} are bounded linear; and inversion of a
bounded linear isomorphism is analytic (locally given by the Neumann
series). Their composition is analytic.

\medskip
\noindent\textbf{Step 2: Smoothness of the core component.}

Given $(L_\alpha)$ (and hence $\Umin{\alpha}{\bw} = G(L_\alpha)$),
the core is determined by
\begin{equation*}
    \bu(E^{(D)})
    = \Bigl(\bigotimes_{\alpha \in D} \exp(-L_\alpha)\Bigr)(\bw),
\end{equation*}
where $\exp(-L_\alpha) = \Id{} - L_\alpha$
(Proposition~\ref{prop:nilpotent_algebra}).
The map $(L_\alpha, \bw) \mapsto (\bigotimes_\alpha (\Id{}-L_\alpha))(\bw)$
is multilinear and continuous in $(L_\alpha)$ and $\bw$, hence
$\mathcal{C}^\infty$ by Proposition~\ref{prop:multilinear_smooth}
(analytic over $\CC$ by Corollary~\ref{cor:multilinear_analytic}).
Reading off the core array $E^{(D)}$ from the result is a bounded
linear operation. The same applies to $E'^{(D)}$ using $(L'_\alpha)$
from Step~1.

\medskip
\noindent\textbf{Conclusion.}

The transition map $(L_\alpha, E^{(D)}) \mapsto (L'_\alpha, E'^{(D)})$
is a composition of $\mathcal{C}^\infty$ (analytic over $\CC$) maps,
hence $\mathcal{C}^\infty$ (analytic over $\CC$).
\end{proof}

\begin{theorem}[$\mathcal{C}^\infty$ Tucker manifold,
{\cite[Theorem~3.17]{FHN2019}}]
\label{thm:Tucker_Banach_Manifold}
Let $(V_\alpha, \|\cdot\|_\alpha)$ be normed spaces for $\alpha \in D$
and let $\|\cdot\|_D$ be a norm on $\VD$ such that the tensor product
map~\eqref{eq:tensor_product_normed} is continuous. Then:
\begin{enumerate}
\item[\textup{(i)}] The collection $\{(\chart(\bv), \chartmap{\bv})\}_{\bv \in
    \Tucker{\fr}{\VD}}$ is a $\mathcal{C}^\infty$-atlas on
    $\Tucker{\fr}{\VD}$, which is thus a $\mathcal{C}^\infty$-Banach
    manifold modelled on the Banach space
    \begin{equation}
    \label{eq:tucker_model_space}
        \Emodel{D} \coloneqq
        \prod_{\alpha \in D}
        \cL\bigl(\Umin{\alpha}{\bv}, \Wmin{\alpha}{\bv}\bigr)
        \times \RR^{\times_{\alpha \in D} r_\alpha}.
    \end{equation}
\item[\textup{(ii)}] The triple $(\Tucker{\fr}{\VD},\,
    \prod_{\alpha \in D}\Grass{r_\alpha}{V_\alpha},\, \varrho_\fr)$
    is a $\mathcal{C}^\infty$-fibre bundle with typical fibre
    $\RR^{\times r_\alpha}_*$.
\end{enumerate}
If the $V_\alpha$ are complex Banach spaces, the manifold is analytic.
\end{theorem}

\begin{proof}
The atlas conditions AT1--AT3 are verified as follows.
AT1: $\{\chart(\bv)\}$ covers $\Tucker{\fr}{\VD}$ since every
$\bv$ belongs to its own chart domain.
AT2: $\chartmap{\bv}$ is a bijection onto
$\prod_\alpha \cL(\Umin{\alpha}{\bv}, \Wmin{\alpha}{\bv}) \times
\RR^{\times r_\alpha}_*$ by Lemma~\ref{lem:element}.
AT3: Lemma~\ref{lem:premanifold} gives $\mathcal{C}^\infty$ transitions.
The fibre bundle structure follows because $\varrho_\fr$ acts locally
as the projection onto $\prod_\alpha \GrassB{\Wmin{\alpha}{\bv}, V_\alpha}$
(the Grassmannian chart domain), with fibre $\RR^{\times r_\alpha}_*$.
\end{proof}

\begin{remark}[Independence from the norm]
\label{rem:norm_independence}
The manifold structure of $\Tucker{\fr}{\VD}$ is independent of the
choice of $\|\cdot\|_D$: the chart maps $\chartmap{\bv}$ depend only
on the splitting $V_{\alpha,\|\cdot\|_\alpha} = \Umin{\alpha}{\bv}
\oplus \Wmin{\alpha}{\bv}$, and the $\mathcal{C}^\infty$ regularity
follows from the continuity of the tensor product map, which holds for
any crossnorm. Different norms give equivalent smooth structures.
\end{remark}

\begin{corollary}[Hilbert manifold]
\label{cor:hilbert_manifold}
If $V_{\alpha,\|\cdot\|_\alpha}$ is a Hilbert space for each
$\alpha \in D$, then $\Tucker{\fr}{\VD}$ is a $\mathcal{C}^\infty$-Hilbert
manifold modelled on
$\prod_\alpha W^{r_\alpha}_\alpha \times \RR^{\times r_\alpha}$,
where $W^{r_\alpha}_\alpha = (U_\alpha)^\perp \times \cdots$ (each
$\cL(U_\alpha, W_\alpha) \cong W_\alpha^{r_\alpha}$ via the chosen
orthonormal basis of $U_\alpha$).
\end{corollary}

\begin{example}
\label{ex:tucker_mfd_Lp}
For $1 < p < \infty$, the $L^p$ tensor space
$L^p(\prod_\alpha I_\alpha)$ satisfies the hypotheses
of Theorem~\ref{thm:Tucker_Banach_Manifold}. Hence
$\Tucker{\fr}({}_a\bigotimes_\alpha L^p(I_\alpha))$ is a
$\mathcal{C}^\infty$-Banach manifold. For $p=2$ it is a Hilbert
manifold.
\end{example}

\begin{example}
\label{ex:tucker_mfd_Sobolev}
For $V_1 = H^{1,p}(I_1)$, $V_2 = H^{1,p}(I_2)$, and either the
full Sobolev norm $\|\cdot\|_{H^{1,p}}$ or the partial norm
$\|\cdot\|_{(0,1),p}$ on $V_1 \atens V_2$, both norms make the
tensor product map continuous (see~\cite[Examples~4.41,~4.42]{Hackbusch2019}).
Hence $\Tucker{(r,r)}{H^{1,p}(I_1) \atens H^{1,p}(I_2)}$ is a
$\mathcal{C}^\infty$-Banach manifold modelled on
$\cL(U_1, W_1) \times \cL(U_2, W_2) \times \GL{\RR^r}$.
\end{example}

%------------------------------------------------------------
\section{The tangent space}
\label{sec:tucker_manifold:tangent}
%------------------------------------------------------------

\begin{definition}[Tangent subspace $Z^{(D)}(\bv)$]
\label{def:ZD}
For $\bv \in \Tucker{\fr}{\VD}$, define the linear subspace
\begin{equation}
\label{eq:ZD}
    \ZD{\bv} \coloneqq
    {}_a\bigotimes_{\alpha \in D} \Umin{\alpha}{\bv}
    \oplus
    \bigoplus_{\alpha \in D}
    \Wmin{\alpha}{\bv} \atens \Umin{D\setminus\{\alpha\}}{\bv}
    \;\subset\; \VDnorm.
\end{equation}
\end{definition}

$\ZD{\bv}$ is a direct sum of $d+1$ subspaces: the core subspace
${}_a\bigotimes_\alpha \Umin{\alpha}{\bv}$ (dimension $\prod r_\alpha$)
and $d$ ``tangent directions'' $\Wmin{\alpha}{\bv} \atens
\Umin{D\setminus\{\alpha\}}{\bv}$.

\begin{proposition}[Tangent space identification]
\label{prop:tangent_space}
Under the continuity assumption, the tangent map of the inclusion
$\mathfrak{i}: \Tucker{\fr}{\VD} \hookrightarrow \VDnorm$ at $\bv$ is
injective, and
\begin{equation*}
    T_\bv \mathfrak{i}\bigl(\mathbb{T}_\bv(\Tucker{\fr}{\VD})\bigr)
    = \ZD{\bv}.
\end{equation*}
More explicitly, for $(\dot{\mathfrak{L}}, \dot{C}^{(D)}) \in
\mathbb{T}_\bv(\Tucker{\fr}{\VD})$:
\begin{equation}
\label{eq:tangent_formula}
    T_\bv \mathfrak{i}(\dot{\mathfrak{L}}, \dot{C}^{(D)})
    = \sum_{(i_\alpha)} \dot{C}^{(D)}_{(i_\alpha)}
      \bigotimes_\alpha u^{(\alpha)}_{i_\alpha}
    + \sum_{\alpha \in D} \sum_{i_\alpha=1}^{r_\alpha}
      \dot{u}^{(\alpha)}_{i_\alpha} \otimes \bU^{(\alpha)}_{i_\alpha},
\end{equation}
where $\dot{u}^{(\alpha)}_{i_\alpha} = \dot{L}_\alpha(u^{(\alpha)}_{i_\alpha})
\in \Wmin{\alpha}{\bv}$ and $\bU^{(\alpha)}_{i_\alpha} \in
\Umin{D\setminus\{\alpha\}}{\bv}$ is given by~\eqref{eq:tucker_Ualpha}.
\end{proposition}

\begin{proof}
We compute $T_\bv \mathfrak{i}$ by differentiating
$\mathfrak{i} \circ \chartmap{\bv}^{-1}$ at the base point
$(0, C_\bv^{(D)}) = \chartmap{\bv}(\bv)$.

\medskip
\noindent\textbf{Step 1: The inverse chart map is degree-one polynomial.}

By Lemma~\ref{lem:element} and Proposition~\ref{prop:nilpotent_algebra},
$\exp(L_\alpha) = \Id{} + L_\alpha$, so:
\begin{align}
\label{eq:chart_inv_poly}
    (\mathfrak{i} \circ \chartmap{\bv}^{-1})(L_\alpha, E^{(D)})
    &= \Bigl(\bigotimes_\alpha (\Id{} + L_\alpha)\Bigr)
       \Bigl(\sum_{(i_\alpha)} E^{(D)}_{(i_\alpha)}
       \bigotimes_\alpha u^{(\alpha)}_{i_\alpha}\Bigr).
\end{align}
Expanding $\bigotimes_\alpha(\Id{}+L_\alpha)$ by multilinearity and
using the nilpotency $L_\alpha \circ L_\beta = 0$ for all $\alpha, \beta$
(Proposition~\ref{prop:nilpotent_algebra}), all terms of degree $\geq 2$
in $(L_\alpha)$ vanish:
\begin{equation*}
    \bigotimes_\alpha (\Id{} + L_\alpha)
    = \bigotimes_\alpha \Id{}
    + \sum_{\alpha \in D}
      L_\alpha \otimes \bigotimes_{\beta \neq \alpha} \Id{\beta}
    \pmod{\text{degree} \geq 2}.
\end{equation*}
Hence~\eqref{eq:chart_inv_poly} is exactly affine-linear in
$(L_\alpha, E^{(D)})$.

\medskip
\noindent\textbf{Step 2: Computing the Fréchet derivative.}

Differentiating~\eqref{eq:chart_inv_poly} at $(0, C_\bv^{(D)})$ in
direction $(\dot{L}_\alpha, \dot{C}^{(D)})$:
\begin{align*}
    D(\mathfrak{i} \circ \chartmap{\bv}^{-1})\big|_{(0,C_\bv^{(D)})}
    (\dot{L}_\alpha, \dot{C}^{(D)})
    &= \sum_{(i_\alpha)} \dot{C}^{(D)}_{(i_\alpha)}
       \bigotimes_\alpha u^{(\alpha)}_{i_\alpha}\\
    &\quad+ \sum_{\alpha \in D}
       \sum_{(i_\alpha)} C^{(D)}_{\bv,(i_\alpha)}\,
       \dot{L}_\alpha(u^{(\alpha)}_{i_\alpha})
       \otimes \bigotimes_{\beta \neq \alpha} u^{(\beta)}_{i_\beta}.
\end{align*}
Using the definition of $\bU^{(\alpha)}_{i_\alpha}$
(formula~\eqref{eq:tucker_Ualpha}), the second line equals
$\sum_\alpha \sum_{i_\alpha}
\dot{L}_\alpha(u^{(\alpha)}_{i_\alpha}) \otimes \bU^{(\alpha)}_{i_\alpha}$,
giving formula~\eqref{eq:tangent_formula} with
$\dot{u}^{(\alpha)}_{i_\alpha} = \dot{L}_\alpha(u^{(\alpha)}_{i_\alpha})
\in \Wmin{\alpha}{\bv}$.

\medskip
\noindent\textbf{Step 3: Image equals $\ZD{\bv}$, and injectivity.}

Formula~\eqref{eq:tangent_formula} shows $T_\bv\mathfrak{i}$ maps into
$\ZD{\bv}$: the first term lies in
${}_a\bigotimes_\alpha \Umin{\alpha}{\bv}$ and the $\alpha$-th summand
in the second term lies in
$\Wmin{\alpha}{\bv} \atens \Umin{D\setminus\{\alpha\}}{\bv}$.

For surjectivity: given $\bz \in \ZD{\bv}$ decomposed as
$\bz = \bz_0 + \sum_\alpha \bz_\alpha$ (cf.~Definition~\ref{def:ZD}),
write $\bz_0 = \sum_{(i_\alpha)} \dot{C}^{(D)}_{(i_\alpha)}
\bigotimes_\alpha u^{(\alpha)}_{i_\alpha}$ to determine $\dot{C}^{(D)}$,
and $\bz_\alpha = \sum_{i_\alpha} w^{(\alpha)}_{i_\alpha} \otimes
\bU^{(\alpha)}_{i_\alpha}$ with $w^{(\alpha)}_{i_\alpha} \in
\Wmin{\alpha}{\bv}$, defining $\dot{L}_\alpha$ by
$\dot{L}_\alpha(u^{(\alpha)}_{i_\alpha}) = w^{(\alpha)}_{i_\alpha}$.
Then $T_\bv\mathfrak{i}(\dot{L}_\alpha, \dot{C}^{(D)}) = \bz$.

Injectivity: if $T_\bv\mathfrak{i}(\dot{L}_\alpha, \dot{C}^{(D)}) = 0$,
then $\bz_0 = 0$ gives $\dot{C}^{(D)} = 0$, and $\bz_\alpha = 0$
for each $\alpha$ gives $\dot{L}_\alpha = 0$.
\end{proof}

%------------------------------------------------------------
\section{Immersion in the ambient Banach space}
\label{sec:tucker_manifold:immersion}
%------------------------------------------------------------

To show that $\Tucker{\fr}{\VD}$ is an immersed submanifold of
$\VDnorm$, it remains to show that $\ZD{\bv} \in \GrassB{\VDnorm}$,
i.e., that $\ZD{\bv}$ is a split subspace of the ambient Banach space.
This requires condition~\condINJ.

\begin{lemma}
\label{lem:two_faces}
Assume \condINJ\ holds. Let $\beta \in D$ and $W_\beta \in
\GrassB{V_{\beta,\|\cdot\|_\beta}}$ with $V_{\beta,\|\cdot\|_\beta}
= U_\beta \oplus W_\beta$ for some finite-dimensional $U_\beta$.
Then for every finite-dimensional subspace $U_{[\beta]} \subset
\Vbigalpha^{[\beta]}$:
\begin{equation*}
    W_\beta \atens U_{[\beta]} \in \GrassB{\VDnorm}.
\end{equation*}
\end{lemma}

\begin{proof}
We construct an explicit bounded idempotent with image
$W_\beta \atens U_{[\beta]}$.

\medskip
\noindent\textbf{Step 1: Bounded projection onto $U_{[\beta]}$.}

Since $U_{[\beta]}$ is finite-dimensional,
Proposition~\ref{prop:finitedim_split} provides a complement
$W_{[\beta]}$ with $\Vbigalpha^{[\beta]} = U_{[\beta]} \oplus W_{[\beta]}$.
Let $Q_{[\beta]} \coloneqq \Proj{U_{[\beta]}}{W_{[\beta]}}$.
The map $\overline{\Id{\beta} \otimes Q_{[\beta]}} \in \cL(\VDnorm)$
is well-defined and bounded: on elementary tensors,
$\|v_\beta \otimes Q_{[\beta]}(\bU)\|_D = \|v_\beta\|_\beta
\|Q_{[\beta]}(\bU)\|_{[\beta]} \leq \|v_\beta\|_\beta
\|Q_{[\beta]}\|\,\|\bU\|_{[\beta]}$
using \condCROSS, and the extension to $\VDnorm$ by continuity is
bounded with norm $\leq \|Q_{[\beta]}\|$.

\medskip
\noindent\textbf{Step 2: Construction of $\mathcal{P}$.}

Let $P_\beta \coloneqq \Proj{W_\beta}{U_\beta} \in
\cL(V_{\beta,\|\cdot\|_\beta})$ be bounded (since $W_\beta$ is split).
Define
\begin{equation*}
    \mathcal{P}
    \coloneqq \overline{P_\beta \otimes \Id{[\beta]}}
    \circ
    \overline{\Id{\beta} \otimes Q_{[\beta]}}
    \in \cL(\VDnorm).
\end{equation*}
This is a composition of bounded operators, hence bounded.

\textit{Idempotency.} On elementary tensors:
$\mathcal{P}^2(v_\beta \otimes \bU)
= P_\beta^2(v_\beta) \otimes Q_{[\beta]}^2(\bU)
= P_\beta(v_\beta) \otimes Q_{[\beta]}(\bU)
= \mathcal{P}(v_\beta \otimes \bU)$.
By density and continuity, $\mathcal{P}^2 = \mathcal{P}$.

\textit{Image.} On elementary tensors:
$\mathcal{P}(v_\beta \otimes \bU)
= P_\beta(v_\beta) \otimes Q_{[\beta]}(\bU)
\in W_\beta \atens U_{[\beta]}$.
Conversely, for $w_\beta \in W_\beta$ and $u \in U_{[\beta]}$:
$\mathcal{P}(w_\beta \otimes u) = w_\beta \otimes u$,
since $P_\beta w_\beta = w_\beta$ and $Q_{[\beta]} u = u$.
Hence $\mathrm{Im}\,\mathcal{P} = W_\beta \atens U_{[\beta]}$.

\medskip
By Proposition~\ref{prop:split_equiv}(b),
$W_\beta \atens U_{[\beta]} \in \GrassB{\VDnorm}$.
\end{proof}

\begin{lemma}[{\cite[Lemma~4.14]{FHN2019}}]
\label{lem:direct_sum_G}
Let $X$ be a Banach space and $U, V \in \GrassB{X}$ with
$U \cap V = \{0\}$. Then $U \oplus V \in \GrassB{X}$ and
$U \cap V \in \GrassB{X}$.
\end{lemma}

\begin{proof}
Since $U, V \in \GrassB{X}$ there exist $U', V' \in \GrassB{X}$ with
$X = U \oplus U' = V \oplus V'$. Then
$U \oplus V \oplus (U' \cap V') = (U \cap V') \oplus (V \cap U')
\oplus (U' \cap V') = X$, proving $U \oplus V \in \GrassB{X}$.
The second claim follows from $X = (U \cap V) \oplus (U \cap V')
\oplus (V \cap U') \oplus (U' \cap V')$.
\end{proof}

\begin{theorem}[Immersion theorem, {\cite[Theorem~4.15]{FHN2019}}]
\label{thm:tucker_immersion}
Assume \condINJ\ holds. Then for each $\bv \in \Tucker{\fr}{\VD}$:
\begin{equation}
\label{eq:ZD_split}
    \ZD{\bv} \in \GrassB{\VDnorm}.
\end{equation}
Consequently, $\Tucker{\fr}{\VD}$ is an immersed submanifold of
$\VDnorm$.
\end{theorem}

\begin{proof}
For each $\alpha \in D$: $\Wmin{\alpha}{\bv} \in
\GrassB{V_{\alpha,\|\cdot\|_\alpha}}$ and $\Umin{D\setminus\{\alpha\}}{\bv}$
is finite-dimensional. By Lemma~\ref{lem:two_faces},
$\Wmin{\alpha}{\bv} \atens \Umin{D\setminus\{\alpha\}}{\bv}
\in \GrassB{\VDnorm}$.
Also, ${}_a\bigotimes_\alpha \Umin{\alpha}{\bv} \in \GrassB{\VDnorm}$
(finite-dimensional, hence split by
Proposition~\ref{prop:finitedim_split}).
The $d+1$ summands in~\eqref{eq:ZD} have pairwise trivial intersection
(by the linear independence properties of minimal subspaces
and the direct sum structure), so Lemma~\ref{lem:direct_sum_G}
applied iteratively gives $\ZD{\bv} \in \GrassB{\VDnorm}$.
\end{proof}

\begin{corollary}[Best projection onto tangent space]
\label{cor:best_proj_tangent}
If $\VDnorm$ is additionally reflexive, then for any $\bv \in
\Tucker{\fr}{\VD}$ and $\dot{\bu} \in \VDnorm$, there exists
$\dot{\bv}_\mathrm{best} \in \ZD{\bv}$ such that
\begin{equation*}
    \|\dot{\bu} - \dot{\bv}_\mathrm{best}\|_D
    = \min_{\dot{\bv} \in \ZD{\bv}} \|\dot{\bu} - \dot{\bv}\|_D.
\end{equation*}
\end{corollary}

\begin{proof}
Since $\VDnorm$ is reflexive, every closed subspace is proximinal
(proximinal = best approximations exist). $\ZD{\bv}$ is closed
(it is a split subspace by Theorem~\ref{thm:tucker_immersion}).
\end{proof}

%------------------------------------------------------------
\section{The tensor space as a Banach manifold}
\label{sec:tucker_manifold:full}
%------------------------------------------------------------

\begin{corollary}[{\cite[Corollary~3.19]{FHN2019}}]
\label{cor:VD_manifold}
Under the continuity assumption on the tensor product map, the algebraic
tensor space $\VD$ itself is a $\mathcal{C}^\infty$-Banach manifold
(not modelled on a single Banach space), with atlas given by the
disjoint union of the Tucker manifold atlases over all $\fr \in \AD{\VD}$:
\begin{equation*}
    \VD = \bigsqcup_{\fr \in \AD{\VD}} \Tucker{\fr}{\VD}.
\end{equation*}
\end{corollary}

%------------------------------------------------------------
\section{Historical notes and comparison}
\label{sec:tucker_manifold:history}
%------------------------------------------------------------

The $\mathcal{C}^\infty$-Banach manifold structure of
$\Tucker{\fr}{\VD}$ in infinite-dimensional spaces was established
in~\cite{FHN2019}. Prior work had treated only finite-dimensional or
Hilbert space settings:

Koch and Lubich~\cite{KochLubich2007} proved that the Tucker manifold
is a quotient manifold in finite dimensions and used it for the
Dirac--Frenkel variational principle for the MCTDH method. Their
approach, however, uses the Hilbert structure (metric projection)
and does not extend to Banach spaces.

Uschmajew and Vandereycken~\cite{Uschmajew2012} studied the HT
manifold in finite dimensions as a quotient manifold
$\GL{r_1} \times \cdots \times \GL{r_d} \backslash \FTv / G$
for a gauge group $G$. The description in~\cite{FHN2019} is
fundamentally different: it provides an \emph{explicit atlas}
with local charts given by formula~\eqref{eq:local_rep}, valid
in the infinite-dimensional Banach setting.

The key novelties of~\cite{FHN2019} are:
\begin{enumerate}
\item The chart maps $\chartmap{\bv}$ are given explicitly in terms
    of the Grassmann chart maps $\Psi^{(\alpha)}_\bv$, making the
    geometric structure completely transparent.
\item The immersion theorem (Theorem~\ref{thm:tucker_immersion})
    holds in general Banach spaces under condition~\condINJ, without
    requiring a Hilbert or reflexive structure.
\item The fibre bundle structure identifies the typical fibre
    as $\RR^{\times r_\alpha}_*$, a finite-dimensional
    smooth manifold (an open subset of a Euclidean space).
\end{enumerate}


%============================================================
%  ch09_TB_manifold.tex
%  Chapter 9: The Tree-Based Manifold
%
%  Tensor Formats in Banach Spaces
%  Antonio Falcó, CEU Universities
%
%  Based on: FHN2023 (Secs. 4-5)
%  Estimated length: ~40 pp.
%============================================================

\chapter{The Tree-Based Manifold}
\label{chap:TB_manifold}

This chapter establishes the geometric counterpart of
Chapter~\ref{chap:tucker_manifold} for general tree-based tensor
formats. The central result is that $\FTv$ — when it is a
\emph{proper} set of tree-based tensors (Definition~\ref{def:proper_TBT})
— carries a natural $\mathcal{C}^\infty$-Banach manifold structure
that is compatible with the Tucker manifold structures at each
level of the tree, and is an immersed submanifold of the ambient
Banach tensor space $\VDnorm$.

The key idea, due to~\cite{FHN2023}, is that
$\FTv = \bigcap_{k=1}^{\dep{\TD}} \Tucker{\fr_k}{\VD, \PkT{k}{\TD}}$
(Theorem~\ref{thm:intersection_representation}) is an intersection of
Tucker manifolds, and its chart system is obtained by taking the
intersection of the Tucker chart systems at each level. The resulting
charts are natural and compatible with those of each Tucker factor.

Throughout, $(V_\alpha, \|\cdot\|_\alpha)$ are normed spaces,
$\|\cdot\|_D$ satisfies \condINJ\ at each level of $\TD$, and
$\fr \in \ADTD$.

%------------------------------------------------------------
\section{Proper sets of tree-based tensors}
\label{sec:TB_manifold:proper}
%------------------------------------------------------------

Not every $\FTv$ is a genuine intersection of Tucker manifolds — in
degenerate cases it may coincide with a single Tucker manifold.
The following definition isolates the non-degenerate case.

\begin{definition}[Proper tree-based set]
\label{def:proper_TBT}
The set $\FTv$ is called a \emph{proper set of tree-based tensors}
with rank $\fr$ if
\begin{equation*}
    \FTv \neq \Tucker{\fr_k}{\VD, \PkT{k}{\TD}}
    \quad \text{for all } 1 \leq k \leq \dep{\TD}.
\end{equation*}
\end{definition}

\begin{example}[Degenerate case]
\label{ex:degenerate}
For $D = \{1,\ldots,6\}$ and the tree $\TD$ with
$\Sons{D} = \{\{1\}, \{2,3,4,5,6\}\}$ (depth~2):
\begin{align*}
    \mathcal{P}_1(\TD) &= \bigl\{\{1\}, \{2,3,4,5,6\}\bigr\}, \\
    \mathcal{P}_2(\TD) &= \Leaves.
\end{align*}
Since $\dim \Umin{\{1\}}{\bv} = \dim \Umin{\{2,3,4,5,6\}}{\bv}$ for all
$\bv \in \VD$ (Theorem~\ref{thm:min_existence}), the constraint at level~1
is automatically satisfied whenever the leaf-level constraints hold.
Hence $\FTv = \Tucker{\fr_2}{\VD, \Leaves}$ and $\FTv$ is \emph{not}
a proper set. The geometry reduces to the Tucker manifold.
\end{example}

In all examples arising from balanced binary trees (HT format),
Tensor Train, and most practical trees with $\dep{\TD} \geq 2$,
the set $\FTv$ is proper.

%------------------------------------------------------------
\section{The Laplacian-like map and the intersection of model spaces}
\label{sec:TB_manifold:laplacian}
%------------------------------------------------------------

For each level $k$ of the tree, the Tucker manifold
$\Tucker{\fr_k}{\VD, \PkT{k}{\TD}}$ has model space
\begin{equation*}
    \Emodel{k} \coloneqq
    \Bigl(\bigoplus_{\alpha \in \PkT{k}{\TD}}
    \cL(\Umin{\alpha}{\bv}, \Wmin{\alpha}{\bv})
    \otimes \spann\{\mathbf{id}_{[\alpha]}\}\Bigr)
    \times \RR^{\times_{\alpha \in \PkT{k}{\TD}} r_\alpha},
\end{equation*}
where $\mathbf{id}_{[\alpha]} = \bigotimes_{\beta \in \PkT{k}{\TD}
\setminus \{\alpha\}} \Id{\Vbigbeta}$ is the identity on all factors
except $\alpha$.

The \emph{Laplacian-like map} $\Delta_k$ embeds the product of
$\cL(\Umin{\alpha}{\bv}, \Wmin{\alpha}{\bv})$ spaces into $\cL(\VDnorm)$:
\begin{equation}
\label{eq:laplacian}
    \Delta_k : \prod_{\alpha \in \PkT{k}{\TD}}
    \cL(\Umin{\alpha}{\bv}, \Wmin{\alpha}{\bv})
    \longrightarrow \cL(\VDnorm),
    \quad
    (L_\alpha)_\alpha \mapsto \sum_{\alpha \in \PkT{k}{\TD}}
    L_\alpha \otimes \mathbf{id}_{[\alpha]}.
\end{equation}

\begin{proposition}[{\cite[Proposition~4.5]{FHN2023}}]
\label{prop:laplacian_exp}
For $(L_\alpha)_{\alpha \in \PkT{k}{\TD}} \in \prod_\alpha
\cL(\Umin{\alpha}{\bv}, \Wmin{\alpha}{\bv})$:
\begin{equation}
\label{eq:exp_laplacian}
    \exp\bigl(\Delta_k((L_\alpha)_\alpha)\bigr)
    = \bigotimes_{\alpha \in \PkT{k}{\TD}} \exp(L_\alpha).
\end{equation}
\end{proposition}

\begin{proof}
Put $L = \sum_\alpha L_\alpha \otimes \mathbf{id}_{[\alpha]}$.
For $\alpha \neq \beta$:
$(L_\alpha \otimes \mathbf{id}_{[\alpha]}) \circ
(L_\beta \otimes \mathbf{id}_{[\beta]}) =
L_\alpha \otimes L_\beta \otimes \mathbf{id}_{[{\alpha,\beta}]}$,
and by the nilpotency of $L_\alpha$ (Proposition~\ref{prop:nilpotent_algebra}),
$\exp(L_\alpha \otimes \mathbf{id}_{[\alpha]}) = \exp(L_\alpha) \otimes
\mathbf{id}_{[\alpha]}$. Since the summands commute,
$\exp(L) = \bigodot_\alpha \exp(L_\alpha \otimes \mathbf{id}_{[\alpha]})
= \bigotimes_\alpha \exp(L_\alpha)$.
\end{proof}

The critical consequence is that the local representation
$\bw = \exp(L)(\bu)$ for $L \in \Emodel{k}$ is
\emph{independent of $k$}: for any $k$ and any $L \in \Emodel{k}$
with the same underlying tensor $L = \Delta_k((L_\alpha)_\alpha)$,
the formula $\bw = (\bigotimes_\alpha \exp(L_\alpha))(\bu)$ gives
the same $\bw$.

%------------------------------------------------------------
\section{Charts for the tree-based manifold}
\label{sec:TB_manifold:charts}
%------------------------------------------------------------

\subsection*{Intersection of chart domains}

For $\bv \in \FTv$, define the \emph{tree-based chart domain}:
\begin{equation}
\label{eq:TBT_chart_domain}
    \mathcal{W}(\bv) \coloneqq
    \bigcap_{k=1}^{\dep{\TD}} \chart^{(k)}(\bv),
\end{equation}
where $\chart^{(k)}(\bv)$ is the chart domain of
$\Tucker{\fr_k}{\VD, \PkT{k}{\TD}}$ at $\bv$
(Section~\ref{sec:tucker_manifold:charts}).

\subsection*{Model space}

Define the \emph{tree-based model space}:
\begin{equation}
\label{eq:TBT_model}
    \mathbf{E}(\bv) \coloneqq
    \bigcap_{k=1}^{\dep{\TD}} \Emodel{k}
    \subset \cL(\VDnorm),
\end{equation}
and the \emph{tree-based core set}:
\begin{equation*}
    \mathcal{O}(\bv) \coloneqq
    \bigcap_{k=1}^{\dep{\TD}}
    \Tucker{\fr_k}\Bigl({}_a\bigotimes_{\alpha \in \PkT{k}{\TD}}
    \Umin{\alpha}{\bv}\Bigr).
\end{equation*}

\begin{lemma}[{\cite[Lemma~4.6]{FHN2023}}]
\label{lem:TBT_chart_open}
Under condition \condINJ\ at each level, $\mathbf{E}(\bv) \times
\mathcal{O}(\bv)$ is an open subset of the Banach space
$\mathbf{E}(\bv) \times {}_a\bigotimes_{\alpha \in \PkT{1}{\TD}}
\Umin{\alpha}{\bv}$.
\end{lemma}

\begin{proof}[Sketch]
Each $\mathcal{O}_k \coloneqq \Tucker{\fr_k}({}_a\bigotimes_\alpha
\Umin{\alpha}{\bv})$ is open in its finite-dimensional ambient space
(it is defined by rank conditions, cf.\ Remark~\ref{rem:Rstar_open}).
The intersection $\mathcal{O}(\bv) = \bigcap_k \mathcal{O}_k$ is
taken in ${}_a\bigotimes_{\alpha \in \PkT{1}{\TD}} \Umin{\alpha}{\bv}$
(since all factors have the same underlying space by the tree structure),
and is an open subset. Each $\Emodel{k}$ is a closed subspace of
$\cL(\VDnorm)$ (Lemma~\ref{lem:two_faces}), so $\mathbf{E}(\bv)$ is
a Banach space and $\mathbf{E}(\bv) \times \mathcal{O}(\bv)$ is open.
\end{proof}

\subsection*{Tree-based chart map}

Since $L \in \mathbf{E}(\bv)$ satisfies $L \in \Emodel{k}$ for all
$k$, the formula $\bw = \exp(L)(\bu)$ is well-defined and independent
of the level index $k$ by Proposition~\ref{prop:laplacian_exp}.
This allows us to define:

\begin{definition}[Tree-based chart]
\label{def:TBT_chart}
For $\bv \in \FTv$, the \emph{tree-based chart map}
\begin{equation}
\label{eq:TBT_chart_map}
    \boldsymbol{\xi}_\bv : \mathcal{W}(\bv)
    \longrightarrow \mathbf{E}(\bv) \times \mathcal{O}(\bv)
\end{equation}
is defined by $\boldsymbol{\xi}_\bv(\bw) = (L, \bu)$ where
$\bw = \exp(L)(\bu)$, with $L \in \mathbf{E}(\bv)$ and
$\bu \in \mathcal{O}(\bv)$ uniquely determined by any Tucker chart
$\chartmap{\bv}^{(k)}$.
\end{definition}

%------------------------------------------------------------
\section{The $\mathcal{C}^\infty$-Banach manifold structure}
\label{sec:TB_manifold:cinfty}
%------------------------------------------------------------

\begin{theorem}[$\mathcal{C}^\infty$ tree-based manifold,
{\cite[Theorem~4.7]{FHN2023}}]
\label{thm:TBT_Banach_Manifold}
Let $\TD$ be a dimension partition tree over $D$ with $\dep{\TD} \geq 2$,
and let $\fr \in \ADTD$ be such that $\FTv$ is a proper set of
tree-based tensors. Assume \condINJ\ holds at each level of $\TD$.
Then:
\begin{enumerate}
\item[\textup{(i)}] The collection
    $\{(\mathcal{W}(\bv), \boldsymbol{\xi}_\bv)\}_{\bv \in \FTv}$
    is a $\mathcal{C}^\infty$-atlas on $\FTv$, making it a
    $\mathcal{C}^\infty$-Banach manifold modelled on
    $\mathbf{E}(\bv) \times {}_a\bigotimes_{\alpha \in \PkT{1}{\TD}}
    \Umin{\alpha}{\bv}$.

\item[\textup{(ii)}] The atlas is compatible with the Tucker manifold
    atlas at each level: the inclusion
    \begin{equation*}
        \mathfrak{i}_k : \FTv \hookrightarrow
        \Tucker{\fr_k}{\VD, \PkT{k}{\TD}}
    \end{equation*}
    is a $\mathcal{C}^\infty$-immersion for each $1 \leq k \leq \dep{\TD}$.
\end{enumerate}
\end{theorem}

\begin{proof}[Sketch]
The atlas conditions follow from Lemma~\ref{lem:TBT_chart_open}
and the fact that the transition maps of the tree-based atlas are
restrictions of the transition maps of the Tucker atlases (which are
$\mathcal{C}^\infty$ by Lemma~\ref{lem:premanifold}). Compatibility
with the Tucker atlases is a consequence of the identity
$\boldsymbol{\xi}_\bv = \chartmap{\bv}^{(k)}|_{\mathcal{W}(\bv)}$
for each $k$. See~\cite{FHN2023} for the complete proof.
\end{proof}

%------------------------------------------------------------
\section{Tangent space of the tree-based manifold}
\label{sec:TB_manifold:tangent}
%------------------------------------------------------------

\begin{proposition}[Tangent space]
\label{prop:TBT_tangent}
Under the assumptions of Theorem~\ref{thm:TBT_Banach_Manifold},
for $\bv \in \FTv$:
\begin{equation}
\label{eq:TBT_tangent}
    T_\bv \mathfrak{i}\bigl(\mathbb{T}_\bv(\FTv)\bigr)
    = \bigcap_{k=1}^{\dep{\TD}}
    T_\bv \mathfrak{i}_k
    \bigl(\mathbb{T}_\bv(\Tucker{\fr_k}{\VD, \PkT{k}{\TD}})\bigr)
    = \bigcap_{k=1}^{\dep{\TD}} \ZD[\PkT{k}{\TD}]{\bv},
\end{equation}
where $\ZD[\mathcal{P}]{\bv}$ is the tangent subspace
\eqref{eq:ZD} for the partition $\mathcal{P}$.
\end{proposition}

\begin{proof}
The model space of $\FTv$ is $\mathbf{E}(\bv) \times \mathcal{O}_1$,
where $\mathbf{E}(\bv) = \bigcap_k \Emodel{k}$. Since $T_\bv \mathfrak{i}_k$
is injective (Proposition~\ref{prop:tangent_space}), the tangent map
$T_\bv \mathfrak{i}$ is also injective and its image is
$\bigcap_k T_\bv \mathfrak{i}_k(\mathbb{T}_\bv \Tucker{\fr_k}{\VD, \PkT{k}})$.
The last equality follows from the explicit formula
\eqref{eq:tangent_formula} for each level.
\end{proof}

%------------------------------------------------------------
\section{Immersion in the ambient Banach space}
\label{sec:TB_manifold:immersion}
%------------------------------------------------------------

\begin{theorem}[Immersion of the tree-based manifold,
{\cite[Corollary~4.9]{FHN2023}}]
\label{thm:TBT_immersion}
Under the assumptions of Theorem~\ref{thm:TBT_Banach_Manifold}:
\begin{equation}
\label{eq:TBT_tangent_split}
    T_\bv \mathfrak{i}\bigl(\mathbb{T}_\bv(\FTv)\bigr)
    \in \GrassB{\VDnorm}
    \quad \text{for all } \bv \in \FTv.
\end{equation}
Consequently, $\FTv$ is an immersed submanifold of $\VDnorm$.
\end{theorem}

\begin{proof}
From Theorem~\ref{thm:tucker_immersion} applied at each level,
$\ZD[\PkT{k}]{\bv} \in \GrassB{\VDnorm}$ for $1 \leq k \leq \dep{\TD}$.
By Proposition~\ref{prop:TBT_tangent}, the tangent space of $\FTv$
is $\bigcap_k \ZD[\PkT{k}]{\bv}$. Applying
Lemma~\ref{lem:direct_sum_G} iteratively (intersection of split
subspaces is split), we obtain
$\bigcap_k \ZD[\PkT{k}]{\bv} \in \GrassB{\VDnorm}$.
\end{proof}

\begin{remark}[Embedded vs.\ immersed]
\label{rem:embedded_vs_immersed}
Theorem~\ref{thm:TBT_immersion} shows that $\FTv$ is an
\emph{immersed} submanifold of $\VDnorm$. Moreover,
Theorem~\ref{thm:TBT_Banach_Manifold}(ii) shows it is an
\emph{embedded} submanifold of each Tucker manifold
$\Tucker{\fr_k}{\VD, \PkT{k}{\TD}}$. Whether $\FTv$ is embedded
in $\VDnorm$ (i.e., the manifold topology coincides with the
subspace topology) is a more delicate question related to the
global topology discussed in Section~\ref{sec:TB_manifold:topology}.
\end{remark}

%------------------------------------------------------------
\section{Comparison with alternative descriptions}
\label{sec:TB_manifold:comparison}
%------------------------------------------------------------

The geometric description of $\FTv$ in~\cite{FHN2023} differs from
earlier approaches in several respects.

\paragraph{Uschmajew--Vandereycken \cite{UschmajewVandereycken2013}.}
For the HT format in finite-dimensional spaces, they construct
$\FTv$ as a quotient manifold
$P \backslash \bigl(\GL{r_\alpha} \times \cdots\bigr)$
for a gauge group $P$. This gives a finite-dimensional manifold
and is well-suited for numerical algorithms (Riemannian gradient
methods). However, it does not extend to infinite-dimensional
spaces and does not make the compatibility with Tucker manifolds
at each level explicit.

\paragraph{Holtz--Rohwedder--Schneider \cite{HoltzRohwedderSchneider2012}.}
For the TT format in finite dimensions, a similar quotient
construction is used. The connection to the Tucker manifold at
each level is not made explicit.

\paragraph{Falcó--Hackbusch--Nouy (2015 preprint \cite{FHN2015}).}
An earlier version of the atlas in~\cite{FHN2023} was proposed
using a different chart system. The description in~\cite{FHN2023}
is more natural because the chart maps are directly inherited from
the Tucker manifold charts via the Laplacian map (Proposition~\ref{prop:laplacian_exp}).

\paragraph{Key novelties of \cite{FHN2023}.}
\begin{enumerate}
\item The atlas $\{(\mathcal{W}(\bv), \boldsymbol{\xi}_\bv)\}$
    is \emph{natural}: it is the restriction of the Tucker atlases
    to the intersection $\FTv$, without any additional construction.
\item The \emph{compatibility} with Tucker manifold structures at
    each level is built into the definition, making the geometric
    relationships transparent.
\item The description holds in \emph{infinite-dimensional Banach
    spaces} under the mild condition \condINJ.
\item The tangent space formula \eqref{eq:TBT_tangent} is explicit
    and computable.
\end{enumerate}

%------------------------------------------------------------
\section{Topology and connected components}
\label{sec:TB_manifold:topology}
%------------------------------------------------------------

\begin{proposition}
\label{prop:TBT_connected}
For the Tucker manifold, $\Tucker{\fr}{\VD}$ is connected for any
$\fr \in \AD{\VD}$ with $\fr \neq \mathbf{0}$. This follows because
$\prod_\alpha \Grass{r_\alpha}{V_\alpha}$ is connected (each
$\Grass{r}{V}$ is connected as a Banach manifold) and the fibre
$\RR^{\times r_\alpha}_*$ is connected.
\end{proposition}

For general tree-based formats, the connectedness of $\FTv$ is an
open problem in general. The following is known:

\begin{proposition}
\label{prop:TBT_connected_cases}
\begin{enumerate}
\item[\textup{(i)}] For the TT format and binary HT formats in
    \emph{finite dimensions}, $\FTv$ is connected
    (see~\cite{UschmajewVandereycken2013, HoltzRohwedderSchneider2012}).
\item[\textup{(ii)}] For general $\TD$ and infinite-dimensional
    spaces, connectedness of $\FTv$ is an open question.
\end{enumerate}
\end{proposition}

\begin{remark}[Stratification]
\label{rem:stratification}
The decomposition
$\VD = \bigsqcup_{\fr \in \ADTD} \FT{\fr}{\VD}{\TD}$
(Theorem~\ref{thm:TBT_decomposition}) is a stratification of $\VD$
by tree-based rank. Each stratum $\FTv$ is an immersed manifold,
and $\FTleq{\fr}{\VD}{\TD} = \overline{\FTv}^{\text{weak}}$
is the weak closure. The boundary strata
$\partial \FTleq{\fr}{\VD}{\TD} = \FTleq{\fr}{\VD}{\TD} \setminus \FTv$
consist of tensors with strictly smaller rank.
This stratification is the tensor-format analogue of the rank
stratification of matrices.
\end{remark}

%------------------------------------------------------------
\section{Historical notes}
\label{sec:TB_manifold:history}
%------------------------------------------------------------

The geometric structure of $\FTv$ as a Banach manifold was established
in~\cite{FHN2023}. The fundamental observation — that the tree-based
manifold is an intersection of Tucker manifolds and that its atlas can
be obtained by intersecting the Tucker atlases — was made possible by
the reformulation of the Tucker manifold charts using the Laplacian-like
map $\Delta_k$ (Proposition~\ref{prop:laplacian_exp}).

The compatibility condition between the tree-based atlas and the Tucker
atlases at each level (Theorem~\ref{thm:TBT_Banach_Manifold}(ii)) is
the key property that distinguishes the description of~\cite{FHN2023}
from earlier work: it allows one to apply the Tucker manifold geometry
(tangent spaces, projections, Dirac--Frenkel principle) at each level
of the tree independently, and then combine the results.

The Dirac--Frenkel variational principle for tree-based formats
(Chapter~\ref{chap:DF}) exploits this compatibility in an essential
way: the reduced ODE on $\FTv$ is obtained by projecting the vector
field onto the tangent space $\bigcap_k \ZD[\PkT{k}]{\bv}$, and the
existence of this projection (Corollary~\ref{cor:best_proj_tangent}
applied at each level) follows from the fact that each $\ZD[\PkT{k}]{\bv}$
is a split subspace of $\VDnorm$.



%%  PART IV: VARIATIONAL PRINCIPLES AND APPLICATIONS
\part{Variational Principles and Applications}
\label{part:variational}

%============================================================
%  ch10_DF.tex
%  Chapter 10: The Dirac--Frenkel Variational Principle
%
%  Tensor Formats in Banach Spaces
%  Antonio Falcó, CEU Universities
%
%  Based on: FHN2019 (Sec. 5), FHN2023 (Sec. 5)
%  Estimated length: ~35 pp.
%============================================================

\chapter{The Dirac--Frenkel Variational Principle}
\label{chap:DF}

The geometric framework developed in Chapters~\ref{chap:tucker_manifold}
and~\ref{chap:TB_manifold} has a direct dynamical application: it
provides the rigorous foundation for \emph{model reduction} of
differential equations on tensor Banach spaces. Given an abstract
Cauchy problem on $\VDnorm$, one looks for an approximate solution
that evolves on the manifold $\Tucker{\fr}{\VD}$ (or $\FTv$).
The natural way to constrain the dynamics to the manifold is the
\emph{Dirac--Frenkel variational principle}: at each time $t$, the
velocity $\dot{\bv}_r(t)$ is chosen as the best approximation of the
vector field $\bF(t, \bv_r(t))$ in the tangent space $\ZD{\bv_r(t)}$.

This principle was originally proposed by Dirac and Frenkel in 1930 for
the time-dependent Schrödinger equation in Hilbert spaces. The present
chapter extends it to general Banach tensor spaces, establishing the
existence of the reduced ODE and characterising its solutions.
Sections~\ref{sec:DF:cauchy}--\ref{sec:DF:hilbert} set up the
abstract framework. Section~\ref{sec:DF:tucker} proves the Tucker
version (Theorems~5.2 and~5.4 of~\cite{FHN2019}).
Section~\ref{sec:DF:treebased} extends to the tree-based setting
(\cite{FHN2023}).

%------------------------------------------------------------
\section{The abstract Cauchy problem on tensor Banach spaces}
\label{sec:DF:cauchy}
%------------------------------------------------------------

Let $(V_\alpha, \|\cdot\|_\alpha)$ be normed spaces for $\alpha \in D$,
and let $\VDnorm = {}_{\|\cdot\|_D}\bigotimes_{\alpha \in D} V_\alpha$
be a reflexive Banach tensor space satisfying condition~\condINJ.
Consider the abstract Cauchy problem:
\begin{align}
    \dot{\bu}(t) &= \bF(t, \bu(t)), \quad t \geq 0,
    \label{eq:DF:ode} \\
    \bu(0) &= \bu_0,
    \label{eq:DF:ic}
\end{align}
where $\bu_0 \neq \mathbf{0}$, and $\bF: [0, \infty) \times \VDnorm
\to \VDnorm$ satisfies the usual conditions for existence and uniqueness
of solutions (e.g., Lipschitz continuity in the second argument).

The goal is to approximate $\bu(t)$ by a \emph{differentiable curve}
$t \mapsto \bv_r(t)$ lying in $\Tucker{\fr}{\VD}$ for some fixed
$\fr \in \AD{\VD}$ with $\fr \neq \mathbf{0}$, with initial condition
$\bv_r(0) = \bv_0 \in \Tucker{\fr}{\VD}$ chosen as an approximation
of $\bu_0$.

\begin{remark}
\label{rem:DF:initial}
By Theorem~\ref{thm:best_approx_tucker}, a best approximation
$\bv_0 \in \Tucker{\leq\fr}{\VD}$ of $\bu_0$ exists when $\VDnorm$
is reflexive. The initial condition $\bv_0 \in \Tucker{\fr}{\VD}$
(with exact rank $\fr$) can always be arranged by taking
$\bv_0 \in \Tucker{\fr'}{\VD}$ for some $\fr' \leq \fr$ and
adjusting the rank.
\end{remark}

%------------------------------------------------------------
\section{The principle in Hilbert spaces}
\label{sec:DF:hilbert}
%------------------------------------------------------------

When $\VDnorm$ is a Hilbert space with inner product
$\langle \cdot, \cdot \rangle_D$, the constraint to $\Tucker{\fr}{\VD}$
takes the classical form. Since $\ZD{\bv_r(t)} \in \GrassB{\VDnorm}$
(Theorem~\ref{thm:tucker_immersion}), the orthogonal projection
$\Pmetric{\bv_r(t)}$ onto $\ZD{\bv_r(t)}$ is well-defined and the
\emph{Dirac--Frenkel equation} reads:
\begin{equation}
\label{eq:DF:hilbert}
    \dot{\bv}_r(t) = \Pmetric{\bv_r(t)}\bigl(\bF(t, \bv_r(t))\bigr),
\end{equation}
which is equivalent to the variational condition:
\begin{equation}
\label{eq:DF:variational}
    \langle \dot{\bv}_r(t) - \bF(t, \bv_r(t)),\, \dot{\bv} \rangle_D = 0
    \quad \text{for all } \dot{\bv} \in \ZD{\bv_r(t)}.
\end{equation}
The MCTDH (Multi-Configuration Time-Dependent Hartree) method of
quantum dynamics~\cite{Lubich2008} is precisely~\eqref{eq:DF:hilbert}
with $\VDnorm = L^2(\RR^{3d})$ and $\Tucker{\fr}{\VD}$ the Hartree
manifold ($\fr = (1,\ldots,1)$).

%------------------------------------------------------------
\section{Projections in Banach spaces}
\label{sec:DF:projection}
%------------------------------------------------------------

In a general Banach space, there is no canonical inner product and hence
no canonical projection. Two notions replace the orthogonal projection.

\subsection*{Metric projection}

\begin{definition}
\label{def:metric_projection}
For $\bv_r(t) \in \Tucker{\fr}{\VD}$ and $t \in I$, the
\emph{metric projection} $\Pmetric{\bv_r(t)}$ maps
$\dot{\bu} \in \VDnorm$ to the nearest point in $\ZD{\bv_r(t)}$:
\begin{equation*}
    \Pmetric{\bv_r(t)}(\dot{\bu})
    \coloneqq \argmin_{\dot{\bv} \in \ZD{\bv_r(t)}}
    \|\dot{\bu} - \dot{\bv}\|_D.
\end{equation*}
\end{definition}

By Corollary~\ref{cor:best_proj_tangent}, when $\VDnorm$ is reflexive
the metric projection exists. When $\VDnorm$ is strictly convex it is
unique (since $\ZD{\bv_r(t)}$ is convex). The metric projection is,
in general, a \emph{nonlinear} operator on Banach spaces.

\subsection*{Normalised duality mapping}

The duality mapping $J: \VDnorm \to 2^{\dual{\VDnorm}}$ is defined by
\begin{equation*}
    J(\bu) \coloneqq \bigl\{f \in \dual{\VDnorm} :
    \langle \bu, f \rangle = \|\bu\|_D^2 = \|f\|_*^2\bigr\}.
\end{equation*}
It is single-valued when $\VDnorm$ is smooth (Definition~\ref{def:smooth}),
and uniformly continuous on bounded sets when $\VDnorm$ is uniformly smooth.
In a Hilbert space, $J$ coincides with the Riesz isomorphism.

\begin{theorem}[Metric projection characterisation,
{\cite[Theorem~5.2]{FHN2019}}]
\label{thm:metric_projection}
Let $\VDnorm$ be reflexive and strictly convex, and assume \condINJ.
For $\bv_r(t) \in \Tucker{\fr}{\VD}$ and fixed $t \in I$, the
following are equivalent:
\begin{enumerate}
\item[\textup{(a)}] $\dot{\bv}_r(t) = \Pmetric{\bv_r(t)}
    (\bF(t,\bv_r(t)))$.
\item[\textup{(b)}] $\langle \dot{\bv}_r(t) - \dot{\bv},\,
    J(\bF(t,\bv_r(t)) - \dot{\bv}_r(t)) \rangle \geq 0$
    for all $\dot{\bv} \in \ZD{\bv_r(t)}$.
\item[\textup{(c)}] $\langle \dot{\bv},\,
    J(\bF(t,\bv_r(t)) - \dot{\bv}_r(t)) \rangle = 0$
    for all $\dot{\bv} \in \ZD{\bv_r(t)}$.
\end{enumerate}
\end{theorem}

\begin{proof}
This is a direct consequence of the classical characterisation of the
metric projection onto a closed convex set in a reflexive, strictly
convex Banach space, combined with~\cite[Proposition~2.10]{Alber1996}.
Condition~(c) recovers the Hilbert-space variational form: when $J$
is the Riesz map, (c) becomes the orthogonality condition~\eqref{eq:DF:variational}.
\end{proof}

\subsection*{Generalised projection}

When $\VDnorm$ is additionally smooth, an alternative projection is
available. Define $\phi: \VDnorm \times \VDnorm \to \RR$ by
\begin{equation*}
    \phi(\bu, \bv) = \|\bu\|_D^2 - 2\langle \bu, J(\bv) \rangle + \|\bv\|_D^2.
\end{equation*}
The \emph{generalised projection} $\Pi_{\bv_r(t)}$ minimises $\phi$
over $\ZD{\bv_r(t)}$:
\begin{equation*}
    \Pi_{\bv_r(t)}(\dot{\bu})
    \coloneqq \argmin_{\dot{\bv} \in \ZD{\bv_r(t)}}
    \phi(\dot{\bv}, \dot{\bu}).
\end{equation*}
It coincides with the metric projection when $\VDnorm$ is a Hilbert
space.

%------------------------------------------------------------
\section{The Dirac--Frenkel principle for the Tucker manifold}
\label{sec:DF:tucker}
%------------------------------------------------------------

\begin{theorem}[Dirac--Frenkel for Tucker, {\cite[Theorems~5.2,~5.4]{FHN2019}}]
\label{thm:DF_tucker}
Let $\VDnorm$ be reflexive and strictly convex, and assume \condINJ.
Let $\fr \in \AD{\VD}$ with $\fr \neq \mathbf{0}$, and let
$\bv_0 \in \Tucker{\fr}{\VD}$.
\begin{enumerate}
\item[\textup{(Existence)}] There exists a differentiable curve
    $t \mapsto \bv_r(t)$ in $\Tucker{\fr}{\VD}$, defined for $t$ in a
    neighbourhood of $0$, satisfying
    \begin{equation}
    \label{eq:DF_tucker}
        \dot{\bv}_r(t) = \Pmetric{\bv_r(t)}\bigl(\bF(t, \bv_r(t))\bigr),
        \qquad \bv_r(0) = \bv_0.
    \end{equation}

\item[\textup{(Characterisation)}] The curve $t \mapsto \bv_r(t)$
    satisfies~\eqref{eq:DF_tucker} if and only if $\dot{\bv}_r(t)
    \in \ZD{\bv_r(t)}$ and
    \begin{equation}
    \label{eq:DF_tucker_char}
        \|\dot{\bv}_r(t) - \bF(t, \bv_r(t))\|_D
        = \min_{\dot{\bv} \in \ZD{\bv_r(t)}}
        \|\dot{\bv} - \bF(t, \bv_r(t))\|_D.
    \end{equation}
\end{enumerate}
\end{theorem}

\begin{proof}[Sketch]
\textbf{Existence.} The reduced ODE~\eqref{eq:DF_tucker} is an ODE on
the Banach manifold $\Tucker{\fr}{\VD}$. By
Theorem~\ref{thm:tucker_immersion}, $\ZD{\bv_r(t)}$ is a split
subspace of $\VDnorm$ for each $\bv_r(t) \in \Tucker{\fr}{\VD}$.
The right-hand side of~\eqref{eq:DF_tucker} depends smoothly on
$\bv_r(t)$ (the metric projection onto a smoothly-varying split
subspace is smooth), so local existence and uniqueness follow from the
ODE existence theorem on Banach manifolds~\cite[Theorem~4.1.5]{Lang1999}.

\textbf{Characterisation.} Since $\ZD{\bv_r(t)}$ is closed and
convex, Corollary~\ref{cor:best_proj_tangent} gives existence of the
minimum in~\eqref{eq:DF_tucker_char}. The equivalence with the metric
projection condition~\eqref{eq:DF_tucker} follows from the uniqueness
(strict convexity) and Theorem~\ref{thm:metric_projection}.
\end{proof}

\begin{remark}[Kinematic equations]
\label{rem:kinematic}
In local chart coordinates $(L_\alpha, C^{(D)})$ at $\bv_r(t)$,
the reduced ODE~\eqref{eq:DF_tucker} takes the form of a system
of coupled equations for the subspace parameters $L_\alpha(t)$ and
the core tensor $C^{(D)}(t)$. The tangent formula~\eqref{eq:tangent_formula}
makes this explicit:
\begin{equation*}
    \dot{\bv}_r(t)
    = \sum_{(i_\alpha)} \dot{C}^{(D)}_{(i_\alpha)} \bigotimes_\alpha u^{(\alpha)}_{i_\alpha}
    + \sum_\alpha \sum_{i_\alpha} \dot{u}^{(\alpha)}_{i_\alpha}
      \otimes \bU^{(\alpha)}_{i_\alpha}.
\end{equation*}
The Dirac--Frenkel condition~\eqref{eq:DF_tucker_char} then yields
coupled evolution equations for the basis vectors
$u^{(\alpha)}_{i_\alpha}(t) \in V_\alpha$ and the core tensor
$C^{(D)}(t)$, which generalise the MCTDH equations to arbitrary
Tucker ranks in Banach spaces.
\end{remark}

\subsection*{The Hartree manifold}

\begin{example}[Time-dependent Hartree method]
\label{ex:hartree}
For $\fr = \mathbf{1} = (1, \ldots, 1)$, the Tucker manifold
$\Tucker{\mathbf{1}}{\VD}$ consists of elementary tensors
$\bv = \lambda \bigotimes_\alpha v_\alpha$ ($v_\alpha \in V_\alpha$,
$\lambda \in \RR$), called the \emph{Hartree manifold}~\cite{Lubich2008}.
The Dirac--Frenkel ODE~\eqref{eq:DF_tucker} on the Hartree manifold
gives the \emph{time-dependent Hartree equations}. In a Hilbert space
with $\|v_\alpha\|_\alpha = 1$ (unit sphere constraint):
\begin{align*}
    \dot{\lambda}(t) &= \langle \bigotimes_\alpha v_\alpha(t),
        \bF(t, \bv_r(t)) \rangle_D, \\
    \lambda(t)\,\dot{v}_\alpha(t) &= P_{W_\alpha(t)}\bigl[
        (\Id{\alpha} \otimes \langle \cdot, \bigotimes_{\beta \neq \alpha}
        v_\beta(t) \rangle)\bF(t, \bv_r(t))\bigr],
\end{align*}
where $P_{W_\alpha(t)}$ is the orthogonal projection onto the
complement of $\spann\{v_\alpha(t)\}$ in $V_\alpha$. This is
the system solved numerically by MCTDH.
\end{example}

%------------------------------------------------------------
\section{Extension to tree-based formats}
\label{sec:DF:treebased}
%------------------------------------------------------------

The extension of Theorem~\ref{thm:DF_tucker} to $\FTv$ is immediate
given the immersion result of Chapter~\ref{chap:TB_manifold}.

\begin{theorem}[Dirac--Frenkel for tree-based formats,
{\cite[Section~5]{FHN2023}}]
\label{thm:DF_treebased}
Let $\TD$ be a dimension partition tree, $\fr \in \ADTD$ such that
$\FTv$ is a proper set of tree-based tensors. Assume $\VDnorm$ is
reflexive and strictly convex, and that \condINJ\ holds at each
level of $\TD$. Let $\bv_0 \in \FTv$.
\begin{enumerate}
\item[\textup{(Existence)}] There exists a differentiable curve
    $t \mapsto \bv_r(t)$ in $\FTv$ satisfying
    \begin{equation}
    \label{eq:DF_TBT}
        \dot{\bv}_r(t) = \Pmetric{\bv_r(t)}\bigl(\bF(t, \bv_r(t))\bigr),
        \qquad \bv_r(0) = \bv_0,
    \end{equation}
    where $\Pmetric{\bv_r(t)}$ is the metric projection onto
    $T_{\bv_r(t)}\mathfrak{i}(\mathbb{T}_{\bv_r(t)}(\FTv))
    = \bigcap_k \ZD[\PkT{k}{\TD}]{\bv_r(t)}$.
\item[\textup{(Characterisation)}] The curve satisfies~\eqref{eq:DF_TBT}
    if and only if $\dot{\bv}_r(t) \in \bigcap_k \ZD[\PkT{k}{\TD}]{\bv_r(t)}$
    and
    \begin{equation}
    \label{eq:DF_TBT_char}
        \|\dot{\bv}_r(t) - \bF(t, \bv_r(t))\|_D
        = \min_{\dot{\bv} \in \bigcap_k \ZD[\PkT{k}]{\bv_r(t)}}
        \|\dot{\bv} - \bF(t, \bv_r(t))\|_D.
    \end{equation}
\end{enumerate}
\end{theorem}

\begin{proof}
The proof is identical to that of Theorem~\ref{thm:DF_tucker}, with
$\Tucker{\fr}{\VD}$ replaced by $\FTv$ and $\ZD{\bv_r(t)}$ replaced
by $\bigcap_k \ZD[\PkT{k}]{\bv_r(t)}$. The key inputs are:
(i) $\FTv$ is a $\mathcal{C}^\infty$-Banach manifold
(Theorem~\ref{thm:TBT_Banach_Manifold}); (ii) its tangent space is
a split subspace of $\VDnorm$ (Theorem~\ref{thm:TBT_immersion}); and
(iii) the metric projection onto a split subspace in a reflexive
strictly convex space is well-defined and unique.
\end{proof}

\begin{remark}[Algorithmic interpretation]
\label{rem:DF_algorithms}
Theorem~\ref{thm:DF_treebased} provides the mathematical foundation
for a class of algorithms that integrate tensor-format ODEs by evolving
the transfer tensors $\{C^{(\alpha)}\}_{\alpha \in \TD}$ and the leaf
basis vectors $\{u^{(j)}_{i_j}\}_{j \in \Leaves}$ along the
Dirac--Frenkel trajectories. For the TT format, this is the starting
point of the \emph{projector-splitting integrator} of Lubich, Oseledets,
and Vandereycken~\cite{LubichOseledetsVandereycken2015}. For general
tree-based formats, the corresponding algorithm was introduced by
Ceruti, Lubich, and Walach~\cite{CerutiLubichWalach2020}.
\end{remark}

%------------------------------------------------------------
\section{Historical notes}
\label{sec:DF:history}
%------------------------------------------------------------

The Dirac--Frenkel variational principle was proposed independently by
Dirac~\cite{Dirac1930} and Frenkel~\cite{Frenkel1934} in 1930 and 1934
for the time-dependent Schrödinger equation in Hilbert space. The
formulation used here — projection of the vector field onto the tangent
space of a manifold — is that of McLachlan et al.\ and Lubich.

Koch and Lubich~\cite{KochLubich2007} proved existence and uniqueness
of solutions to the Dirac--Frenkel ODE on Tucker manifolds in
finite-dimensional Hilbert spaces, and used it to analyse the MCTDH
method. The extension to infinite-dimensional Banach spaces
(Theorem~\ref{thm:DF_tucker}) was established in~\cite{FHN2019},
where the key difficulty — showing that the tangent space $\ZD{\bv_r(t)}$
is a split subspace of $\VDnorm$, without the Hilbert space structure —
was resolved by the immersion theorem (Theorem~\ref{thm:tucker_immersion}).

The extension to tree-based formats (Theorem~\ref{thm:DF_treebased})
is the main application result of~\cite{FHN2023}: having established
the geometry of $\FTv$ in Chapter~\ref{chap:TB_manifold}, the
Dirac--Frenkel principle follows with no additional effort.


%============================================================
%  ch12_quantum_states.tex
%  Chapter 12: Quantum States as Tensor Formats
%
%  Tensor Formats in Banach Spaces
%  Antonio Falcó, CEU Universities
%
%  Based on: Acevedo--Falcó (2026), FH2012, FHN2019
%  Estimated length: ~25 pp.
%============================================================

\chapter{Quantum States as Tensor Formats}
\label{chap:quantum_states}

The theory developed in Parts~I--IV connects naturally to quantum
mechanics through a single observation: for a non-zero vector $v$ in a
Hilbert tensor space $H = \bigotimes_{j=1}^d H_j$, the minimal subspace
$\Umin{\alpha}{v}$ coincides with the support of the reduced density
operator of the quantum state associated to $v$. This chapter makes
this connection precise and develops the geometric theory of quantum
states from the perspective of tensor formats and minimal subspaces.

\begin{remark}[Notation convention for Part~\ref{part:quantum}]
\label{rem:qs:notation}
In Parts~\ref{part:foundations}--\ref{part:variational}, tensors in
$\VDnorm$ are written in boldface ($\bv, \bu, \bw$) to distinguish
them from their component vectors ($v_\alpha \in V_\alpha$). In this
part, the primary objects are density operators $\rho \in \cTsa{1}(H)$
and Hilbert-space vectors $v \in H$; we follow the standard convention
of the operator-theory literature and use lightface notation throughout.
When a pure state $\rho_v = v \otimes J(v)/\|v\|^2$ is considered as
a tensor in $H \cong \bigotimes_j H_j$ and the tensor structure is the
explicit focus, the boldface notation $\bv$ is used.
\end{remark}

The state space of a quantum system on $H$ is the set $\Dstate{H}$
of positive trace-class operators of unit trace (see
Appendix~\ref{app:quantum} for the relevant operator-theoretic background).
Each fixed-rank stratum $\Drank{r}{H}$ carries the structure of a smooth
Banach manifold, fibred over the Grassmannian $\Grass{r}{H}$ — precisely
the structure of the Tucker manifold of Chapter~\ref{chap:tucker_manifold}
specialised to the trace-class Banach space $\cTsa{1}(H)$.

Section~\ref{sec:qs:pure} establishes the fundamental connection between
minimal subspaces and partial traces for pure states.
Section~\ref{sec:qs:mixed} extends to mixed states and identifies
the rank stratification of $\Dstate{H}$ with the Tucker rank stratification.
Section~\ref{sec:qs:entanglement} connects the Tucker rank vector with
the entanglement structure of multipartite systems.
Section~\ref{sec:qs:manifold} develops the Banach manifold structure of
$\Drank{r}{H}$ as a special case of the Tucker manifold theory.
Section~\ref{sec:qs:approx} applies the best-approximation theory of
Chapter~\ref{chap:minsubspaces} to quantum state compression.
Section~\ref{sec:qs:history} contains historical notes.

Throughout this chapter, $H = \bigotimes_{j=1}^d H_j$ is a Hilbert
tensor product of separable complex Hilbert spaces $H_j$,
and $\cTsa{1}(H)$ is the real Banach space of self-adjoint trace-class
operators on $H$, equipped with the trace norm $\|\cdot\|_1$
(see Appendix~\ref{app:quantum}).

%------------------------------------------------------------
\section{Pure states and minimal subspaces}
\label{sec:qs:pure}
%------------------------------------------------------------

\subsection*{The pure state associated to a vector}

For a non-zero vector $v \in H$, the \emph{pure state associated to $v$} is
\begin{equation}
\label{eq:qs:pure_state}
    \rho_v \coloneqq \frac{v \otimes J(v)}{\|v\|^2}
    \in \Dstate{H},
\end{equation}
where $J: H \to H^*$ is the Riesz isomorphism and $v \otimes J(v)$
is the rank-one operator $w \mapsto \langle v, w\rangle v / \|v\|^2$.
Two vectors $v$ and $\lambda v$ ($\lambda \in \CC^*$) give the same
pure state. The map $v \mapsto \rho_v$ thus induces a bijection between
the projective space $\mathbb{P}(H)$ and the set $\Drank{1}{H}$ of
pure states.

\subsection*{The fundamental connection}

The following theorem is the cornerstone of this chapter: it identifies
the minimal subspaces of $v$ (in the sense of
Definition~\ref{def:minimal_subspace}) with the supports of the
partial traces of $\rho_v$.

\begin{theorem}[Minimal subspaces and partial traces]
\label{thm:qs:minimal_partial_trace}
Let $v \in H = \bigotimes_{j=1}^d H_j$ with $\|v\|=1$ and let
$\rho_v = v \otimes J(v)$ be the associated pure state. For each
$\alpha \in D = \{1, \ldots, d\}$, let
$\rho_\alpha \coloneqq \TrS{\alpha^c}(\rho_v) \in \cTsa{1}(H_\alpha)$
be the partial trace over all factors except $\alpha$. Then:
\begin{enumerate}
\item[\textup{(a)}] $\rho_\alpha$ is a positive trace-class operator
    with $\Tr(\rho_\alpha) = 1$ and
    \begin{equation}
    \label{eq:qs:supp_min}
        \supp(\rho_\alpha) = \Umin{\alpha}{v}.
    \end{equation}
\item[\textup{(b)}] $\mathrm{rank}(\rho_\alpha) = r_\alpha(v)$,
    the Tucker $\alpha$-rank of $v$.
\item[\textup{(c)}] The non-zero eigenvalues of $\rho_\alpha$ are
    $\lambda_k^{(\alpha)}(v) > 0$, $k = 1, \ldots, r_\alpha(v)$,
    with $\sum_k \lambda_k^{(\alpha)}(v) = 1$.
    For $d=2$, $\lambda_k^{(1)}(v) = \lambda_k^{(2)}(v) = s_k(v)^2$
    where $s_k(v)$ are the Schmidt coefficients of $v$.
\end{enumerate}
\end{theorem}

\begin{proof}
By the definition of $\Umin{\alpha}{v}$ (Definition~\ref{def:minimal_subspace}),
$\Umin{\alpha}{v}$ is the smallest subspace $U \subset H_\alpha$ such
that $v \in U \otimes_H H_{\alpha^c}$. This is exactly the range of the
map $T: H_\alpha^* \to H_{\alpha^c}$ defined by $T(\varphi) = (\varphi
\otimes \Id{\alpha^c})(v)$, hence $\Umin{\alpha}{v} =
\overline{\mathrm{Im}(T^*)}$ where $T^*: H_{\alpha^c}^* \to H_\alpha$
is the adjoint.

On the other hand, $\rho_\alpha = \TrS{\alpha^c}(\rho_v)$ is the
operator on $H_\alpha$ characterised by
$\langle u, \rho_\alpha w\rangle_{H_\alpha} = \langle u \otimes
\Id{\alpha^c}, \rho_v (w \otimes \Id{\alpha^c})\rangle$ for all
$u, w \in H_\alpha$. A direct computation using
$\rho_v = v \otimes J(v)$ gives
\begin{equation*}
    \rho_\alpha = \TrS{\alpha^c}(v \otimes J(v))
    = T T^*: H_\alpha \to H_\alpha
\end{equation*}
(the Gram operator of $T$ composed with its adjoint). Hence
$\mathrm{Im}(\rho_\alpha) = \mathrm{Im}(TT^*) = \overline{\mathrm{Im}(T^*)}
= \Umin{\alpha}{v}$, proving~\eqref{eq:qs:supp_min}.
Part~(b) follows since $\mathrm{rank}(TT^*) = \mathrm{rank}(T) = r_\alpha(v)$.
Part~(c) follows from the spectral theorem for compact operators
(Theorem~\ref{thm:app:spectral}): the eigenvalues of $\rho_\alpha = TT^*$
are the squared singular values of $T$, and $\Tr(\rho_\alpha) = 1$
gives $\sum_k \lambda_k^{(\alpha)} = 1$.
\end{proof}

\begin{remark}[The $d=2$ case: Schmidt decomposition]
\label{rem:qs:schmidt}
For $d=2$, Theorem~\ref{thm:qs:minimal_partial_trace} reduces to the
classical Schmidt decomposition. Every $v \in H_1 \otimes_H H_2$ can
be written uniquely (up to phases) as
\begin{equation*}
    v = \sum_{k=1}^r s_k\, e_k \otimes f_k,
    \quad s_1 \geq s_2 \geq \cdots \geq s_r > 0,
\end{equation*}
with $\{e_k\}_{k=1}^r \subset H_1$ and $\{f_k\}_{k=1}^r \subset H_2$
orthonormal, $r = r_1(v) = r_2(v)$. Then
$\Umin{1}{v} = \spann\{e_k\}$ and $\Umin{2}{v} = \spann\{f_k\}$,
and $\rho_1 = \sum_k s_k^2\, e_k \otimes J(e_k)$,
$\rho_2 = \sum_k s_k^2\, f_k \otimes J(f_k)$.
\end{remark}

\begin{remark}[Tucker format as the natural language for multipartite states]
\label{rem:qs:tucker_natural}
Theorem~\ref{thm:qs:minimal_partial_trace} shows that for a pure state
$\rho_v$, the Tucker rank vector $(r_\alpha(v))_{\alpha \in D}$ is
an intrinsic property of the state: it is determined by the ranks of
the reduced density operators $\{\rho_\alpha\}$, which are physical
observables. The Tucker manifold $\Tucker{\fr}{H}$ is therefore the
natural variational class for quantum systems with prescribed
entanglement structure $(r_\alpha)_{\alpha \in D}$.
\end{remark}

\begin{figure}[ht]
\centering
\begin{tikzpicture}[>=Latex, node distance=1.8cm,
  state/.style={draw, rounded corners=4pt, minimum width=2.4cm,
                minimum height=0.75cm, align=center, font=\small},
  obs/.style={draw, dashed, rounded corners=4pt, minimum width=2.4cm,
              minimum height=0.75cm, align=center, font=\small, fill=orange!8},
  arr/.style={->, thick},
  label/.style={font=\scriptsize, draw=none}
]
% Central: pure state
\node[state, fill=blue!8] (v)
  {$v \in \Tucker{\fr}{H}$\\[1pt]\scriptsize pure state};

% Tucker rank annotation
\node[label, above=0.4cm of v]
  {Tucker rank $\fr = (r_\alpha)_{\alpha \in D}$};

% Partial traces
\node[obs, left=3.0cm of v] (r1)
  {$\rho_1 = \TrS{1^c}(\rho_v)$\\[1pt]\scriptsize red.\ density op.\ on $H_1$};
\node[obs, right=3.0cm of v] (r2)
  {$\rho_2 = \TrS{2^c}(\rho_v)$\\[1pt]\scriptsize red.\ density op.\ on $H_2$};
\node[obs, below=1.6cm of v] (ralpha)
  {$\rho_\alpha = \TrS{\alpha^c}(\rho_v)$\\[1pt]\scriptsize
   for each $\alpha \in D$};

% Minimal subspaces
\node[state, fill=green!8, left=3.0cm of r1, yshift=0cm] (U1)
  {$\Umin{1}{v} = \supp(\rho_1)$\\[1pt]\scriptsize $\dim = r_1$};
\node[state, fill=green!8, right=3.0cm of r2, yshift=0cm] (U2)
  {$\Umin{2}{v} = \supp(\rho_2)$\\[1pt]\scriptsize $\dim = r_2$};
\node[state, fill=green!8, below=1.4cm of ralpha] (Ualpha)
  {$\Umin{\alpha}{v} = \supp(\rho_\alpha)$\\[1pt]\scriptsize $\dim = r_\alpha$};

% Arrows: pure state -> partial traces
\draw[arr, blue] (v) -- node[above, font=\scriptsize]{$\TrS{1^c}$} (r1);
\draw[arr, blue] (v) -- node[above, font=\scriptsize]{$\TrS{2^c}$} (r2);
\draw[arr, blue] (v) -- node[right, font=\scriptsize]{$\TrS{\alpha^c}$} (ralpha);

% Arrows: partial traces -> minimal subspaces
\draw[arr, orange!70!black] (r1)
  -- node[above, font=\scriptsize]{$\supp(\cdot)$} (U1);
\draw[arr, orange!70!black] (r2)
  -- node[above, font=\scriptsize]{$\supp(\cdot)$} (U2);
\draw[arr, orange!70!black] (ralpha)
  -- node[right, font=\scriptsize]{$\supp(\cdot)$} (Ualpha);

% Curved arrow: direct connection
\draw[arr, dashed, green!60!black] (v.south west)
  to[bend right=30]
  node[left, font=\scriptsize, green!50!black]{Def.~of $\Umin{\alpha}{\cdot}$}
  (U1.north east);

% Tucker manifold annotation at bottom
\node[label, below=0.3cm of Ualpha, align=center]
  {$v \in \bigotimes_\alpha \Umin{\alpha}{v}$,\quad
   $r_\alpha(v) = \mathrm{rank}(\rho_\alpha)$};
\end{tikzpicture}
\caption{The correspondence between minimal subspaces and partial traces
(Theorem~\ref{thm:qs:minimal_partial_trace}), illustrated for
$v \in \Tucker{\fr}{H} \subset H = \bigotimes_{j \in D} H_j$.
Blue arrows: formation of reduced density operators by partial tracing.
Orange arrows: extraction of minimal subspaces as supports of the
reduced operators. The dashed green arrow shows the direct definition
of $\Umin{\alpha}{v}$ as the smallest $U \subset H_\alpha$ with
$v \in U \otimes_H H_{\alpha^c}$.
The Tucker $\alpha$-rank $r_\alpha(v) = \dim\Umin{\alpha}{v}$ equals
the rank of the corresponding reduced density operator.}
\label{fig:minimal_partial_trace}
\end{figure}

%------------------------------------------------------------
\section{Mixed states and the rank stratification}
\label{sec:qs:mixed}
%------------------------------------------------------------

\subsection*{Hilbert--Schmidt operators as tensors}

By Theorem~\ref{thm:app:hs_tensor}, $\cT{2}(H) \cong H \otimes_H H^*$
isometrically as Hilbert spaces. Under this identification, a
Hilbert--Schmidt operator $A \in \cT{2}(H)$ of rank $r$ corresponds
to a tensor in $H \otimes_H H^*$ of Tucker rank $(r, r)$ with respect
to the bipartition $\{H, H^*\}$. The SVD of $A$,
\begin{equation*}
    A = \sum_{k=1}^r s_k(A)\, v_k \otimes J(u_k),
\end{equation*}
is the Tucker decomposition of this tensor, with minimal subspaces
$\Umin{1}{A} = \overline{\mathrm{Im}(A^*)} \subset H$ and
$\Umin{2}{A} = J(\overline{\mathrm{Im}(A)}) \subset H^*$.

\subsection*{Mixed states of rank $r$}

For a mixed state $\rho \in \Drank{r}{H}$ with spectral decomposition
\begin{equation}
\label{eq:qs:mixed_spectral}
    \rho = \sum_{k=1}^r \lambda_k\, v_k \otimes J(v_k),
    \quad \lambda_k > 0, \quad \sum_{k=1}^r \lambda_k = 1,
\end{equation}
the support is $\supp(\rho) = \spann\{v_k : k=1,\ldots,r\} \in \Grass{r}{H}$.
The positive core $\sigma \coloneqq V^*\rho V \in \Dpos{\CC^r}$
(where $V = [v_1|\cdots|v_r]: \CC^r \to H$ is the column matrix of
eigenvectors) is the finite-dimensional density operator on $\CC^r$
with eigenvalues $\lambda_1, \ldots, \lambda_r$.

\begin{proposition}[Rank stratification as Tucker decomposition]
\label{prop:qs:rank_tucker}
The decomposition
\begin{equation}
\label{eq:qs:rank_strat}
    \Dstate{H} = \bigsqcup_{r=1}^\infty \Drank{r}{H}
    \cup \mathcal{D}_\infty(H)
\end{equation}
is the rank stratification of $\Dstate{H}$ by Tucker rank in the
Banach tensor space $(\cTsa{1}(H), \|\cdot\|_1)$. Specifically, viewing
$\rho \in \Drank{r}{H}$ as a tensor in $H \otimes_\pi H^*$ via
$\cT{1}(H) \cong H \otimes_\pi H^*$:
\begin{enumerate}
\item[\textup{(i)}] $r_1(\rho) = r_2(\rho) = r = \mathrm{rank}(\rho)$.
\item[\textup{(ii)}] $\Umin{1}{\rho} = \supp(\rho) \subset H$ and
    $\Umin{2}{\rho} = J(\supp(\rho)) \subset H^*$.
\item[\textup{(iii)}] The trace norm $\|\rho\|_1 = \Tr(\rho) = 1$
    is the projective tensor norm, satisfying \condINJ\ with equality
    at all density operators \textup{(}Remark~\ref{rem:app:trace_projective}\textup{)}.
\end{enumerate}
\end{proposition}

\begin{proof}
Parts~(i)--(ii) follow from the identification $\cT{1}(H) \cong H
\otimes_\pi H^*$ and the definition of Tucker rank
(Definition~\ref{def:bounded_tucker_rank}): $r_\alpha(\rho)$ is the dimension
of $\Umin{\alpha}{\rho}$, which for $\alpha = 1$ is the support in $H$
and for $\alpha = 2$ is the support in $H^*$ (both have dimension $r$
by~\eqref{eq:qs:mixed_spectral}). Part~(iii) is
Remark~\ref{rem:app:trace_projective}.
\end{proof}

\begin{corollary}[Weak lower semicontinuity of quantum rank]
\label{cor:qs:rank_lsc}
In the trace-norm topology, rank is weakly lower semicontinuous on
$\Dstate{H}$: if $\rho_n \in \Drank{r}{H}$ and $\|\rho_n - \rho\|_1
\to 0$, then $\mathrm{rank}(\rho) \leq r$.
\end{corollary}

\begin{proof}
This is Theorem~\ref{thm:weak_closed} applied to the Tucker manifold
$\Drank{r}{H} \subset \cTsa{1}(H)$: the trace norm satisfies \condINJ\
(Remark~\ref{rem:app:trace_projective}), and $\cTsa{1}(H)$ is reflexive
(it is a closed subspace of the Banach space $\cT{1}(H)$, which is
reflexive when $H$ is finite-dimensional, and the statement holds in
infinite dimensions by the same argument with the weak topology on
$\cTsa{1}(H)$; see~\cite{BratteliRobinson1987}).
\end{proof}

%------------------------------------------------------------
\section{Entanglement and minimal subspace dimension}
\label{sec:qs:entanglement}
%------------------------------------------------------------

\subsection*{Separability and Tucker rank one}

For a multipartite system $H = \bigotimes_{j=1}^d H_j$ and a pure
state $\rho_v = v \otimes J(v)$, Theorem~\ref{thm:qs:minimal_partial_trace}
gives a complete characterisation of separability:

\begin{proposition}[Separability and Tucker rank]
\label{prop:qs:separability}
Let $v \in \bigotimes_{j=1}^d H_j$ with $\|v\|=1$ and let $\alpha
\subset D$ be a subset of parties. The following are equivalent:
\begin{enumerate}
\item[\textup{(a)}] $v$ is separable across the bipartition
    $\alpha | \alpha^c$: $v = u \otimes w$ for some $u \in
    \bigotimes_{j \in \alpha} H_j$, $w \in \bigotimes_{j \in \alpha^c} H_j$.
\item[\textup{(b)}] $r_\alpha(v) = 1$.
\item[\textup{(c)}] $\mathrm{rank}(\rho_\alpha) = 1$, i.e., the reduced
    state $\rho_\alpha = \TrS{\alpha^c}(\rho_v)$ is a pure state.
\item[\textup{(d)}] $\Umin{\alpha}{v}$ is one-dimensional.
\end{enumerate}
\end{proposition}

\begin{proof}
(a)$\Leftrightarrow$(b): by definition of Tucker rank
(Definition~\ref{def:bounded_tucker_rank}).
(b)$\Leftrightarrow$(c): Theorem~\ref{thm:qs:minimal_partial_trace}(b).
(b)$\Leftrightarrow$(d): $r_\alpha(v) = \dim\Umin{\alpha}{v}$.
\end{proof}

\subsection*{Entanglement entropy and minimal subspaces}

The \emph{entanglement entropy} of the bipartition $\alpha|\alpha^c$
is the von Neumann entropy of the reduced state:
\begin{equation}
\label{eq:qs:entanglement_entropy}
    S_\alpha(v) \coloneqq -\Tr\bigl(\rho_\alpha \log \rho_\alpha\bigr)
    = -\sum_{k=1}^{r_\alpha(v)} \lambda_k^{(\alpha)}(v)
    \log \lambda_k^{(\alpha)}(v),
\end{equation}
where $\{\lambda_k^{(\alpha)}(v)\}$ are the eigenvalues of $\rho_\alpha$
from Theorem~\ref{thm:qs:minimal_partial_trace}(c). Since
$\lambda_k^{(\alpha)} > 0$ and $\sum_k \lambda_k^{(\alpha)} = 1$,
we have $0 \leq S_\alpha(v) \leq \log r_\alpha(v)$, with:
\begin{itemize}
\item $S_\alpha(v) = 0$ iff $r_\alpha(v) = 1$ (separable, no entanglement).
\item $S_\alpha(v) = \log r_\alpha(v)$ iff all eigenvalues are equal
    $\lambda_k = 1/r_\alpha(v)$ (maximally entangled).
\end{itemize}
The Tucker rank $r_\alpha(v)$ is thus the \emph{Schmidt rank} of the
bipartition, which upper-bounds the entanglement entropy:
$S_\alpha(v) \leq \log r_\alpha(v)$.

\begin{remark}[The Tucker manifold as a variational class for entanglement]
\label{rem:qs:tucker_entanglement}
The Tucker manifold $\Tucker{\fr}{H}$ for $H = \bigotimes_j H_j$
and $\fr = (r_\alpha)_{\alpha \in D}$ is the set of pure states $v$
with prescribed bipartite Schmidt ranks. States on $\Tucker{\fr}{H}$
have entanglement entropies bounded by $S_\alpha \leq \log r_\alpha$
for all bipartitions $\alpha$. The Dirac--Frenkel principle of
Chapter~\ref{chap:DF}, applied on $\Tucker{\fr}{H}$, yields a
variational dynamics that preserves the entanglement structure
$(r_\alpha)_\alpha$ in the sense that the Schmidt ranks remain
equal to $\fr$ along the trajectory.
\end{remark}

%------------------------------------------------------------
\section{The Banach manifold structure of $\Drank{r}{H}$}
\label{sec:qs:manifold}
%------------------------------------------------------------

The geometry of $\Drank{r}{H}$ is a special case of the Tucker manifold
theory of Chapter~\ref{chap:tucker_manifold}, applied to the Banach
tensor space $(\cTsa{1}(H), \|\cdot\|_1) \cong (H \otimes_\pi H^*,
\|\cdot\|_\pi)$.

\subsection*{The atlas}

Fix $\rho_0 \in \Drank{r}{H}$ and let $S_0 = \supp(\rho_0) \in
\Grass{r}{H}$. Choose an isometry $V_0: \CC^r \to H$ with
$\mathrm{Im}(V_0) = S_0$, and set $P_0 = V_0 V_0^*$ (orthogonal
projection onto $S_0$). For $A \in \cL(\CC^r, S_0^\perp)$ with
$\|A\| < 1$, define the isometry
\begin{equation}
\label{eq:qs:VA}
    V_A \coloneqq V_0(I + A^*A)^{-1/2} + \iota A (I + A^*A)^{-1/2}
    \in \cL(\CC^r, H),
    \quad V_A^* V_A = I_{\CC^r},
\end{equation}
where $\iota: S_0^\perp \hookrightarrow H$ is the inclusion. The chart
map
\begin{equation}
\label{eq:qs:chart}
    \Phi_{S_0}: \cL(\CC^r, S_0^\perp) \times \Dpos{\CC^r}
    \longrightarrow \Drank{r}{H},
    \quad (A, \sigma) \longmapsto V_A \sigma V_A^*,
\end{equation}
is the quantum analogue of the Tucker chart map $\chartmap{\bv}$
of Section~\ref{sec:tucker_manifold:charts}: the Grassmann variable
$A \in \cL(\CC^r, S_0^\perp)$ encodes the motion of the support
(analogous to $L_\alpha \in \cL(\Umin{\alpha}{\bv}, \Wmin{\alpha}{\bv})$),
and the core variable $\sigma \in \Dpos{\CC^r}$ encodes the positive
density on the support (analogous to the core tensor $C^{(D)} \in
\RR^{\times r_\alpha}_*$).

\begin{theorem}[Banach manifold structure of $\Drank{r}{H}$,
{\cite[Theorem~3.5]{AcevedoFalco2026}}]
\label{thm:qs:manifold}
For each $r \in \NN$, the set $\Drank{r}{H}$ admits a smooth
$\mathcal{C}^\infty$-Banach manifold structure in the sense of
Lang~\cite{Lang1999}. Specifically:
\begin{enumerate}
\item[\textup{(i)}] For each $\rho_0 \in \Drank{r}{H}$, the map
    $\Phi_{S_0}$ in~\eqref{eq:qs:chart} is a $\mathcal{C}^\infty$
    diffeomorphism from a neighbourhood of $(0, V_0^*\rho_0 V_0)$ in
    $\cL(\CC^r, S_0^\perp) \times \Dpos{\CC^r}$ onto a neighbourhood
    of $\rho_0$ in $\Drank{r}{H}$.
\item[\textup{(ii)}] The transition maps between overlapping charts are
    $\mathcal{C}^\infty$ (analytic over $\CC$).
\item[\textup{(iii)}] The tangent space at $\rho \in \Drank{r}{H}$ is
    \begin{equation}
    \label{eq:qs:tangent}
        T_\rho \Drank{r}{H}
        = \bigl\{ \dot\rho \in \cTsa{1}(H) :
        (I - P)\dot\rho(I - P) = 0,\;
        \Tr(\dot\rho) = 0 \bigr\},
    \end{equation}
    where $P = \PS{\supp(\rho)}$ is the orthogonal projection onto
    $\supp(\rho)$. This is a complemented closed subspace of
    $\cTsa{1}(H)$.
\item[\textup{(iv)}] The inclusion $\iota: \Drank{r}{H} \hookrightarrow
    \cTsa{1}(H)$ is a smooth immersion.
\end{enumerate}
\end{theorem}

\begin{proof}
We prove each part in turn.

\medskip
\noindent\textbf{Part~(i): The chart map is a $\mathcal{C}^\infty$ diffeomorphism.}

Fix $\rho_0 \in \Drank{r}{H}$, $S_0 = \supp(\rho_0)$, and an isometry
$V_0: \CC^r \to H$ with $\mathrm{Im}(V_0) = S_0$.

\emph{Step~1: $\Phi_{S_0}$ is smooth.}
For $\|A\| < 1$, the operator $I_{\CC^r} + A^*A$ is invertible in
$\cL(\CC^r, \CC^r)$ (since $\CC^r$ is finite-dimensional and $A^*A \geq 0$).
The map $A \mapsto (I + A^*A)^{-1/2}$ is real-analytic on operator
norm balls, hence $A \mapsto V_A = (V_0 + \iota A)(I + A^*A)^{-1/2}$
is $\mathcal{C}^\infty$ by composition. Since $\sigma \mapsto V_A
\sigma V_A^*$ is linear in $\sigma$ for fixed $A$, the full map
$(A, \sigma) \mapsto V_A \sigma V_A^*$ is $\mathcal{C}^\infty$.

\emph{Step~2: Values in $\Drank{r}{H}$.}
For any $(A, \sigma)$ in the domain: $V_A \sigma V_A^* \geq 0$ (since
$\sigma > 0$ and $V_A$ is an isometry); $\Tr(V_A \sigma V_A^*) =
\Tr(\sigma V_A^* V_A) = \Tr(\sigma) = 1$ (since $V_A^* V_A =
I_{\CC^r}$); and $\mathrm{rank}(V_A \sigma V_A^*) = r$ (since
$\sigma \succ 0$ and $V_A$ is injective). Hence $\Phi_{S_0}(A,\sigma)
\in \Drank{r}{H}$.

\emph{Step~3: Local bijectivity and smooth inverse.}
At $A = 0$: $V_0 = V_0 (I)^{-1/2} = V_0$ and $\Phi_{S_0}(0,\sigma) =
V_0 \sigma V_0^*$; in particular $\Phi_{S_0}(0, V_0^* \rho_0 V_0) =
V_0 (V_0^* \rho_0 V_0) V_0^* = \rho_0$. For $\rho'$ near $\rho_0$
in $\Drank{r}{H}$, $\supp(\rho')$ is near $S_0$ in $\Grass{r}{H}$,
so there is a unique small $A$ with $\mathrm{Im}(V_A) = \supp(\rho')$.
The inverse is $\Phi_{S_0}^{-1}(\rho') = (A(\rho'), V_A^* \rho' V_A)$,
where $A(\rho') = \chi_{S_0, S_0^\perp}(\supp(\rho'))$ is the
Grassmannian chart map. Since the Grassmannian chart is $\mathcal{C}^\infty$
(Theorem~\ref{thm:grassmann_manifold}) and $\rho' \mapsto V_A^* \rho'
V_A$ is linear in $\rho'$, the inverse is $\mathcal{C}^\infty$.

\medskip
\noindent\textbf{Part~(ii): Smooth transition maps.}

Let $\rho_0, \tilde\rho_0 \in \Drank{r}{H}$ with chart domains
$\mathcal{B}(\rho_0)$ and $\mathcal{B}(\tilde\rho_0)$ overlapping.
The transition map $\Phi_{\tilde S_0}^{-1} \circ \Phi_{S_0}$ is the
composition:
\begin{equation*}
    (A, \sigma) \xmapsto{\Phi_{S_0}} \rho
    \xmapsto{\Phi_{\tilde S_0}^{-1}} (\tilde A, \tilde\sigma),
\end{equation*}
where $\tilde A = \chi_{\tilde S_0, \tilde S_0^\perp}(\supp(\rho))$
is smooth in $\rho$ (Grassmannian transition map, Theorem~\ref{thm:grassmann_manifold}),
and $\tilde\sigma = \tilde V_{\tilde A}^* \rho \tilde V_{\tilde A}$
is smooth in both $\tilde A$ and $\rho$. The composition is therefore
$\mathcal{C}^\infty$.

\medskip
\noindent\textbf{Part~(iii): Tangent space characterisation.}

Let $t \mapsto \rho(t) = \Phi_{S_0}(A(t), \sigma(t))$ be a
$\mathcal{C}^1$ curve with $A(0) = 0$ and $\sigma(0) = V_0^* \rho_0
V_0$. Differentiating $\rho(t) = V_{A(t)} \sigma(t) V_{A(t)}^*$ at
$t=0$:
\begin{equation*}
    \dot\rho(0) = \dot V_0 \sigma_0 V_0^* + V_0 \dot\sigma_0 V_0^*
    + V_0 \sigma_0 \dot V_0^*,
\end{equation*}
where $\dot V_0 = \frac{d}{dt}V_{A(t)}|_{t=0}$. Since $V_{A(t)}^*
V_{A(t)} = I$ implies $\dot V_0^* V_0 + V_0^* \dot V_0 = 0$ (so
$V_0^* \dot V_0$ is skew-Hermitian), and $\mathrm{Im}(V_0) = S_0$
(so $\dot V_0 \in \cL(\CC^r, S_0^\perp)$ by the Grassmannian chart),
one computes:
\begin{equation*}
    (I - P)\dot\rho(0)(I - P)
    = (I-P)\dot V_0 \sigma_0 V_0^*(I-P) = 0,
\end{equation*}
since $V_0^*(I-P) = 0$ (the range of $V_0$ is $S_0$). Also
$\Tr(\dot\rho(0)) = \Tr(\dot\sigma_0) + 2\mathrm{Re}\,\Tr(\sigma_0
V_0^* \dot V_0) = 0$ if $\Tr(\dot\sigma_0) = 0$ and $V_0^*\dot V_0$
is skew-Hermitian. This shows the inclusion
$\dot\rho(0) \in T_\rho \Drank{r}{H} \supset$ the set in~\eqref{eq:qs:tangent}.
The reverse inclusion follows by constructing explicit curves realising
each tangent vector of the given form (Step~3 of the proof
in~\cite{AcevedoFalco2026}).

\medskip
\noindent\textbf{Part~(iv): Smooth immersion.}

The differential of $\iota: \Drank{r}{H} \hookrightarrow \cTsa{1}(H)$
at $\rho$ is the map $d\iota_\rho: T_\rho \Drank{r}{H} \to \cTsa{1}(H)$,
$d\iota_\rho(\dot\rho) = \dot\rho$, which is injective by definition.
The image $T_\rho \Drank{r}{H}$ is a complemented subspace of
$\cTsa{1}(H)$ by the explicit bounded projection $\hat P_\rho$ of
Proposition~\ref{prop:qs:projector} below, satisfying
$\|\hat P_\rho\| \leq 4 < \infty$. Hence $d\iota_\rho$ has closed
split image (Proposition~\ref{prop:split_equiv}), i.e., $\iota$ is
an immersion.
\end{proof}

\subsection*{The tangent projector}

\begin{proposition}[Bounded projector onto $T_\rho \Drank{r}{H}$,
{\cite[Proposition~3.4]{AcevedoFalco2026}}]
\label{prop:qs:projector}
Let $\rho \in \Drank{r}{H}$ and $P = \PS{\supp(\rho)}$. The map
\begin{equation}
\label{eq:qs:projector}
    \hat{P}_\rho: \cTsa{1}(H) \longrightarrow T_\rho \Drank{r}{H},
    \quad
    \hat{P}_\rho(X) \coloneqq PXP + PX(I-P) + (I-P)XP
    - \frac{\Tr(PXP)}{r} P,
\end{equation}
is a bounded linear projection with $\|\hat{P}_\rho\|_{\cL(\cTsa{1}(H))}
\leq 4$, and $\hat{P}_\rho(X) = X$ for all $X \in T_\rho \Drank{r}{H}$.
\end{proposition}

\begin{proof}
Write $X \in \cTsa{1}(H)$ in block form with respect to $H = S \oplus S^\perp$:
\begin{equation*}
    X = \begin{pmatrix} X_{11} & X_{12}^* \\ X_{12} & X_{22}
    \end{pmatrix},
    \quad X_{11} \in \cTsa{1}(S),\; X_{12} \in \cT{1}(S,S^\perp),\;
    X_{22} \in \cTsa{1}(S^\perp).
\end{equation*}
By direct computation using $P^2 = P = P^*$:
\begin{align*}
    PXP &= \begin{pmatrix} X_{11} & 0 \\ 0 & 0 \end{pmatrix}, &
    PX(I-P) &= \begin{pmatrix} 0 & X_{12}^* \\ 0 & 0 \end{pmatrix}, \\
    (I-P)XP &= \begin{pmatrix} 0 & 0 \\ X_{12} & 0 \end{pmatrix}, &
    \frac{\Tr(PXP)}{r}P &= \frac{\Tr(X_{11})}{r}
    \begin{pmatrix} I_S & 0 \\ 0 & 0 \end{pmatrix}.
\end{align*}
Hence
\begin{equation}
\label{eq:qs:projector_block}
    \hat P_\rho(X)
    = \begin{pmatrix} X_{11} - \frac{\Tr(X_{11})}{r}I_S & X_{12}^* \\
    X_{12} & 0 \end{pmatrix}.
\end{equation}
\emph{Idempotency:} Applying $\hat P_\rho$ to~\eqref{eq:qs:projector_block}:
the $(S,S)$ block is $X_{11} - \frac{\Tr(X_{11})}{r}I_S - \frac{1}{r}
\Tr(X_{11} - \frac{\Tr(X_{11})}{r}I_S)I_S = X_{11} - \frac{\Tr(X_{11})}{r}I_S$
(since $\Tr(X_{11} - \frac{\Tr(X_{11})}{r}I_S) = \Tr(X_{11}) - \Tr(X_{11}) = 0$),
and the off-diagonal blocks are unchanged. So $\hat P_\rho^2 = \hat P_\rho$.

\emph{Image:} From~\eqref{eq:qs:projector_block}, $\hat P_\rho(X)$ has
$(S^\perp,S^\perp)$ block zero and $\Tr(\hat P_\rho(X)) = \Tr(X_{11}
- \frac{\Tr(X_{11})}{r}I_S) + 0 = 0$. So $\hat P_\rho(X) \in
T_\rho \Drank{r}{H}$.

\emph{$\hat P_\rho = \mathrm{Id}$ on $T_\rho\Drank{r}{H}$:} If $Y \in
T_\rho\Drank{r}{H}$, its block form is $Y_{22} = 0$ and $\Tr(Y_{11})=0$
(by~\eqref{eq:qs:tangent}). Then $\frac{\Tr(PYP)}{r} = \frac{\Tr(Y_{11})}{r} = 0$,
so $\hat P_\rho(Y) = Y$.

\emph{Norm bound:} Using the ideal property $\|BAC\|_1 \leq \|B\|\|A\|_1\|C\|$:
\begin{align*}
    \|\hat P_\rho(X)\|_1
    &\leq \|PXP\|_1 + \|PX(I-P)\|_1 + \|(I-P)XP\|_1
    + \tfrac{|\Tr(PXP)|}{r}\|P\|_1 \\
    &\leq \|X\|_1 + \|X\|_1 + \|X\|_1
    + \tfrac{\|PXP\|_1}{r} \cdot r \\
    &\leq 4\|X\|_1,
\end{align*}
where we used $|\Tr(PXP)| \leq \|PXP\|_1 \leq \|X\|_1$ and
$\|P\|_1 = \Tr(P) = r$.
\end{proof}

\begin{remark}[Comparison with the Tucker projector]
\label{rem:qs:projector_comparison}
The projector $\hat{P}_\rho$ in~\eqref{eq:qs:projector} is the
quantum analogue of the metric projection $\Pmetric{\bv_r}$ onto
$\ZD{\bv_r}$ from Chapter~\ref{chap:DF}. The explicit formula in
terms of blocks with respect to $H = S \oplus S^\perp$ is the
quantum version of the tangent space formula~\eqref{eq:tangent_formula}:
the $PXP$ block corresponds to the core variation $\dot{C}^{(D)}$,
and the off-diagonal blocks correspond to the subspace velocity
$\dot{u}^{(\alpha)}_{i_\alpha} \in \Wmin{\alpha}{\bv_r}$.
The trace correction $-\frac{\Tr(PXP)}{r}P$ is the additional
constraint $\Tr(\dot\rho) = 0$ that appears because $\Drank{r}{H}$
lies in the affine hyperplane $\{\Tr = 1\}$, not in the full Banach
space.
\end{remark}

\subsection*{The fibre bundle structure}

The support map $\pi: \Drank{r}{H} \to \Grass{r}{H}$,
$\pi(\rho) = \supp(\rho)$, is the quantum analogue of the rank map
$\varrho_\fr$ of Section~\ref{sec:tucker_manifold:cinfty}.

\begin{theorem}[Fibre bundle, {\cite[Theorem~3.8]{AcevedoFalco2026}}]
\label{thm:qs:bundle}
The support map $\pi: \Drank{r}{H} \to \Grass{r}{H}$ is a smooth
fibre bundle with typical fibre $\Dpos{\CC^r}$, a smooth manifold of
real dimension $r^2 - 1$. The structure group is $U(r)$ (the unitary
group), acting on $\Dpos{\CC^r}$ by $U \cdot \sigma = U \sigma U^*$.
\end{theorem}

This is the quantum specialisation of Theorem~\ref{thm:Tucker_Banach_Manifold}(ii)
and Propositions~\ref{prop:tucker_triv}--\ref{prop:tucker_transition_fns}:
the structure group reduces from $\GL{r}{\RR}$ to $U(r)$ because the
core variable $\sigma \in \Dpos{\CC^r}$ transforms by unitary
conjugation (preserving positivity and trace), rather than by
an arbitrary invertible change of basis.

%------------------------------------------------------------
\section{Best approximation and quantum state compression}
\label{sec:qs:approx}
%------------------------------------------------------------

The best-approximation theorem (Theorem~\ref{thm:best_approx_tucker})
applied to the quantum setting gives the mathematical foundation for
\emph{low-rank approximation of quantum states}.

\begin{theorem}[Existence of best low-rank approximation]
\label{thm:qs:best_approx}
Let $\rho \in \Dstate{H}$ and $r \in \NN$. There exists
$\rho_r \in \bigcup_{k=1}^r \Drank{k}{H}$ such that
\begin{equation}
\label{eq:qs:best_approx}
    \|\rho - \rho_r\|_1
    = \inf_{\sigma \in \bigcup_{k=1}^r \Drank{k}{H}} \|\rho - \sigma\|_1.
\end{equation}
\end{theorem}

\begin{proof}
The set $\bigcup_{k=1}^r \Drank{k}{H}$ is weakly closed in
$\cTsa{1}(H)$ by Corollary~\ref{cor:qs:rank_lsc} (weak lower
semicontinuity of rank). The trace-class space $\cTsa{1}(H)$ is
reflexive when $H$ is finite-dimensional; in infinite dimensions,
the result follows from the explicit formula: the best rank-$r$
approximation of $\rho$ in trace norm is given by the \emph{spectral
truncation}
\begin{equation}
\label{eq:qs:spectral_trunc}
    \rho_r \coloneqq \frac{1}{\sum_{k=1}^r \lambda_k}
    \sum_{k=1}^r \lambda_k\, v_k \otimes J(v_k),
\end{equation}
where $\lambda_1 \geq \lambda_2 \geq \cdots$ are the eigenvalues of
$\rho$ in decreasing order and $\{v_k\}$ the corresponding eigenvectors.
The error is $\|\rho - \rho_r\|_1 = 2\sum_{k > r} \lambda_k$
(Ky Fan's theorem).
\end{proof}

\begin{remark}[Spectral truncation as Tucker approximation]
\label{rem:qs:ky_fan}
The spectral truncation~\eqref{eq:qs:spectral_trunc} is the quantum
analogue of the HOSVD (Higher-Order SVD) truncation for tensor formats
(Algorithm~\ref{app:alg:hosvd} in Appendix~\ref{app:algorithms}):
in both cases, the best rank-$r$ approximation in the ambient norm is
obtained by projecting onto the dominant eigenspaces/singular subspaces.
For tensors with $d \geq 3$, the analogue of~\eqref{eq:qs:spectral_trunc}
gives a near-best Tucker approximation (quasi-optimality with
constant $\sqrt{d}$), not an exact best approximation; see
Appendix~\ref{app:alg:hosvd}.
\end{remark}

%------------------------------------------------------------
\section{Historical notes}
\label{sec:qs:history}
%------------------------------------------------------------

\paragraph{Schmidt decomposition and entanglement.}
The Schmidt decomposition for bipartite systems was introduced by
Schmidt~(1907) in the context of integral equations, long before its
quantum-mechanical interpretation. Its role in quantum entanglement
was recognised in the 1990s with the development of quantum information
theory. The generalisation to multipartite systems via minimal subspaces
is the contribution of Falcó and Hackbusch~\cite{FH2012}.

\paragraph{Tucker format and quantum many-body physics.}
The connection between the Tucker format and tree tensor networks (TTN)
was developed in the quantum chemistry and condensed matter literature
under the name MCTDH (Multi-Configuration Time-Dependent Hartree) and
its matrix product state (MPS/DMRG) variants. White's DMRG algorithm
\cite{White1992} is algorithmically the DMRG of
Appendix~\ref{app:algorithms}, implemented for quantum spin chains.
The mathematical equivalence between Tucker formats and tensor network
states was clarified in~\cite{Hackbusch2019}.

\paragraph{Banach manifold structure of $\Drank{r}{H}$.}
The explicit Grassmann-Banach manifold model for fixed-rank density
operators in the trace-norm topology, with the support-core charts and
bounded tangent projectors, was developed in the companion
paper~\cite{AcevedoFalco2026}. Earlier work in the Hilbert--Schmidt
(T$_2$) setting includes~\cite{KochLubich2007}; the move to the
trace-class (T$_1$) setting is motivated by the operational
interpretation of the trace norm as the natural metric for state
distinguishability (Remark~\ref{rem:app:trace_projective}).

%============================================================
%  ch13_lindblad.tex
%  Chapter 13: The Dirac--Frenkel Principle for Open Quantum
%  Systems
%
%  Tensor Formats in Banach Spaces
%  Antonio Falcó, CEU Universities
%
%  Based on: Acevedo--Falcó (2026), FHN2019
%  Estimated length: ~30 pp.
%============================================================

\chapter{The Dirac--Frenkel Principle for Open Quantum Systems}
\label{chap:lindblad}

Chapter~\ref{chap:DF} established the Dirac--Frenkel variational
principle for abstract Cauchy problems on tensor Banach manifolds.
Chapter~\ref{chap:quantum_states} identified the fixed-rank stratum
$\Drank{r}{H}$ as a smooth Banach manifold within the trace-class
Banach space $\cTsa{1}(H)$, inheriting all the geometric properties
developed in Chapters~\ref{chap:tucker_manifold}--\ref{chap:TB_manifold}.
This chapter combines these two lines to obtain the Dirac--Frenkel
principle for quantum master equations on $\cTsa{1}(H)$: the Lindblad
(GKSL) equation for open quantum systems.

The Lindblad equation describes the evolution of the state $\rho(t)
\in \Dstate{H}$ of a quantum system coupled to an environment. Its
generator $\Lind$ maps $\cTsa{1}(H)$ to itself and preserves
$\Dstate{H}$. Since exact simulation of $\rho(t)$ in
infinite-dimensional systems is generally intractable, one seeks a
reduced dynamics on the finite-rank manifold $\Drank{r}{H}$,
obtained by projecting the Lindblad vector field onto the tangent space
of $\Drank{r}{H}$.

Section~\ref{sec:lind:equation} formulates the Lindblad equation as an
ODE on $\cTsa{1}(H)$ and identifies its vector field.
Section~\ref{sec:lind:tangent} develops the tangent space of $\Drank{r}{H}$
via minimal subspaces and establishes the bounded projector.
Section~\ref{sec:lind:df} states and proves the main well-posedness theorem.
Section~\ref{sec:lind:preservation} analyses which structural properties
(trace, rank, positivity) are preserved by the projected dynamics.
Section~\ref{sec:lind:kinematic} derives the kinematic equations in
factorised coordinates, recovering the equations of Chapter~\ref{chap:DF}
as a special case.
Section~\ref{sec:lind:history} contains historical notes.

%------------------------------------------------------------
\section{The Lindblad equation on $\cTsa{1}(H)$}
\label{sec:lind:equation}
%------------------------------------------------------------

\subsection*{The GKSL generator}

The most general Markovian master equation for the state $\rho(t)
\in \Dstate{H}$ of an open quantum system is the
\emph{Gorini--Kossakowski--Sudarshan--Lindblad} (GKSL) equation:
\begin{equation}
\label{eq:lind:GKSL}
    \dot\rho(t) = \Lind(\rho(t))
    \coloneqq -i[H, \rho(t)]
    + \sum_{j \in J} \Bigl(
    L_j \rho(t) L_j^*
    - \tfrac{1}{2}\{L_j^* L_j, \rho(t)\}
    \Bigr),
\end{equation}
where:
\begin{itemize}
\item $H = H^* \in \cL(H)$ is the \emph{Hamiltonian} (self-adjoint,
    bounded);
\item $\{L_j\}_{j \in J} \subset \cL(H)$ are the \emph{Lindblad
    operators} (noise operators), with $J$ a finite or countable index set;
\item $[H, \rho] = H\rho - \rho H$ is the commutator and
    $\{A, B\} = AB + BA$ is the anticommutator.
\end{itemize}
When $J$ is finite and all $L_j \in \cL(H)$ are bounded, $\Lind$ is a
bounded operator on $\cT{1}(H)$. The solution to~\eqref{eq:lind:GKSL}
is a strongly continuous one-parameter semigroup $(T_t)_{t \geq 0}$ on
$\cT{1}(H)$ that is \emph{completely positive and trace preserving}
(CPTP): $T_t(\Dstate{H}) \subset \Dstate{H}$ for all $t \geq 0$.

\begin{remark}[The Lindblad equation as a tensor Banach ODE]
\label{rem:lind:tensor_ODE}
The equation~\eqref{eq:lind:GKSL} is a special case of the abstract
Cauchy problem~\eqref{eq:DF:ode} of Chapter~\ref{chap:DF}, with ambient
Banach space $\VDnorm = \cTsa{1}(H)$ and vector field
$\bF(t, \rho) = \Lind(\rho)$ (time-independent in the homogeneous case).
The trace norm $\|\cdot\|_1$ plays the role of $\|\cdot\|_D$, and the
identification $\cTsa{1}(H) \cong H \otimes_\pi H^*$ (Remark~\ref{rem:app:trace_projective})
makes $(\cTsa{1}(H), \|\cdot\|_1)$ a Banach tensor space with $d=2$,
$V_1 = H$, $V_2 = H^*$, satisfying condition~\condINJ\ with the
projective norm.
\end{remark}

\begin{assumption}[Regularity of $\Lind$]
\label{ass:lind:regular}
Throughout this chapter we assume that:
\begin{enumerate}
\item[\textup{(R1)}] $H \in \cL(H)$ is bounded and $J$ is finite,
    so $\Lind: \cTsa{1}(H) \to \cTsa{1}(H)$ is a bounded linear operator.
\item[\textup{(R2)}] The CPTP semigroup $(T_t)_{t \geq 0}$ generated
    by $\Lind$ satisfies $T_t(\Dstate{H}) \subset \Dstate{H}$.
\item[\textup{(R3)}] For each $\rho_0 \in \Drank{r}{H}$, the restriction
    $\Lind|_{\Drank{r}{H}}$ is locally Lipschitz in $\|\cdot\|_1$.
\end{enumerate}
\end{assumption}

Under Assumption~\ref{ass:lind:regular}, the full Lindblad equation has
a unique global solution $\rho: [0,\infty) \to \Dstate{H}$.

\begin{remark}[Unbounded generators]
\label{rem:lind:unbounded}
Assumption~\ref{ass:lind:regular}(R1) excludes unbounded Lindblad
operators (e.g., position or momentum operators in infinite dimensions).
The extension to unbounded generators requires domain-adapted fixed-rank
strata and is the subject of Open Problem~\ref{prob:lindblad_unbounded}.
\end{remark}

%------------------------------------------------------------
\section{Tangent space of $\Drank{r}{H}$ via minimal subspaces}
\label{sec:lind:tangent}
%------------------------------------------------------------

The tangent space characterisation of Theorem~\ref{thm:qs:manifold}(iii)
takes the following explicit form, connecting with the Tucker tangent
space of Chapter~\ref{chap:tucker_manifold}.

\begin{proposition}[Tangent space decomposition]
\label{prop:lind:tangent}
Let $\rho \in \Drank{r}{H}$, $S = \supp(\rho)$, $P = \PS{S}$.
With respect to the orthogonal decomposition $H = S \oplus S^\perp$,
every $\dot\rho \in T_\rho \Drank{r}{H}$ has a unique block form
\begin{equation}
\label{eq:lind:tangent_block}
    \dot\rho =
    \begin{pmatrix} C & B^* \\ B & 0 \end{pmatrix}
    \in \cTsa{1}(H),
    \quad C \in \cTsa{1}(S),\; \Tr(C) = 0,\;
    B \in \cT{1}(S, S^\perp),
\end{equation}
where $\cT{1}(S, S^\perp)$ denotes the trace-class operators from
$S$ to $S^\perp$ (automatically finite-rank since $\dim S = r$).
The decomposition $T_\rho \Drank{r}{H} \cong \{C \in \cTsa{1}(S) :
\Tr(C) = 0\} \oplus \cT{1}(S, S^\perp)$ identifies:
\begin{itemize}
\item The \emph{core direction} $C$: variation of $\rho$ within fixed
    support $S$ — the quantum analogue of $\dot{C}^{(D)}$ in~\eqref{eq:DF:tangent_decomp}.
\item The \emph{support direction} $B$: variation of the support
    subspace $S$ — the quantum analogue of $\dot{u}^{(\alpha)}_{i_\alpha}
    \in \Wmin{\alpha}{\bv_r}$ in~\eqref{eq:DF:tangent_decomp}.
\end{itemize}
\end{proposition}

\begin{proof}
This is Lemma~3.1 and Corollary~3.2 of~\cite{AcevedoFalco2026},
proved by differentiating the chart curve $(A(t), \sigma(t))
\mapsto V_{A(t)}\sigma(t)V_{A(t)}^*$ at $t=0$ in local chart
coordinates~\eqref{eq:qs:chart}. The block structure follows from
the orthogonality of $H = S \oplus S^\perp$ and the chain rule,
exactly as in the proof of Proposition~\ref{prop:tangent_space}.
\end{proof}

\begin{remark}[The trace constraint]
\label{rem:lind:trace_constraint}
The condition $\Tr(\dot\rho) = 0$ in~\eqref{eq:lind:tangent_block}
is the infinitesimal form of $\Tr(\rho(t)) = 1$. It corresponds to
the fact that $\Drank{r}{H}$ lies in the affine hyperplane
$\{\Tr = 1\} \subset \cTsa{1}(H)$, which is an additional constraint
not present in the Tucker manifold $\Tucker{\fr}{\VD} \subset \VDnorm$.
In the kinematic equations of Section~\ref{sec:lind:kinematic},
this constraint is enforced automatically by the trace-removal term
in the projector $\hat{P}_\rho$ (Proposition~\ref{prop:qs:projector}).
\end{remark}

%------------------------------------------------------------
\section{The Dirac--Frenkel principle for Lindblad dynamics}
\label{sec:lind:df}
%------------------------------------------------------------

\begin{definition}[Dirac--Frenkel projected Lindblad dynamics]
\label{def:lind:df}
The \emph{Dirac--Frenkel projected Lindblad dynamics} on $\Drank{r}{H}$
is the ODE
\begin{equation}
\label{eq:lind:df}
    \dot\rho(t) = \hat{P}_{\rho(t)}\bigl(\Lind(\rho(t))\bigr),
    \qquad \rho(0) = \rho_0 \in \Drank{r}{H},
\end{equation}
where $\hat{P}_\rho: \cTsa{1}(H) \to T_\rho \Drank{r}{H}$ is the
bounded projector of Proposition~\ref{prop:qs:projector}.
\end{definition}

The equation~\eqref{eq:lind:df} is the quantum specialisation of the
Dirac--Frenkel equation~\eqref{eq:DF_tucker} of Chapter~\ref{chap:DF},
with $\Tucker{\fr}{\VD}$ replaced by $\Drank{r}{H}$, $\VDnorm$ by
$\cTsa{1}(H)$, $\bF$ by $\Lind$, and $\Pmetric{\bv_r}$ by $\hat{P}_\rho$.

\begin{theorem}[Local well-posedness,
{\cite[Theorem~4.6]{AcevedoFalco2026}}]
\label{thm:lind:wellposed}
Under Assumption~\ref{ass:lind:regular}, for every
$\rho_0 \in \Drank{r}{H}$ there exist $T > 0$ and a unique
$\mathcal{C}^1$-curve
\begin{equation*}
    \rho: [0, T) \longrightarrow \Drank{r}{H}
\end{equation*}
satisfying~\eqref{eq:lind:df} with $\rho(0) = \rho_0$.
The solution depends continuously on $\rho_0$ in the trace norm.
\end{theorem}

\begin{proof}
We verify the three conditions of Theorem~\ref{thm:app:ode} (Picard--Lindelöf
on Banach manifolds, Appendix~\ref{app:manifolds}).

\medskip
\noindent\textbf{Step~1: $\Drank{r}{H}$ is a $\mathcal{C}^\infty$ Banach manifold.}

This is Theorem~\ref{thm:qs:manifold}. The model space for a chart at
$\rho_0$ is $\cL(\CC^r, S_0^\perp) \times \cTsa{1}(\CC^r)$, a Banach
space (infinite-dimensional in the first factor, finite-dimensional in
the second).

\medskip
\noindent\textbf{Step~2: The reduced vector field is locally Lipschitz.}

In the chart $\Phi_{S_0}$ at $\rho_0$, the reduced vector field
in coordinates is
\begin{equation*}
    F(A, \sigma)
    \coloneqq d\Phi_{S_0}^{-1}\bigl(\Phi_{S_0}(A,\sigma)\bigr)
    \bigl[\hat P_{\rho(A,\sigma)}\bigl(\Lind(\rho(A,\sigma))\bigr)\bigr],
\end{equation*}
where $\rho(A,\sigma) = V_A \sigma V_A^*$.

We verify local Lipschitz continuity in three sub-steps:

\emph{(a) $\Lind(\rho)$ is Lipschitz in $\|\cdot\|_1$.}
By Assumption~\ref{ass:lind:regular}(R1), $\Lind$ is bounded on
$\cT{1}(H)$: $\|\Lind(\rho) - \Lind(\rho')\|_1 \leq \|\Lind\|_{\cL(\cT{1})}
\|\rho - \rho'\|_1$.

\emph{(b) $\rho \mapsto \hat P_\rho$ is locally Lipschitz.}
Writing $\hat P_\rho(X) = P_{\supp(\rho)} X P_{\supp(\rho)} + \ldots$,
the dependence on $\rho$ enters only through the orthogonal projection
$P = P_{\supp(\rho)}$. In the chart, $\supp(\rho(A,\sigma)) =
\mathrm{Im}(V_A)$ and $P = V_A V_A^*$. The map $A \mapsto V_A V_A^*$
is $\mathcal{C}^\infty$ by Step~1 (Theorem~\ref{thm:qs:manifold}(i)),
hence Lipschitz on bounded sets.

\emph{(c) The chart map is Lipschitz.}
$d\Phi_{S_0}^{-1}$ is bounded on bounded sets since $\Phi_{S_0}$ is a
$\mathcal{C}^\infty$ diffeomorphism with bounded differential.

Combining (a)--(c) by the chain rule: $F$ is locally Lipschitz on
a neighbourhood of $(0, V_0^*\rho_0 V_0)$ in $\cL(\CC^r, S_0^\perp)
\times \cTsa{1}(\CC^r)$.

\medskip
\noindent\textbf{Step~3: Local existence and uniqueness by Picard--Lindelöf.}

By Theorem~\ref{thm:app:picard} (Picard--Lindelöf in Banach spaces),
applied to the coordinate ODE $\dot z(t) = F(z(t))$ in the Banach space
$\cL(\CC^r, S_0^\perp) \times \cTsa{1}(\CC^r)$, there exist $T > 0$ and
a unique $\mathcal{C}^1$ solution $z: [0,T) \to \cL(\CC^r, S_0^\perp)
\times \cTsa{1}(\CC^r)$ with $z(0) = (0, V_0^*\rho_0 V_0)$.
Pushing forward by $\Phi_{S_0}$, we obtain $\rho(t) = \Phi_{S_0}(z(t))$,
a $\mathcal{C}^1$ curve in $\Drank{r}{H}$ satisfying~\eqref{eq:lind:df}.

\medskip
\noindent\textbf{Step~4: Continuous dependence on $\rho_0$.}

Continuous dependence of $z(t)$ on $z(0)$ follows from Theorem~\ref{thm:app:picard}.
Transported through the $\mathcal{C}^\infty$ chart map $\Phi_{S_0}$,
this gives continuous dependence of $\rho(t)$ on $\rho_0$ in the
trace norm.
\end{proof}

\begin{remark}[Choice of projector]
\label{rem:lind:projector_choice}
The projector $\hat{P}_\rho$ in Definition~\ref{def:lind:df} is the
explicit block projector of Proposition~\ref{prop:qs:projector}. Other
choices of bounded projector $\pi_\rho: \cTsa{1}(H) \to T_\rho
\Drank{r}{H}$ lead to equivalent well-posedness results; the choice
affects the trajectory $\rho(t)$ but not the structural preservation
properties. In particular, in the Hilbert--Schmidt setting one may use
the orthogonal projection in $\cT{2}(H)$, which gives the TDVP
(time-dependent variational principle) projection of~\cite{LeBrisRouchon2013}.
\end{remark}

%------------------------------------------------------------
\section{Structure preservation}
\label{sec:lind:preservation}
%------------------------------------------------------------

\begin{proposition}[Trace, rank, and positivity]
\label{prop:lind:preservation}
Let $\rho: [0,T) \to \Drank{r}{H}$ be the solution of~\eqref{eq:lind:df}.
Then for all $t \in [0, T)$:
\begin{enumerate}
\item[\textup{(i)}] \emph{Trace preservation:}
    $\Tr(\rho(t)) = 1$.
\item[\textup{(ii)}] \emph{Rank preservation:}
    $\mathrm{rank}(\rho(t)) = r$.
\item[\textup{(iii)}] \emph{Positivity:}
    $\rho(t) \geq 0$.
\end{enumerate}
\end{proposition}

\begin{proof}
\textup{(i)}: Differentiating, $\frac{d}{dt}\Tr(\rho(t)) =
\Tr(\dot\rho(t)) = \Tr(\hat{P}_{\rho(t)}(\Lind(\rho(t)))) = 0$,
since $\hat{P}_\rho(X) \in T_\rho \Drank{r}{H}$ satisfies
$\Tr(\hat{P}_\rho(X)) = 0$ by~\eqref{eq:lind:tangent_block}.

\textup{(ii)} and \textup{(iii)}: The solution $\rho(t)$ is a curve
in $\Drank{r}{H}$ by Theorem~\ref{thm:lind:wellposed}. Since
$\Drank{r}{H} \subset \Dstate{H}$, every point satisfies
$\rho(t) \geq 0$ and $\mathrm{rank}(\rho(t)) = r$.
\end{proof}

\begin{remark}[Complete positivity is not guaranteed]
\label{rem:lind:cptp}
Proposition~\ref{prop:lind:preservation} guarantees that $\rho(t)
\in \Drank{r}{H} \subset \Dstate{H}$ (positive, trace one, rank $r$)
along the trajectory. It does \emph{not} guarantee that the reduced
flow $\rho_0 \mapsto \rho(t)$ is a CPTP map (quantum channel).
The projected dynamics is a nonlinear ODE on $\Drank{r}{H}$; the
linearity and complete positivity of the full Lindblad semigroup
$(T_t)_{t \geq 0}$ are not inherited by the projected flow.
Constructing a CPTP-preserving rank-$r$ approximation requires
additional structure, such as Kraus-form parametrisations;
see Open Problem~\ref{prob:lindblad_cptp}.
\end{remark}

%------------------------------------------------------------
\section{Kinematic equations in factorised coordinates}
\label{sec:lind:kinematic}
%------------------------------------------------------------

\subsection*{Stiefel-core parametrisation}

Every $\rho \in \Drank{r}{H}$ admits a factorisation
\begin{equation}
\label{eq:lind:stiefel_core}
    \rho = V \sigma V^*,
    \quad V \in \Stiefel{r}{H},\;
    \sigma \in \Dpos{\CC^r},
\end{equation}
where $\Stiefel{r}{H} = \{V \in \cL(\CC^r, H) : V^*V = I_{\CC^r}\}$
is the Stiefel manifold of isometries. The parametrisation is unique up
to the unitary gauge action $(V, \sigma) \sim (VU, U^*\sigma U)$ for
$U \in U(r)$.

Differentiating~\eqref{eq:lind:stiefel_core}:
\begin{equation}
\label{eq:lind:diff_stiefel}
    \dot\rho = \dot{V} \sigma V^* + V \dot\sigma V^* + V \sigma \dot{V}^*.
\end{equation}
The Stiefel constraint $V^*V = I$ differentiates to give $V^*\dot{V} +
\dot{V}^*V = 0$, so $V^*\dot{V}$ is skew-Hermitian. Imposing the
\emph{horizontal (parallel transport) gauge}
\begin{equation}
\label{eq:lind:gauge}
    V^* \dot{V} = 0
    \quad\Longleftrightarrow\quad
    \dot{V} = (I - VV^*) \dot{V},
\end{equation}
the variation $\dot{V}$ lies in the complement of $\mathrm{Im}(V) = S$,
i.e., $\dot{V}(\CC^r) \subset S^\perp$.

\subsection*{Derivation of the kinematic system}

Under the gauge condition~\eqref{eq:lind:gauge}, projecting the
Lindblad vector field $\Lind(\rho)$ onto $T_\rho \Drank{r}{H}$ via
$\hat{P}_\rho$ gives, using the block formula~\eqref{eq:qs:projector}:

\medskip
\noindent\textbf{Equation for the core $\sigma(t)$.}
Multiplying $\dot\rho = \hat{P}_\rho(\Lind(\rho))$ on the left by $V^*$
and on the right by $V$, and using $V^*\dot{V} = 0$:
\begin{equation}
\label{eq:lind:core_eq}
    \dot\sigma(t) = V(t)^*\, \hat{P}_{\rho(t)}\bigl(\Lind(\rho(t))\bigr)\, V(t).
\end{equation}
Since $\hat{P}_\rho(X) = PXP + \text{off-diagonal} - \frac{\Tr(PXP)}{r}P$
and $P = VV^*$, this simplifies to
\begin{equation}
\label{eq:lind:core_explicit}
    \dot\sigma(t)
    = V^* \Lind(\rho) V
    - \frac{\Tr(V^* \Lind(\rho) V)}{r}\, I_{\CC^r}
    = V^* \Lind(\rho) V,
\end{equation}
where the last equality uses $\Tr(\Lind(\rho)) = 0$ (trace preservation
of the full Lindblad semigroup), which gives $\Tr(V^*\Lind(\rho)V) =
\Tr(VV^*\Lind(\rho)) = \Tr(P\Lind(\rho)) = 0$ when the $(S^\perp,S^\perp)$
block of $\Lind(\rho)$ vanishes on the tangent space.

\medskip
\noindent\textbf{Equation for the frame $V(t)$.}
Applying $(I - VV^*)$ on the left and $V\sigma^{-1}$ on the right to
$\dot\rho = \hat{P}_\rho(\Lind(\rho))$:
\begin{equation}
\label{eq:lind:frame_eq}
    \dot{V}(t)
    = (I - V(t)V(t)^*)\,
    \hat{P}_{\rho(t)}\bigl(\Lind(\rho(t))\bigr)\,
    V(t)\, \sigma(t)^{-1}.
\end{equation}
This is well-defined since $\sigma(t) \in \Dpos{\CC^r}$ is invertible
(positive definite) along the trajectory.

\begin{remark}[Equations~\eqref{eq:lind:core_explicit}--\eqref{eq:lind:frame_eq}
as a special case of the MCTDH system]
\label{rem:lind:mctdh_special}
Equations~\eqref{eq:lind:core_explicit}--\eqref{eq:lind:frame_eq} are
the quantum analogue of the MCTDH kinematic
system~\eqref{eq:DF:MCTDH_core}--\eqref{eq:DF:MCTDH_basis} of
Chapter~\ref{chap:DF}. Precisely:
\begin{itemize}
\item $\sigma(t) \in \Dpos{\CC^r}$ plays the role of the core tensor
    $C^{(D)}(t)$: the finite-dimensional degrees of freedom within the
    fixed support.
\item $V(t) \in \Stiefel{r}{H}$ plays the role of the basis
    $\{u^{(\alpha)}_{i_\alpha}(t)\}$ of $\Umin{\alpha}{\bv_r(t)}$:
    the isometric frame of the support subspace.
\item The gauge condition $V^*\dot{V} = 0$ is the quantum analogue of
    the tangential constraint~\eqref{eq:DF:tangential}.
\item The mean-field force $F^{(\alpha)}$ of~\eqref{eq:DF:mean_field_hilbert}
    becomes here $V^*\Lind(\rho)V$ in the core equation and
    $(I-VV^*)\Lind(\rho)V\sigma^{-1}$ in the frame equation.
\end{itemize}
The key difference is the additional constraint $\Tr(\dot\sigma) = 0$
(trace zero on the core), which has no analogue in the pure-state
MCTDH case.
\end{remark}

\subsection*{The closed kinematic system}

Combining~\eqref{eq:lind:core_explicit} and~\eqref{eq:lind:frame_eq},
the Dirac--Frenkel projected Lindblad dynamics in Stiefel-core
coordinates reads:
\begin{subequations}
\label{eq:lind:kinematic}
\begin{align}
    \dot\sigma(t)
    &= V(t)^*\, \Lind\bigl(V(t)\sigma(t)V(t)^*\bigr)\, V(t),
    \label{eq:lind:kinematic_core}\\[4pt]
    \dot{V}(t)
    &= \bigl(I - V(t)V(t)^*\bigr)\,
    \Lind\bigl(V(t)\sigma(t)V(t)^*\bigr)\,
    V(t)\, \sigma(t)^{-1},
    \label{eq:lind:kinematic_frame}
\end{align}
\end{subequations}
with initial conditions $V(0) = V_0 \in \Stiefel{r}{H}$ and
$\sigma(0) = V_0^* \rho_0 V_0 \in \Dpos{\CC^r}$.

\begin{proposition}[Well-posedness of the kinematic system]
\label{prop:lind:kinematic_wellposed}
Under Assumption~\ref{ass:lind:regular}, the system~\eqref{eq:lind:kinematic}
has a unique local solution $(V(t), \sigma(t))$ with
$V(t) \in \Stiefel{r}{H}$ and $\sigma(t) \in \Dpos{\CC^r}$ for
$t \in [0, T)$. The reconstructed density operator
$\rho(t) = V(t)\sigma(t)V(t)^*$ satisfies~\eqref{eq:lind:df}.
\end{proposition}

\begin{proof}
In the chart $\Phi_{S_0}$ centred at $\rho_0$, the system~\eqref{eq:lind:kinematic}
is a locally Lipschitz ODE on
$\cL(\CC^r, S_0^\perp) \times \cTsa{1}(\CC^r)$
(a Banach space, finite-dimensional in the core direction and infinite-
dimensional in the frame direction). Local existence and uniqueness
follow from Theorem~\ref{thm:app:picard}. Positivity and trace of $\sigma(t)$
are preserved as in Proposition~\ref{prop:lind:preservation}.
\end{proof}

\begin{remark}[Complexity and numerical implementation]
\label{rem:lind:complexity}
In a finite-dimensional truncation $H \cong \CC^d$ (dimension $d$,
rank $r \ll d$), the kinematic system~\eqref{eq:lind:kinematic}
requires storing only $V(t) \in \CC^{d \times r}$ and $\sigma(t) \in
\CC^{r \times r}$ at cost $O(dr + r^2)$, compared to $O(d^2)$ for the
full state $\rho(t)$. Each time step of~\eqref{eq:lind:kinematic}
costs $O(|J|dr^2 + dr^2 + r^3)$, where $|J|$ is the number of Lindblad
operators. The dominant cost is the evaluation of $\Lind(\rho)$
projected into the frame coordinates, which requires computing
$L_j V(t) \in \CC^{d \times r}$ for each $j$. This is a factor
$d/r$ reduction compared to the full simulation.
\end{remark}

%------------------------------------------------------------
\section{Historical notes}
\label{sec:lind:history}
%------------------------------------------------------------

\paragraph{The Lindblad equation.}
The GKSL master equation was derived independently by
Gorini, Kossakowski, and Sudarshan~(1976) and by Lindblad~(1976) as
the most general Markovian, completely positive, trace-preserving
quantum master equation. It is the quantum analogue of the Fokker--Planck
equation for classical Markov processes. The bounded-generator case
(Assumption~\ref{ass:lind:regular}) is completely understood;
the unbounded case remains an active research area.

\paragraph{Dirac--Frenkel for open systems.}
The application of the Dirac--Frenkel variational principle to the
Lindblad equation — projecting onto a finite-rank manifold — appears
in the numerical literature as the time-dependent variational principle
(TDVP) for density operators, used in the simulation of open quantum
spin chains via matrix product operators (MPO). The mathematical
foundations, including the well-posedness theorem
(Theorem~\ref{thm:lind:wellposed}) and the explicit bounded projector
(Proposition~\ref{prop:qs:projector}), are established in the companion
paper~\cite{AcevedoFalco2026}.

\paragraph{Low-rank methods for Lindblad equations.}
Numerical low-rank methods for the Lindblad equation were developed
by Le Bris, Rouchon, and Roussel~\cite{LeBrisRouchon2013,LeBrisRouchonRoussel2015}
in the Hilbert--Schmidt setting, and by Gravina and Savona~\cite{GravinaSavona2024}
using adaptive rank strategies. The CPTP-preserving approach via Kraus
parametrisations was developed by Appelö and Cheng~\cite{AppeloChengo2024}.
The present chapter provides a unified geometric framework for all
these methods.

\paragraph{Connection to tensor network methods.}
For multipartite open systems $H = \bigotimes_{j=1}^d H_j$, the density
operator $\rho(t)$ is a positive trace-class operator in
$\cTsa{1}(\bigotimes_j H_j)$. Approximating $\rho(t)$ by a matrix
product operator (MPO) of bond dimension $r$ is a special case of
the tree-based tensor format of Chapters~\ref{chap:topTBT}--\ref{chap:TB_manifold},
applied to the operator space $\cTsa{1}(\bigotimes_j H_j) \cong
\bigotimes_j \cL(H_j)$. The Dirac--Frenkel principle on this space
gives a rigorous foundation for TDVP-MPO algorithms for open systems.
This connection is explored in Open Problem~\ref{prob:lindblad_tensor}.

%============================================================
%  ch11_open.tex
%  Chapter 11: Open Problems and Perspectives
%
%  Tensor Formats in Banach Spaces
%  Antonio Falcó, CEU Universities
%
%  Estimated length: ~20 pp.
%============================================================

\chapter{Open Problems and Perspectives}
\label{chap:open}

The results established in this monograph open several directions for
further research. This chapter collects the most natural open problems,
organised by theme, together with a discussion of connections to other
active areas of mathematics and its applications.

%------------------------------------------------------------
\section{Global topology of tensor manifolds}
\label{sec:open:topology}
%------------------------------------------------------------

\subsection*{Connectedness of tree-based manifolds}

We showed in Proposition~\ref{prop:TBT_connected} that each Tucker
manifold $\Tucker{\fr}{\VD}$ is connected. For tree-based manifolds
with $\dep{\TD} \geq 2$, this question is open in the
infinite-dimensional setting.

\begin{problem}
\label{prob:connected}
Let $\TD$ be a dimension partition tree with $\dep{\TD} \geq 2$ and
let $\fr \in \ADTD$ with $\FTv$ a proper set. Is $\FTv$ connected?
\end{problem}

In finite dimensions, connectedness has been established for the HT
format~\cite{UschmajewVandereycken2013} and the TT format~\cite{HoltzRohwedderSchneider2012},
but the methods used there (algebraic geometry of determinantal varieties)
do not directly extend to the Banach setting.

\subsection*{Higher homotopy groups}

Even in finite dimensions, the homotopy type of $\Tucker{\fr}{\VD}$
and $\FTv$ is not fully understood. Since $\Tucker{\fr}{\VD}$ is
a fibre bundle over the product of Grassmannians
$\prod_\alpha \Grass{r_\alpha}{V_\alpha}$ (Theorem~\ref{thm:Tucker_Banach_Manifold}),
the long exact sequence of homotopy groups relates
$\pi_n(\Tucker{\fr}{\VD})$ to $\pi_n(\prod_\alpha \Grass{r_\alpha}{\RR^{n_\alpha}})$
and $\pi_n(\RR^{\times r_\alpha}_*)$.

\begin{problem}
\label{prob:homotopy}
Compute the fundamental group $\pi_1(\Tucker{\fr}{\VD})$ and
$\pi_1(\FTv)$ in the finite-dimensional case. Are these manifolds
simply connected?
\end{problem}

\subsection*{Rank stratification}

The decomposition $\VD = \bigsqcup_\fr \Tucker{\fr}{\VD}$ is a
stratification of the algebraic tensor space. The boundary
$\overline{\Tucker{\fr}{\VD}} \setminus \Tucker{\fr}{\VD}$ consists of
tensors with strictly smaller rank, and the closure
$\overline{\Tucker{\fr}{\VD}}^{\|\cdot\|}$ in the Banach topology is the
set $\Tucker{\leq\fr}{\VD}$ of tensors of bounded rank
(Theorem~\ref{thm:weak_closed}).

\begin{problem}
\label{prob:stratification}
Describe the local geometry of the boundary strata: what is the tangent
cone to $\Tucker{\leq\fr}{\VD}$ at a boundary point
$\bv \in \Tucker{\leq\fr}{\VD} \setminus \Tucker{\fr}{\VD}$?
Is the stratification Whitney regular?
\end{problem}

%------------------------------------------------------------
\section{Riemannian and Finsler geometry on tensor manifolds}
\label{sec:open:riemannian}
%------------------------------------------------------------

In the Hilbert case, the Tucker manifold inherits a Riemannian metric
from the ambient Hilbert tensor space. Riemannian gradient methods
on $\Tucker{\fr}{\VD}$ (ALS, Riemannian gradient, retraction-based
methods) have been developed by Uschmajew and Vandereycken~\cite{UschmajewVandereycken2020}.
In the Banach setting, a natural replacement is a Finsler metric (a
continuously varying norm on each tangent space).

\begin{problem}
\label{prob:finsler}
Equip $\Tucker{\fr}{\VD}$ with the Finsler metric induced by
$\|\cdot\|_D$ (via the isomorphism
$T_\bv\Tucker{\fr}{\VD} \cong \ZD{\bv}$). Study the geodesics,
curvature, and exponential map of this Finsler manifold. In particular:
does the Dirac--Frenkel flow~\eqref{eq:DF_tucker} preserve the
Finsler structure in any sense?
\end{problem}

%------------------------------------------------------------
\section{Algorithms and convergence analysis}
\label{sec:open:algorithms}
%------------------------------------------------------------

\subsection*{Projector-splitting integrators on tree-based manifolds}

For the TT format, Lubich, Oseledets, and
Vandereycken~\cite{LubichOseledetsVandereycken2015} introduced the
\emph{projector-splitting integrator}, which integrates the
Dirac--Frenkel ODE by splitting the projection onto $\ZD{\bv_r(t)}$
into a sequence of rank-one updates. For general tree-based formats,
the corresponding integrator was introduced in~\cite{CerutiLubichWalach2020}.

The convergence analysis of these algorithms in infinite-dimensional
Banach spaces (as opposed to Hilbert spaces) requires understanding
the stability of the Dirac--Frenkel flow on $\FTv$ and the
approximation properties of the splitting.

\begin{problem}
\label{prob:integrator}
Extend the convergence analysis of the projector-splitting integrator
to the setting of Banach tensor spaces. In particular, establish
error bounds that depend explicitly on the geometry of the tree-based
manifold $\FTv$ (curvature, injectivity radius, condition number of
the rank map $\varrho_\fr$).
\end{problem}

\subsection*{Alternating least squares in Banach spaces}

The \emph{alternating least squares} (ALS) algorithm and its variants
(DMRG, HOSVD) minimise a functional over $\Tucker{\leq\fr}{\VD}$ by
alternating optimisation over one factor at a time. In Hilbert spaces,
convergence of ALS is well understood~\cite{UschmajewVandereycken2020}.
In Banach spaces, convergence requires different tools.

\begin{problem}
\label{prob:ALS}
Prove convergence of ALS for minimisation over $\Tucker{\leq\fr}{\VD}$
in a reflexive Banach tensor space $\VDnorm$. What role does the
geometry of the Tucker manifold (specifically, the structure of the
tangent spaces $\ZD{\bv}$) play in the convergence rate?
\end{problem}

%------------------------------------------------------------
\section{Beyond Banach spaces}
\label{sec:open:beyond}
%------------------------------------------------------------

\subsection*{Quasi-Banach and Fréchet spaces}

The component spaces $V_\alpha$ in this monograph are normed (Banach or
at least normed). Many function spaces of interest — $L^p$ for
$0 < p < 1$, or spaces of smooth functions — are quasi-Banach or Fréchet
but not Banach. The theory of minimal subspaces and Tucker manifolds
should extend to these settings, but the proofs require substantial
modifications.

\begin{problem}
\label{prob:quasi-banach}
Develop the theory of minimal subspaces and tensor manifolds for
quasi-Banach spaces $V_\alpha$ with $0 < p < 1$. In particular,
does the Tucker manifold $\Tucker{\fr}{\VD}$ retain its smooth
structure when the $V_\alpha$ are $L^p$ spaces for $p < 1$?
\end{problem}

\subsection*{Infinite-dimensional index sets}

All results in this monograph assume $D = \{1, \ldots, d\}$ is finite.
For problems in infinite dimensions — e.g., infinite particle quantum
systems or PDEs in infinitely many variables — one would want $D$ to
be infinite. The algebraic tensor product over an infinite index set
requires projective limits or nuclear spaces.

\begin{problem}
\label{prob:infinite_D}
Extend the tree-based tensor framework to infinite index sets
$D = \NN$ (or more general sets). What is the natural analogue of
a dimension partition tree, and what conditions on the norms guarantee
the existence of minimal subspaces and best approximations?
\end{problem}

%------------------------------------------------------------
\section{Connections with quantum information}
\label{sec:open:quantum}
%------------------------------------------------------------

\subsection*{Tensor networks and entanglement}

In quantum information, a \emph{tensor network state} on a lattice
is a tensor $\bv \in \bigotimes_{j \in D} \CC^{d_j}$ whose
transfer tensors $\{C^{(\alpha)}\}$ are specified by a graph (not
necessarily a tree). The tree-based format corresponds to
\emph{tree tensor networks} (TTN), which include MPS/TT as the
linear-chain case and MERA (multi-scale entanglement renormalisation
ansatz) as a more general tree.

The \emph{entanglement entropy} of a bipartition $D = \alpha \cup \alpha^c$
is $S(\alpha) = -\mathrm{tr}(\rho_\alpha \log \rho_\alpha)$, where
$\rho_\alpha$ is the reduced density matrix of the state on $\alpha$.
The \emph{area law} for ground states of local Hamiltonians states that
$S(\alpha)$ grows at most as $|\partial \alpha|$ (the boundary of
$\alpha$), not as $|\alpha|$ (the volume). This motivates the use of
tree-based formats: the rank constraints in $\FTv$ directly control
the entanglement structure.

\begin{problem}
\label{prob:quantum}
Formulate the Dirac--Frenkel variational principle for tree tensor
networks in the language of $C^*$-algebras and operator systems.
What is the correct generalisation of the Tucker manifold in the
non-commutative setting?
\end{problem}

%------------------------------------------------------------
\section{Learning on tensor manifolds}
\label{sec:open:learning}
%------------------------------------------------------------

\subsection*{Statistical learning theory}

In machine learning, tensor network states arise naturally as
\emph{function approximation schemes} for high-dimensional data.
The generalisation capacity of these models depends on the geometry
of the hypothesis class $\FTleq{\fr}{\VD}{\TD}$.

The \emph{Rademacher complexity} of $\FTleq{\fr}{\VD}{\TD}$ controls
the generalisation error via the Dudley entropy integral. For finite
$D$ and fixed ranks, this can be bounded using the Weyl tube formula
once the curvature of the manifold $\FTv$ is understood.

\begin{problem}
\label{prob:learning}
Compute the Rademacher complexity and metric entropy of
$\FTleq{\fr}{\VD}{\TD}$ in terms of the tree-based rank $\fr$,
the depth $\dep{\TD}$, and the dimensions $\dim V_j$. How does the
choice of tree $\TD$ affect the approximation-complexity tradeoff?
\end{problem}

\subsection*{Optimal tree selection}

Given a function $f: \prod_j I_j \to \RR$ (a target function for
approximation), which tree $\TD$ gives the best approximation of $f$
by a tree-based tensor of a given rank? This question connects the
geometry of tensor manifolds to the combinatorics of tree structures.

\begin{problem}
\label{prob:tree_selection}
Given $f \in \VDnorm$ and a bound $r$ on the tree-based rank,
find the dimension partition tree $\TD^*$ that minimises the
approximation error $\inf_{\bv \in \FTleq{\fr}{\VD}{\TD^*}} \|f - \bv\|_D$.
Is this optimisation problem tractable?
\end{problem}

        % Open Problems — final chapter

%------------------------------------------------------------
%  APPENDICES
%------------------------------------------------------------
\appendix

%============================================================
%  app_manifolds.tex — Banach Manifolds: Reference Summary
%  Tensor Formats in Banach Spaces
%  Antonio Falcó, CEU Universities
%============================================================

\chapter{Banach Manifolds: Reference Summary}
\label{app:manifolds}

This appendix is a reference guide to the differential-geometric notions
used in the monograph. The full exposition --- with complete proofs and
examples --- is in Chapter~\ref{chap:banach}. This appendix records
the precise statements for quick lookup, and adds the material on
connections and curvature that is not needed in the main proofs but
is relevant to the open problems of Chapter~\ref{chap:open}.

Standard references: \cite{Lang1999} and~\cite{Upmeier1985}.

%------------------------------------------------------------
\section{Quick-reference index to Chapter~\ref{chap:banach}}
\label{app:manifolds:index}
%------------------------------------------------------------

\begin{center}
\renewcommand{\arraystretch}{1.35}
\begin{tabular}{@{}p{7cm}l@{}}
\toprule
\textbf{Concept} & \textbf{Location} \\
\midrule
Banach manifold, atlas, AT1--AT3
    & Def.~\ref{def:banach_manifold}, \S\ref{sec:banach:manifolds} \\
Tangent space and tangent map
    & Def.~\ref{def:tangent_space}, \S\ref{sec:banach:manifolds} \\
Immersion, split image condition
    & Def.~\ref{def:immersion}, \S\ref{sec:banach:manifolds} \\
Split (complemented) subspaces
    & Prop.~\ref{prop:split_equiv}, \S\ref{sec:banach:complemented} \\
Finite-dimensional subspaces are split
    & Prop.~\ref{prop:finitedim_split}, \S\ref{sec:banach:complemented} \\
Nilpotent algebra $\cL_{(U,W)}(X)$, $\exp(L)=\Id+L$
    & Prop.~\ref{prop:nilpotent_algebra}, \S\ref{sec:banach:complemented} \\
Grassmann--Banach manifold, chart map $\Psi_{U\oplus W}$
    & Thm.~\ref{thm:grassmann_manifold}, \S\ref{sec:banach:grassmann} \\
Tangent space $T_U\Grass{r}{X}\cong\cL(U,W)$
    & Prop.~\ref{prop:grassmann_tangent}, \S\ref{sec:banach:grassmann} \\
Continuous multilinear maps are $\mathcal{C}^\infty$
    & Prop.~\ref{prop:multilinear_smooth}, \S\ref{sec:banach:multilinear} \\
$\mathcal{C}^1 \Leftrightarrow$ analytic (complex)
    & Prop.~\ref{prop:analytic_c1}, \S\ref{sec:banach:multilinear} \\
Tucker fibre bundle, trivialisations, structure group $G_\fr$
    & \S\ref{sec:tucker_manifold:bundle} \\
Quotient construction fails in $\infty$-dim.\
    & \S\ref{sec:tucker_manifold:bundle}, \S\ref{sec:tucker_manifold:history} \\
\bottomrule
\end{tabular}
\end{center}

%------------------------------------------------------------
\section{ODEs on Banach manifolds}
\label{app:manifolds:ode}
%------------------------------------------------------------

The following theorem is used in Chapter~\ref{chap:DF}
(proof of Theorem~\ref{thm:DF_tucker}) but does not appear
explicitly in Chapter~\ref{chap:banach}.

A \emph{smooth vector field} on a $\mathcal{C}^\infty$-Banach manifold
$\mathbb{M}$ is a $\mathcal{C}^\infty$-map
$F: \mathbb{M} \to T\mathbb{M}$ with $F(m) \in T_m\mathbb{M}$.
An \emph{integral curve} of $F$ through $m_0$ is a $\mathcal{C}^1$-curve
$\gamma: (-\varepsilon,\varepsilon) \to \mathbb{M}$ with
$\gamma(0)=m_0$ and $\dot\gamma(t)=F(\gamma(t))$.

\begin{theorem}[Existence and uniqueness; {\cite[Theorem~4.1.5]{Lang1999}}]
\label{thm:app:ode}
Let $\mathbb{M}$ be a $\mathcal{C}^\infty$-Banach manifold and $F$ a
$\mathcal{C}^\infty$ vector field. For every $m_0\in\mathbb{M}$ there
exist $\varepsilon>0$ and a unique integral curve
$\gamma:(-\varepsilon,\varepsilon)\to\mathbb{M}$ through $m_0$.
The flow map $(t,m_0)\mapsto\gamma(t)$ is $\mathcal{C}^\infty$.
\end{theorem}

\begin{remark}
In local coordinates the theorem reduces to Picard--Lindelöf in Banach
spaces (Theorem~\ref{thm:app:picard}). The $\mathcal{C}^\infty$ flow
follows from smooth dependence on initial data; see~\cite[Ch.~4]{Lang1999}.
\end{remark}

%------------------------------------------------------------
\section{The tangent bundle}
\label{app:manifolds:tangentbundle}
%------------------------------------------------------------

The \emph{tangent bundle} of a $\mathcal{C}^p$-manifold $\mathbb{M}$
($p\geq 1$) is
\begin{equation}
\label{eq:app:tangent_bundle}
    T\mathbb{M} \coloneqq \bigsqcup_{m\in\mathbb{M}} T_m\mathbb{M},
\end{equation}
with projection $\pi_{T\mathbb{M}}:(m,v)\mapsto m$ and zero section
$\sigma_0:m\mapsto(m,0)$. Each chart $(M_i,u_i)$ on $\mathbb{M}$
induces a chart $Tu_i:\pi^{-1}(M_i)\to u_i(M_i)\times X_i$,
making $T\mathbb{M}$ a $\mathcal{C}^{p-1}$-Banach manifold modelled
on $X_i\times X_i$.

\begin{proposition}[{\cite[Chapter~3]{Lang1999}}]
\label{prop:app:tangent_bundle}
$T\mathbb{M}$ is a $\mathcal{C}^{p-1}$-Banach manifold; $\pi_{T\mathbb{M}}$
is a $\mathcal{C}^{p-1}$-submersion; $\sigma_0$ is a
$\mathcal{C}^{p-1}$-embedding. For $f:\mathbb{M}\to\mathbb{N}$
of class $\mathcal{C}^p$, the tangent map $Tf:T\mathbb{M}\to T\mathbb{N}$
satisfies $T(g\circ f)=Tg\circ Tf$.
\end{proposition}

\begin{remark}[Tucker manifold]
\label{rem:app:tangent_TM}
The fibre $T_\bv\Tucker{\fr}{\VD}\cong\Emodel{D}$
(Theorem~\ref{thm:Tucker_Banach_Manifold}(i)); the tangent map of
$\mathfrak{i}:\Tucker{\fr}{\VD}\hookrightarrow\VDnorm$ has image
$\ZD{\bv}$ (Proposition~\ref{prop:tangent_space}); and the immersion
theorem (Theorem~\ref{thm:tucker_immersion}) says this map has
closed split image.
\end{remark}

%------------------------------------------------------------
\section{Connections, parallel transport, and curvature}
\label{app:manifolds:connections}
%------------------------------------------------------------

This section is included for completeness and for
Chapter~\ref{chap:open}; none of this material is used in the
proofs of the main text.

\subsection*{Connections on vector bundles}

A \emph{(linear) connection} on $\pi:E\to\mathbb{M}$ is a bilinear
map $\nabla:\mathfrak{X}(\mathbb{M})\times\Gamma(E)\to\Gamma(E)$
satisfying $\nabla_{fX}s=f\nabla_Xs$ and
$\nabla_X(fs)=(Xf)s+f\nabla_Xs$ for all $f\in\mathcal{C}^\infty(\mathbb{M})$.
In a local trivialisation, $\nabla_Xs=X(s)+\omega(X)(s)$ where $\omega$
is the $\cL(F)$-valued connection $1$-form.

\emph{Parallel transport} along $\gamma:[0,1]\to\mathbb{M}$ gives
isomorphisms $P_t:E_{\gamma(0)}\to E_{\gamma(t)}$ by solving
$\nabla_{\dot\gamma}s=0$, $s(0)=e_0$.

\subsection*{Levi-Civita connection and curvature}

If $\mathbb{M}$ is Riemannian (each $T_m\mathbb{M}$ has a smoothly
varying inner product), the unique torsion-free metric connection is
the \emph{Levi-Civita connection}; its curvature tensor is
\begin{equation}
\label{eq:app:curvature_tensor}
    R(X,Y)Z
    = \nabla_X\nabla_YZ - \nabla_Y\nabla_XZ - \nabla_{[X,Y]}Z.
\end{equation}
The \emph{sectional curvature} at $m \in \mathbb{M}$ in the direction
of a two-plane $\Pi = \spann\{X,Y\} \subset T_m\mathbb{M}$ is
\begin{equation}
\label{eq:app:sectional_curvature}
    K(\Pi) = \frac{\langle R(X,Y)Y, X\rangle}
    {\|X\|^2\|Y\|^2 - \langle X,Y\rangle^2}.
\end{equation}

\begin{proposition}[Jacobi fields and curvature,
{\cite[Chapter~10]{Lang1999}}]
\label{prop:app:jacobi}
Let $\gamma$ be a geodesic in a Riemannian Banach manifold $(\mathbb{M},g)$.
A \emph{Jacobi field} $J(t)$ along $\gamma$ satisfies the Jacobi
equation
\begin{equation}
    \nabla_{\dot\gamma}\nabla_{\dot\gamma} J
    + R(\dot\gamma, J)\dot\gamma = 0.
\end{equation}
If $K \leq \kappa$ along $\gamma$, the spread of neighbouring geodesics
is bounded below by the corresponding space-form comparison: for two
geodesics starting at $\gamma(0)$ with angle $\theta$, the distance
after time $t$ is $\geq t\theta \cdot s_\kappa(t)/t$ where $s_\kappa$
is the sine function in the $\kappa$-space-form.
\end{proposition}

\begin{remark}[Sectional curvature and error bounds for Dirac--Frenkel]
\label{rem:app:curvature_DF}
For the Tucker manifold $\Tucker{\fr}{\VD}$ with $\VDnorm =
\bigotimes_j L^2(I_j)$ (Hilbert case), the sectional curvature of the
inherited Riemannian metric from $\VDnorm$ is non-positive in directions
tangent to the orbit of $\Tucker{\fr}{\VD}$ under the action of
$\prod_\alpha U(r_\alpha)$ (unitary group on the minimal subspaces).
Bounded curvature would imply, via Jacobi field estimates, that the
Dirac--Frenkel projected solution cannot deviate from the full solution
faster than a rate controlled by the curvature; this is the geometric
approach to the error bounds of Open Problem~\ref{prob:lind:error_bounds}.
For the density operator manifold $\Drank{r}{H}$, the curvature is
studied in Open Problem~\ref{prob:curvature_entanglement}.
\end{remark}

\subsection*{Holonomy and homotopy groups}

The \emph{holonomy group} $\mathrm{Hol}(m)$ at $m \in \mathbb{M}$ is the
group of parallel transport isomorphisms $P_\gamma: T_m\mathbb{M} \to
T_m\mathbb{M}$ along all loops $\gamma$ based at $m$. By the
Ambrose--Singer theorem, $\mathrm{Hol}(m)$ is a Lie group whose Lie
algebra is generated by the curvature operators $\{R(X,Y) : X, Y \in
T_m\mathbb{M}\}$~\cite{AbrahamMarsdenRatiu1988}.

\begin{proposition}[Holonomy of the Grassmannian,
{\cite[Chapter~3]{MilnorStasheff1974}}]
\label{prop:app:holonomy_grassmann}
For $\Grass{r}{H}$ with $H$ a separable Hilbert space and the Riemannian
metric induced by $\langle \cdot, \cdot\rangle_H$, the holonomy group
at $U \in \Grass{r}{H}$ is $U(r)$ (acting on $T_U\Grass{r}{H} \cong
\cL(U, U^\perp)$ by left multiplication). In particular,
$\Grass{r}{H}$ is an irreducible symmetric space.
\end{proposition}

%------------------------------------------------------------
\section{The cotangent bundle and differential forms}
\label{app:manifolds:cotangent}
%------------------------------------------------------------

The \emph{cotangent bundle} of $\mathbb{M}$ is
\begin{equation}
    T^*\mathbb{M} \coloneqq \bigsqcup_{m \in \mathbb{M}} (T_m\mathbb{M})^*,
\end{equation}
where $(T_m\mathbb{M})^*$ is the Banach dual of the tangent space.
A \emph{$1$-form} on $\mathbb{M}$ is a section $\alpha: \mathbb{M} \to
T^*\mathbb{M}$; in local coordinates, $\alpha = \sum_i \alpha_i(x)\,dx^i$.
The \emph{exterior derivative} $d\alpha$ is a $2$-form defined by
$(d\alpha)(X,Y) = X(\alpha(Y)) - Y(\alpha(X)) - \alpha([X,Y])$.

\begin{remark}[Relevance to the Tucker and Lindblad manifolds]
\label{rem:app:cotangent_lind}
The cotangent bundle $T^*\Tucker{\fr}{\VD}$ is used implicitly when
defining the gradient of a functional $J: \Tucker{\fr}{\VD} \to \RR$
in the non-Hilbert Banach setting: the Fréchet derivative
$dJ(\bv): T_\bv\Tucker{\fr}{\VD} \to \RR$ is a cotangent vector,
and the \emph{Finsler gradient} $\nabla_F J(\bv) \in
T_\bv\Tucker{\fr}{\VD}$ is its pre-image under the duality mapping
$J: T_\bv\Tucker{\fr}{\VD} \to (T_\bv\Tucker{\fr}{\VD})^*$
(Proposition~\ref{prop:app:duality_mapping}). This is the starting
point for Riemannian and Finsler optimisation on tensor manifolds
(Open Problem~\ref{prob:supervised_learning_banach}).

For the density operator manifold $\Drank{r}{H} \subset \cTsa{1}(H)$,
the cotangent space at $\rho$ is a subspace of $(\cTsa{1}(H))^* \cong
\cT{\infty}(H)_{\mathrm{sa}}$ (bounded self-adjoint operators). The
duality pairing $\langle A, \rho\rangle = \Tr(A\rho)$ gives the natural
pairing between observables and states, and the gradient of a functional
$J(\rho) = \Tr(A\rho)$ (linear observable) is the projection of $A$
onto the tangent space of $\Drank{r}{H}$ via $\hat{P}_\rho$.
\end{remark}

%============================================================
%  app_functional.tex  —  Functional Analysis: Selected Results
%  TODO: Write content
%============================================================

\chapter{Functional Analysis: Selected Results}
\label{app:functional}

\section*{To be written}

This chapter is a placeholder. Content will be added in subsequent
revisions.


%============================================================
%  app_quantum.tex
%  Appendix E: Operators on Hilbert Spaces — Reference Summary
%
%  Tensor Formats in Banach Spaces
%  Antonio Falcó, CEU Universities
%
%  This appendix is a reference guide to the operator-theoretic
%  notions used in Part V. The full exposition with proofs is in
%  Chapters 12 and 13. This appendix contains:
%    - A quick-reference index to Chapters 12-13
%    - Self-contained statements of the foundational theorems
%      (spectral theorem, SVD, trace class)
%    - The tensor product identifications connecting operators
%      with the framework of Chapters 1-4
%    - Historical notes
%
%  Standard references: ReedSimon1972, BratteliRobinson1987, Ryan2002
%============================================================

\chapter{Operators on Hilbert Spaces: Reference Summary}
\label{app:quantum}

This appendix is a reference guide to the operator-theoretic vocabulary
used in Part~\ref{part:quantum}. The full mathematical development ---
with complete proofs and motivation --- is in
Chapters~\ref{chap:quantum_states} and~\ref{chap:lindblad}.
This appendix is intended for two purposes: as a \emph{quick reference}
for readers who are familiar with operator theory and want to locate a
specific result, and as a \emph{self-contained introduction} for
readers without prior background in quantum mechanics who wish to
read Part~\ref{part:quantum}.

Throughout this appendix, $H$ denotes a separable complex Hilbert space
of \emph{arbitrary} (possibly infinite) dimension.
Standard references: \cite{ReedSimon1972}, \cite{BratteliRobinson1987},
and~\cite{Ryan2002} for the trace-class tensor-product perspective.

%------------------------------------------------------------
\section{Quick-reference index to Part~\ref{part:quantum}}
\label{app:quantum:index}
%------------------------------------------------------------

\begin{center}
\renewcommand{\arraystretch}{1.35}
\begin{tabular}{@{}p{7.5cm}l@{}}
\toprule
\textbf{Concept} & \textbf{Location} \\
\midrule
Inner product, Riesz isomorphism $J: H \to H^*$
    & \S\ref{app:quantum:hilbert} \\
Dual $H^*$ vs.\ conjugate space $\overline{H}$
    & Remark~\ref{rem:app:dual_vs_conjugate} \\
Dirac notation (not used; comparison only)
    & Remark~\ref{rem:app:dirac} \\
Adjoint $A^*$, self-adjoint, positive, unitary operators
    & \S\ref{app:quantum:bounded} \\
Orthogonal projections and closed subspaces
    & Prop.~\ref{prop:split_equiv} \\
Rank-one operators $u \otimes \varphi \in \cL(H)$
    & \S\ref{app:quantum:bounded},
      Remark~\ref{rem:app:rank_one_tensor} \\
Compact operators, spectral theorem
    & \S\ref{app:quantum:compact},
      Thm.~\ref{thm:app:spectral} \\
Singular values, SVD
    & \S\ref{app:quantum:compact},
      Thm.~\ref{thm:app:svd} \\
Hilbert--Schmidt operators $\cT{2}(H)$,
    $\langle A,B\rangle_{\mathrm{HS}} = \Tr(A^*B)$
    & \S\ref{app:quantum:hs},
      Def.~\ref{def:app:hs} \\
$\cT{2}(H) \cong H \otimes_H H^*$ isometrically
    & Thm.~\ref{thm:app:hs_tensor},
      Remark~\ref{rem:app:hs_tensor} \\
Trace-class operators $\cT{1}(H)$, trace norm $\|\cdot\|_1$
    & \S\ref{app:quantum:trace},
      Def.~\ref{def:app:trace_class} \\
$\cT{1}(H) \cong H \otimes_\pi H^*$, \condINJ\ and $\|\cdot\|_1$
    & Remark~\ref{rem:app:trace_projective} \\
Partial trace $\TrS{\alpha^c}(\rho)$
    & \S\ref{app:quantum:trace} \\
Density operators $\Dstate{H}$, pure/mixed states
    & \S\ref{app:quantum:states},
      Def.~\ref{def:app:density} \\
$\supp(\rho_\alpha) = \Umin{\alpha}{v}$
    — the fundamental connection
    & Thm.~\ref{thm:qs:minimal_partial_trace}
      in Ch.~\ref{chap:quantum_states} \\
Rank stratification $\Dstate{H} = \bigsqcup_r \Drank{r}{H}$
    & Prop.~\ref{prop:qs:rank_tucker}
      in Ch.~\ref{chap:quantum_states} \\
Manifold structure of $\Drank{r}{H}$, tangent space
    & Thm.~\ref{thm:qs:manifold}
      in Ch.~\ref{chap:quantum_states} \\
Bounded projector $\hat{P}_\rho: \cTsa{1}(H) \to T_\rho\Drank{r}{H}$
    & Prop.~\ref{prop:qs:projector}
      in Ch.~\ref{chap:quantum_states} \\
Lindblad (GKSL) equation
    & \S\ref{sec:lind:equation}
      in Ch.~\ref{chap:lindblad} \\
Well-posedness of projected Lindblad dynamics
    & Thm.~\ref{thm:lind:wellposed}
      in Ch.~\ref{chap:lindblad} \\
Kinematic equations in Stiefel-core coordinates
    & \S\ref{sec:lind:kinematic}
      in Ch.~\ref{chap:lindblad} \\
\bottomrule
\end{tabular}
\end{center}

%------------------------------------------------------------
\section{Complex Hilbert spaces and the dual}
\label{app:quantum:hilbert}
%------------------------------------------------------------

A \emph{complex Hilbert space} $H$ carries an inner product
$\langle \cdot, \cdot\rangle: H \times H \to \CC$, conjugate-linear
in the first argument and linear in the second, positive definite, and
complete. The Riesz representation theorem gives the isometric
conjugate-linear isomorphism
\begin{equation}
\label{eq:app:riesz}
    J: H \longrightarrow H^*,
    \qquad J(v) = \langle v, \cdot\rangle,
\end{equation}
where $H^* = \cL(H, \CC)$ is the topological dual. This appendix works
throughout with $H^*$ (defined by the norm alone), not the conjugate
space $\overline{H}$ (defined by the linear structure). The two are
isomorphic as Banach spaces via $J$, but not canonically as complex
Banach spaces.

\begin{remark}[Dual vs.\ conjugate space]
\label{rem:app:dual_vs_conjugate}
In finite dimensions $H^* \cong H \cong \overline{H}$ as complex vector
spaces. In infinite dimensions, the Riesz map $J: H \to H^*$ is
conjugate-linear (not linear), making $H$ and $H^*$ non-isomorphic
as complex Banach spaces. The identification $\cT{1}(H) \cong H
\otimes_\pi H^*$ (Remark~\ref{rem:app:trace_projective}) requires
$H^*$ as the second factor to preserve the $\CC$-bilinearity of the
tensor product; using $\overline{H}$ instead introduces a
conjugate-linear identification that is inconsistent with the
Banach-space tensor framework of Chapter~\ref{chap:algtensor}.
\end{remark}

\begin{remark}[Dirac notation]
\label{rem:app:dirac}
In the physics literature, $v \in H$ is written $|v\rangle$ (a
\emph{ket}) and $J(v) \in H^*$ is written $\langle v|$ (a \emph{bra}),
so $\langle u, v\rangle_H = \langle u | v\rangle$. The rank-one operator
$u \otimes J(v)$ is written $|u\rangle\langle v|$. This monograph uses
exclusively the standard mathematical notation: $\langle u, v\rangle$
for the inner product, $J(v) \in H^*$ for the Riesz dual, and
$u \otimes \varphi \in \cL(H)$ for rank-one operators.
\end{remark}

%------------------------------------------------------------
\section{Bounded operators, adjoint, and special classes}
\label{app:quantum:bounded}
%------------------------------------------------------------

$\cL(H)$ with the operator norm is a C*-algebra. The adjoint
$A^* \in \cL(H)$ of $A$ is characterised by
$\langle A^*u, v\rangle = \langle u, Av\rangle$ for all $u, v \in H$,
with $\|A^*A\| = \|A\|^2$. Key classes:
self-adjoint ($A = A^*$), positive ($A \geq 0$: $\langle u, Au\rangle
\geq 0$), strictly positive ($A > 0$), unitary ($A^*A = AA^* = I$),
orthogonal projection ($A = A^* = A^2$). Every closed subspace $S
\subset H$ has an orthogonal projection $\PS{S}$ and
$H = S \oplus S^\perp$ (cf.\ Proposition~\ref{prop:finitedim_split},
which holds for \emph{all} closed subspaces in Hilbert spaces, not
just finite-dimensional ones).

For $u \in H$ and $\varphi \in H^*$, the rank-one operator
$u \otimes \varphi \in \cL(H)$ is defined by
$(u \otimes \varphi)(w) \coloneqq \varphi(w)\, u$.

\begin{remark}[Rank-one operators and the tensor product $H \otimes H^*$]
\label{rem:app:rank_one_tensor}
The finite-rank operators form the algebraic tensor product
$H \otimes_a H^*$ inside $\cL(H)$. The map $(u, \varphi) \mapsto
u \otimes \varphi$ is $\CC$-bilinear, making this consistent with the
tensor product framework of Chapter~\ref{chap:algtensor}. The
Hilbert--Schmidt completion gives $\cT{2}(H) \cong H \otimes_H H^*$
and the trace-class completion gives $\cT{1}(H) \cong H \otimes_\pi H^*$;
see Sections~\ref{app:quantum:hs}--\ref{app:quantum:trace}.
\end{remark}

%------------------------------------------------------------
\section{Compact operators, spectral theorem, and SVD}
\label{app:quantum:compact}
%------------------------------------------------------------

$A \in \cL(H)$ is \emph{compact} if it maps the unit ball to a
relatively compact set. The set $\cK(H)$ of compact operators is a
closed two-sided ideal, equal to the norm-closure of finite-rank
operators (in separable $H$).

\begin{theorem}[Spectral theorem for compact self-adjoint operators,
{\cite[Thm.\ VI.16]{ReedSimon1972}}]
\label{thm:app:spectral}
Let $A \in \cK(H)$ with $A = A^*$. Then there exist an orthonormal
sequence $\{u_k\}$ in $H$ and real $\lambda_k \to 0$ with
\begin{equation}
\label{eq:app:spectral}
    A = \sum_{k} \lambda_k\, u_k \otimes J(u_k),
\end{equation}
convergent in operator norm. The $\lambda_k$ are the non-zero
eigenvalues of $A$.
\end{theorem}

The \emph{singular values} of $A \in \cK(H)$ are the eigenvalues
$s_k(A) \geq 0$ of $|A| = (A^*A)^{1/2}$.

\begin{theorem}[Singular value decomposition]
\label{thm:app:svd}
For $A \in \cK(H)$ with non-zero singular values $(s_k(A))$, there
exist orthonormal sequences $\{u_k\}$, $\{v_k\}$ in $H$ with
\begin{equation}
\label{eq:app:svd}
    A = \sum_k s_k(A)\, v_k \otimes J(u_k),
\end{equation}
convergent in operator norm. For $d=2$ and $v \in H_1 \otimes_H H_2$,
this is the Schmidt decomposition; the Tucker rank $r_1(v) = r_2(v)$
equals the number of non-zero singular values
(Remark~\ref{rem:qs:schmidt} in Chapter~\ref{chap:quantum_states}).
\end{theorem}

%------------------------------------------------------------
\section{Hilbert--Schmidt operators}
\label{app:quantum:hs}
%------------------------------------------------------------

\begin{definition}
\label{def:app:hs}
$A \in \cL(H)$ is \emph{Hilbert--Schmidt} if
$\|A\|_2^2 \coloneqq \sum_k \|Ae_k\|^2 = \Tr(A^*A) < \infty$
for some (equivalently every) orthonormal basis $\{e_k\}$.
The space $\cT{2}(H)$ with $\langle A, B\rangle_{\mathrm{HS}}
\coloneqq \Tr(A^*B)$ is a Hilbert space, a two-sided ideal in
$\cL(H)$, with $\|A\| \leq \|A\|_2 = (\sum_k s_k(A)^2)^{1/2}$.
\end{definition}

\begin{theorem}[$\cT{2}(H) \cong H \otimes_H H^*$]
\label{thm:app:hs_tensor}
There is a canonical isometric isomorphism of Hilbert spaces
\begin{equation}
\label{eq:app:hs_iso}
    \cT{2}(H) \cong H \otimes_H H^*,
\end{equation}
under which the elementary tensor $u \otimes \varphi \in H \otimes_a H^*$
corresponds to the rank-one operator $w \mapsto \varphi(w)\, u$.
\end{theorem}

\begin{remark}[Why $H^*$ and not $\overline{H}$]
\label{rem:app:hs_tensor}
The identification~\eqref{eq:app:hs_iso} uses $H^*$ to preserve
$\CC$-bilinearity. In finite dimensions (or when $H = \overline{H}$ is
identified via $J$) one also sees $\cT{2}(H) \cong H \otimes_H \overline{H}$
in the literature, but this conflates a linear and a conjugate-linear
isomorphism. The present convention is consistent with the projective
tensor norm identification in the next section.
\end{remark}

%------------------------------------------------------------
\section{Trace-class operators}
\label{app:quantum:trace}
%------------------------------------------------------------

\begin{definition}
\label{def:app:trace_class}
$A \in \cL(H)$ is \emph{trace class} if
$\|A\|_1 \coloneqq \Tr|A| = \sum_k s_k(A) < \infty$.
The space $\cT{1}(H)$ with $\|\cdot\|_1$ is a Banach space (not Hilbert),
a two-sided ideal in $\cL(H)$, with $\cT{1}(H) \subset \cT{2}(H) \subset
\cK(H)$ and $\|A\| \leq \|A\|_2 \leq \|A\|_1$.
The \emph{trace} $\Tr(A) \coloneqq \sum_k \langle e_k, Ae_k\rangle$
is absolutely convergent and basis-independent.
The duality $\cT{1}(H)^* \cong \cL(H)$ holds via $\langle A, X\rangle
= \Tr(A^*X)$.
\end{definition}

\begin{remark}[$\cT{1}(H) \cong H \otimes_\pi H^*$ and \condINJ]
\label{rem:app:trace_projective}
The identification
\begin{equation}
\label{eq:app:tc_iso}
    \cT{1}(H) \cong H \otimes_\pi H^*
\end{equation}
holds isometrically~\cite{Grothendieck1953,Ryan2002}. On $H \otimes_a H^*$:
\begin{equation*}
    \|\cdot\|_\varepsilon \leq \|\cdot\|_H \leq \|\cdot\|_1 = \|\cdot\|_\pi,
\end{equation*}
where $\|\cdot\|_H$ is the Hilbert--Schmidt norm and $\|\cdot\|_\varepsilon$
the injective norm. Both $\|\cdot\|_H$ and $\|\cdot\|_\pi$ satisfy
condition~\condINJ\ of Definition~\ref{def:crossnorm}. The trace norm
$\frac{1}{2}\|\rho - \sigma\|_1$ equals the total variation distance
between quantum states $\rho$ and $\sigma$ — the maximum probability
of distinguishing them by any measurement — which is why it is the
natural norm for state estimation and compression.
\end{remark}

The \emph{partial trace} $\TrS{\alpha^c}(A) \in \cT{1}(H_\alpha)$,
for $A \in \cT{1}(H_1 \otimes_H H_2)$, is characterised by
$\Tr_{H_1}(B\, \TrS{2}(A)) = \Tr_H((B \otimes I_2)\, A)$
for all $B \in \cL(H_1)$.

%------------------------------------------------------------
\section{Density operators and quantum states}
\label{app:quantum:states}
%------------------------------------------------------------

\begin{definition}
\label{def:app:density}
A \emph{density operator} is $\rho \in \cTsa{1}(H) \coloneqq
\{A \in \cT{1}(H) : A = A^*\}$ with $\rho \geq 0$ and $\Tr(\rho)=1$.
The \emph{state space} is $\Dstate{H} \coloneqq \{\rho \in \cTsa{1}(H)
: \rho \geq 0,\, \Tr(\rho)=1\}$.
\end{definition}

$\Dstate{H}$ is convex and trace-norm closed, but not globally a
manifold (rank singularities, boundary phenomena).
A state is \emph{pure} if $\mathrm{rank}(\rho)=1$ (equivalently
$\rho = u \otimes J(u)$ for a unit vector $u \in H$) and \emph{mixed}
otherwise.

\medskip\noindent
The \emph{rank-$r$ stratum} $\Drank{r}{H} \coloneqq \{\rho \in \Dstate{H}:
\mathrm{rank}(\rho)=r\}$ is a smooth Banach manifold (Theorem~\ref{thm:qs:manifold}
in Chapter~\ref{chap:quantum_states}).
The stratification $\Dstate{H} = \bigsqcup_r \Drank{r}{H} \cup
\mathcal{D}_\infty(H)$ mirrors the Tucker rank decomposition
$\VD = \bigsqcup_\fr \Tucker{\fr}{\VD}$ of
Theorem~\ref{thm:tucker_decomposition_full}.

\medskip\noindent
The fundamental connection with minimal subspaces
(Theorem~\ref{thm:qs:minimal_partial_trace}) is:
for $v \in \bigotimes_{j=1}^d H_j$ with $\|v\|=1$,
\begin{equation}
\label{eq:app:min_supp}
    \supp\bigl(\TrS{\alpha^c}(v \otimes J(v))\bigr) = \Umin{\alpha}{v},
    \qquad
    \mathrm{rank}\bigl(\TrS{\alpha^c}(v \otimes J(v))\bigr) = r_\alpha(v).
\end{equation}
This is the cornerstone of Chapter~\ref{chap:quantum_states}.

%------------------------------------------------------------
\section{Historical notes}
\label{app:quantum:history}
%------------------------------------------------------------

\paragraph{Operator-theoretic foundations of quantum mechanics.}
Von Neumann~\cite{vonNeumann1932} established the functional-analytic
framework using density operators on Hilbert spaces. Schatten~\cite{Schatten1950}
introduced the Schatten $p$-classes $\cT{p}(H)$ with norm
$\|A\|_p = (\sum_k s_k(A)^p)^{1/p}$; the cases $p=1$ (trace class)
and $p=2$ (Hilbert--Schmidt) are central to Part~\ref{part:quantum}.

\paragraph{Tensor norms and the trace.}
The identification $\cT{1}(H) \cong H \otimes_\pi H^*$
(Remark~\ref{rem:app:trace_projective}) is due to
Grothendieck~\cite{Grothendieck1953}, connecting the trace norm to the
projective tensor norm of Chapter~\ref{chap:algtensor}. The
Hilbert--Schmidt identification $\cT{2}(H) \cong H \otimes_H H^*$
(Theorem~\ref{thm:app:hs_tensor}) is the non-commutative analogue of
the Fubini isomorphism $L^2(\Omega_1) \otimes_H L^2(\Omega_2) \cong
L^2(\Omega_1 \times \Omega_2)$ (cf.~\cite[Chapter~II]{ReedSimon1972}).

\paragraph{Minimal subspaces and partial traces.}
Equation~\eqref{eq:app:min_supp} makes precise a connection implicit in
the quantum information literature since Schmidt~(1907). Its systematic
treatment in the tensor-Banach-space framework is the contribution of
Falcó--Hackbusch~\cite{FH2012} and Acevedo--Falcó~\cite{AcevedoFalco2026}.

%============================================================
%  app_algorithms.tex
%  Appendix D: Computational Aspects and Algorithms
%
%  Tensor Formats in Banach Spaces
%  Antonio Falcó, CEU Universities
%============================================================

\chapter{Computational Aspects and Algorithms}
\label{app:algorithms}

The theory developed in this monograph --- minimal subspaces, Tucker
and tree-based manifolds, the Dirac--Frenkel variational principle ---
has direct computational consequences. This appendix describes the
main algorithms used in practice, explains their connection to the
mathematical structures of the main text, and identifies precisely
where the Banach-space generalisations of
Chapters~\ref{chap:minsubspaces}--\ref{chap:DF} modify or constrain
the algorithmic picture.

Each algorithm is stated with pseudocode, its cost is analysed, and
its relationship to the abstract theory is made explicit. The principal
references are~\cite{Hackbusch2019} (systematic treatment of tensor
numerics), \cite{Bachmayr2016} (survey of tensor network methods), and
\cite{UschmajewVandereycken2020} (Riemannian optimisation on tensor
manifolds).

%------------------------------------------------------------
\section{Storage complexity of tensor formats}
\label{app:alg:storage}
%------------------------------------------------------------

\subsection*{The curse of dimensionality, quantified}

Let $V_j = \RR^n$ for $j = 1, \ldots, d$. The full tensor space
$\VD = \bigotimes_{j=1}^d \RR^n$ has dimension $n^d$: intractable
for moderate $n$ and $d$. Tensor formats replace this by structured
low-dimensional parametrisations. The table below summarises storage
cost in terms of maximal rank $r$ and mode size $n$:

\medskip
\begin{center}
\renewcommand{\arraystretch}{1.4}
\begin{tabular}{@{}llll@{}}
\toprule
\textbf{Format} & \textbf{Parameters}
    & \textbf{Storage} & \textbf{Tree}\\
\midrule
Full & $n^d$ entries & $O(n^d)$ & ---\\
Tucker & core $r^d$ + $d$ factor matrices $n\!\times\! r$
    & $O(r^d + dnr)$ & $\TDtuck$\\
HT (balanced binary)
    & transfer tensors at $O(d)$ nodes, each $r\!\times\! r^2$
    & $O(dr^3 + dnr)$ & $\TDht$\\
TT (Tensor Train)
    & $d$ core matrices, each $r\!\times\! n\!\times\! r$
    & $O(dr^2n)$ & $\TDtt$\\
\bottomrule
\end{tabular}
\end{center}
\medskip

\begin{remark}[Exponential vs.\ polynomial scaling]
\label{rem:app:scaling}
The Tucker format reduces factor storage to $O(dnr)$ (linear in $d$)
but retains the core at $O(r^d)$. The HT and TT formats remove this
remaining exponential cost by imposing rank constraints at each level
of the tree, reducing core storage to $O(dr^3)$ and $O(dr^2n)$
respectively. This is precisely the content of the intersection
representation (Theorem~\ref{thm:intersection_representation}):
$\FTv$ imposes rank constraints at \emph{every} level of $\TD$,
not only at the leaves.
\end{remark}

\subsection*{Parameter count for tree-based formats}

For a general dimension partition tree $\TD$ and rank tuple $\fr$,
the total parameter count is:
\begin{equation}
\label{eq:app:param_count}
    P(\fr,\TD)
    = \sum_{j \in D} r_{\{j\}} n_j
    + \sum_{\substack{\alpha \in \TD \setminus (\Leaves \cup \{D\})}}
      r_\alpha \prod_{\beta \in \Sons{\alpha}} r_\beta
    + \prod_{\beta \in \Sons{D}} r_\beta,
\end{equation}
where the three terms count, respectively, the leaf factor matrices,
the internal transfer tensors, and the root tensor.

%------------------------------------------------------------
\section{The Higher-Order SVD}
\label{app:alg:hosvd}
%------------------------------------------------------------

\subsection*{Setup and goal}

The \emph{Higher-Order SVD} (HOSVD) of De Lathauwer, De Moor, and
Vandewalle~\cite{DeLathauwer2000} computes a Tucker approximation.
Given $\bv \in \VD = \bigotimes_{j=1}^d \RR^{n_j}$ and target ranks
$\fr = (r_j)_{j \in D}$, it produces:
\begin{enumerate}
\item[(i)] Factor matrices $U_j \in \RR^{n_j \times r_j}$ with
    orthonormal columns, approximating $\Umin{j}{\bv}$.
\item[(ii)] A core tensor
    $C^{(D)} \in \RR^{r_1 \times \cdots \times r_d}$ such that
    $\bv \approx \sum_{(i_j)} C^{(D)}_{(i_j)}
    \bigotimes_j u_{i_j}^{(j)}$.
\end{enumerate}

\subsection*{Algorithm}

\noindent\textbf{Algorithm D.1} (HOSVD).\\
\textit{Input:} tensor $\bv \in \RR^{n_1 \times \cdots \times n_d}$,
target ranks $(r_1,\ldots,r_d)$.\\
\textit{Output:} core tensor $C^{(D)}$, factor matrices
$U_1,\ldots,U_d$.
\begin{enumerate}
\item \textbf{Mode unfoldings.} For each $j = 1,\ldots,d$:
    \begin{enumerate}
    \item Compute the $j$-mode unfolding
        $M_j = \Mmat{\{j\}}(\bv) \in \RR^{n_j \times \prod_{k \neq j} n_k}$.
    \item Compute the truncated SVD: $M_j \approx U_j \Sigma_j V_j^\top$,
        retaining the $r_j$ largest singular values.
        Set $U_j \in \RR^{n_j \times r_j}$ (left singular vectors).
    \end{enumerate}
\item \textbf{Core compression.} Compute
    \begin{equation*}
        C^{(D)}_{i_1,\ldots,i_d}
        = \sum_{k_1,\ldots,k_d}
          (U_1)_{k_1 i_1} \cdots (U_d)_{k_d i_d}\;
          \bv_{k_1,\ldots,k_d},
    \end{equation*}
    i.e., project $\bv$ onto the subspaces $\spann(U_j)$ in each mode.
\item \textbf{Output.} Return $(C^{(D)}, U_1,\ldots,U_d)$.
\end{enumerate}

\subsection*{Cost and approximation quality}

The dominant cost is $d$ SVDs of unfolding matrices of size
$n_j \times \prod_{k \neq j} n_k$, giving arithmetic cost
$O(dn^{d+1})$, which still scales exponentially in $d$.
For tensors already in a low-rank format, recompression costs
$O(dr^{d+1})$ or $O(dr^3)$ depending on the tree structure.

\begin{theorem}[HOSVD quasi-optimality; {\cite[Theorem~4.2]{DeLathauwer2000}}]
\label{thm:app:hosvd}
Let $\bv \in \VD$ and let $\bw$ be the output of \textup{Algorithm~D.1}
with ranks $(r_j)_j$. Then
\begin{equation}
\label{eq:app:hosvd_bound}
    \|\bv - \bw\|_F
    \leq \sqrt{d}\;
    \min_{\bu \in \Tucker{\leq\fr}{\VD}} \|\bv - \bu\|_F.
\end{equation}
The HOSVD is not a best-approximation algorithm, but provides a
quasi-optimal approximation with factor $\sqrt{d}$.
\end{theorem}

\begin{remark}[Connection to minimal subspaces]
\label{rem:app:hosvd_min}
The factor matrix $U_j$ is the best rank-$r_j$ approximation to
the column space of the $j$-mode unfolding $M_j$, whose column space
is exactly $\Umin{j}{\bv}$ when $r_j = r_j(\bv)$ (the exact rank).
Algorithm~D.1 is therefore a computational procedure for approximating
the minimal subspaces of Chapter~\ref{chap:minsubspaces}.

In the Banach setting (non-Hilbert $V_j$), the SVD is replaced by a
best rank-$r_j$ approximation in the operator norm of
$\cL(V_j^*,\Vbigalpha^{[j]})$. Its existence is guaranteed by
Theorem~\ref{thm:best_approx_tucker}, but the quasi-optimality
factor $\sqrt{d}$ may worsen depending on the Banach--Mazur distance
of $\VDnorm$ from a Hilbert space.
\end{remark}

%------------------------------------------------------------
\section{Alternating Least Squares}
\label{app:alg:als}
%------------------------------------------------------------

\subsection*{The ALS problem and microiterations}

Given $\bu \in \VDnorm$ and rank tuple $\fr$, the best Tucker
approximation problem is
\begin{equation}
\label{eq:app:als_problem}
    \min_{\bw \in \Tucker{\leq\fr}{\VD}} \|\bu - \bw\|_D^2.
\end{equation}
\emph{Alternating Least Squares} (ALS) solves this by fixing all
factor matrices but one and optimising over the remaining one,
cycling over all modes.

\medskip
\noindent\textbf{Algorithm D.2} (ALS for Tucker format).\\
\textit{Input:} $\bu \in \VDnorm$, ranks $\fr$, initial factor
matrices $\{U_j^{(0)}\}_{j=1}^d$.\\
\textit{Repeat until convergence:}
\begin{enumerate}
\item For $j = 1,\ldots,d$:
    \begin{enumerate}
    \item \textbf{Fix} all $U_k$ for $k \neq j$ at their current values.
    \item \textbf{Form the reduced right-hand side:} compute the array
        \begin{equation*}
            \bigl(R_j\bigr)_{i_j, (i_k)_{k \neq j}}
            \coloneqq
            \sum_{(k_l)} \bu_{k_1,\ldots,k_d}
            \prod_{l \neq j} (U_l)_{k_l, i_l},
        \end{equation*}
        which is the tensor $\bu$ projected onto all modes except $j$.
        This gives $R_j \in \RR^{n_j \times \prod_{k \neq j} r_k}$.
    \item \textbf{Solve the mode-$j$ subproblem:}
        \begin{equation*}
            U_j^{\mathrm{new}}
            = \argmin_{U \in \RR^{n_j \times r_j}}
              \bigl\|R_j - U\, H_j\bigr\|_F^2,
        \end{equation*}
        where $H_j \in \RR^{r_j \times \prod_{k \neq j} r_k}$ is the
        Gram factor assembled from the other modes.
    \item \textbf{Orthonormalise:} $U_j \leftarrow U_j^{\mathrm{new}}$
        after a QR step.
    \end{enumerate}
\item \textbf{Update the core:}
    $C^{(D)}_{i_1,\ldots,i_d} \leftarrow
    \sum_{k_1,\ldots,k_d} \bu_{k_1,\ldots,k_d}
    \prod_j (U_j)_{k_j,i_j}$.
\end{enumerate}

\subsection*{Convergence}

\begin{theorem}[ALS convergence in Hilbert spaces;
    {\cite[Theorem~4.1]{UschmajewVandereycken2020}}]
\label{thm:app:als_convergence}
Let $\VDnorm$ be a Hilbert tensor space and let $(\bw^{(k)})$ be
the sequence generated by \textup{Algorithm~D.2}. Then:
\begin{enumerate}
\item[\textup{(i)}] The residuals $\|\bu - \bw^{(k)}\|$ are monotonically
    non-increasing.
\item[\textup{(ii)}] Every accumulation point $\bw^*$ satisfies
    the first-order optimality conditions in each mode direction:
    $\langle \bu - \bw^*, v \otimes \bw^*_{[j]}\rangle = 0$
    for all $v \in V_j$, $j = 1,\ldots,d$.
\item[\textup{(iii)}] Global convergence to a best approximation is
    not guaranteed; ALS may converge to a local minimum or saddle point.
\end{enumerate}
\end{theorem}

\begin{remark}[ALS in Banach spaces]
\label{rem:app:als_banach}
In a reflexive Banach tensor space $\VDnorm$, the mode-$j$ subproblem
in step~1(c) becomes: find the best rank-$r_j$ operator approximation
in the norm $\|\cdot\|_D$ projected onto mode~$j$. This is a convex
problem whose unique solution exists by
Theorem~\ref{thm:best_approx_tucker}. Monotone decrease of part~(i)
holds in full generality. The characterisation of accumulation points
in part~(ii) must be rephrased using the duality mapping $J$
(Appendix~\ref{app:functional}) in place of the inner product:
\begin{equation*}
    \langle v \otimes \bw^*_{[j]},\,
    J(\bu - \bw^*)\rangle = 0
    \quad \forall\, v \in V_j,\; j = 1,\ldots,d.
\end{equation*}
The convergence analysis in this Banach setting is Open
Problem~\ref{prob:ALS}.
\end{remark}

%------------------------------------------------------------
\section{The projector-splitting integrator}
\label{app:alg:splitting}
%------------------------------------------------------------

\subsection*{The Dirac--Frenkel ODE}

Consider the Dirac--Frenkel reduced model on $\Tucker{\fr}{\VD}$
(Chapter~\ref{chap:DF}):
\begin{equation}
\label{eq:app:df_reduced}
    \dot{\bv}_r(t)
    = \Pmetric{\bv_r(t)}\bigl(\bF(t,\bv_r(t))\bigr),
    \quad \bv_r(0) = \bv_0.
\end{equation}
Direct integration requires evaluating $\Pmetric{\bv_r(t)}$ at each
step. The \emph{projector-splitting integrator}
of~\cite{LubichOseledetsVandereycken2015} avoids this by decomposing
the projection into a sequence of simpler updates.

\subsection*{The TT case}

For the TT format ($\TD = \TDtt$, bond dimensions
$r_1,\ldots,r_{d-1}$), the tangent space projection splits as
\begin{equation}
\label{eq:app:splitting}
    \Pmetric{\bv_r}
    = \sum_{k=1}^{d} P_k
      - \sum_{k=1}^{d-1} P_{k,k+1},
\end{equation}
where $P_k$ projects onto the subspace with the $k$-th component free
and $P_{k,k+1}$ is the correction at the $(k,k{+}1)$ interface.

\medskip
\noindent\textbf{Algorithm D.3} (Projector-splitting, one step $h$).\\
\textit{Input:} $\bv_r^n$ in TT format, vector field $\bF$, step $h$.\\
\textit{Output:} $\bv_r^{n+1}$ in TT format.
\begin{enumerate}
\item \textbf{Forward sweep} ($k = 1,\ldots,d$): integrate the
    $k$-th component ODE
    \begin{equation*}
        \dot{C}^{(k)}(t)
        = \bigl\langle \bF(t,\bv_r(t)),\,
          P_k(\,\cdot\,) \bigr\rangle
    \end{equation*}
    from $t^n$ to $t^n + h$; left-orthogonalise via QR.
\item \textbf{Correction sweep} ($k = d-1,\ldots,1$): apply the
    correction steps $-P_{k,k+1}$ symmetrically.
\item \textbf{Return} the updated TT tensor $\bv_r^{n+1}$.
\end{enumerate}

\begin{theorem}[Error bound in Hilbert spaces;
    {\cite[Theorem~4.1]{LubichOseledetsVandereycken2015}}]
\label{thm:app:splitting_error}
Let $\VDnorm$ be a Hilbert tensor space and $\bF$ Lipschitz with
constant $L$. Then there exists $C > 0$ (depending on $L$, the
curvature of $\Tucker{\fr}{\VD}$, and the initial data) such that
\begin{equation}
\label{eq:app:splitting_bound}
    \|\bu(t^n) - \bv_r^n\|
    \leq Ch + O\!\bigl(\delta_{\fr}(\bu)\bigr),
\end{equation}
where $\delta_{\fr}(\bu)
= \inf_{\bw \in \Tucker{\leq\fr}{\VD}} \|\bu - \bw\|$ is the
best-approximation error.
\end{theorem}

The curvature term in $C$ is the quantity whose Banach-space
analogue appears in Open Problem~\ref{prob:integrator}. For general
tree-based formats, the splitting integrator
of~\cite{CerutiLubichWalach2020} applies a level-by-level sweep along
$\TD$ at cost $O(\dep{\TD} \cdot dr^3)$ per time step.

%------------------------------------------------------------
\section{DMRG and the Tucker manifold}
\label{app:alg:dmrg}
%------------------------------------------------------------

\subsection*{DMRG as ALS on the TT manifold}

The \emph{Density Matrix Renormalisation Group} (DMRG), introduced by
White~\cite{White1992}, is the most widely used algorithm for
ground-state computation in quantum many-body systems. In its modern
formulation~\cite{HoltzRohwedderSchneider2012}, DMRG is equivalent to
ALS on the TT manifold applied to the eigenvalue problem
$H\bu = \lambda\bu$: each microiteration updates one TT component
while fixing all others, producing an effective eigenvalue problem
of size $r^2 n$.

\begin{remark}[Gauge conditions and the Tucker bundle]
\label{rem:app:dmrg}
The \emph{left-orthogonal gauge} in DMRG requires each component tensor
$C^{(k)} \in \RR^{r_{k-1} \times n_k \times r_k}$ to satisfy
\begin{equation*}
    \sum_{i_k=1}^{n_k}
    \bigl(C^{(k)}\bigr)_{:,i_k,:}^{\!\top}
    \bigl(C^{(k)}\bigr)_{:,i_k,:}
    = I_{r_k}.
\end{equation*}
This is the computational analogue of fixing the complement $W_k$ at
each tree node to be the orthogonal complement of $\spann(C^{(k)})$,
giving a canonical trivialisation of the Tucker bundle
(Section~\ref{sec:tucker_manifold:bundle}).
The gauge group $\prod_k \GL{r_k}{\RR}$ of
Remark~\ref{rem:gauge_freedom} is quotiented out by the QR
orthogonalisation step.
\end{remark}

\subsection*{Two-site DMRG and rank adaptation}

One-site DMRG cannot change the bond dimensions $r_k$. \emph{Two-site
DMRG} optimises two adjacent components simultaneously, forming an
effective tensor of rank $r_k \cdot r_{k+1}$ which is then truncated
by SVD to obtain updated bond dimensions. Geometrically, this
corresponds to moving between strata of the stratification
$\VD = \bigsqcup_\fr \Tucker{\fr}{\VD}$
(Theorem~\ref{thm:tucker_decomposition_full}).

%------------------------------------------------------------
\section{Retraction-based optimisation on tensor manifolds}
\label{app:alg:retraction}
%------------------------------------------------------------

\subsection*{Riemannian gradient methods}

In the Hilbert case, the Tucker manifold is a Riemannian submanifold
of $\VDnorm$. Riemannian gradient descent~\cite{UschmajewVandereycken2020}
for minimising $f: \VDnorm \to \RR$ over $\Tucker{\fr}{\VD}$ reads:
\begin{equation}
\label{eq:app:riem_gradient}
    \bv_{k+1}
    = \mathrm{Ret}_{\bv_k}\!\bigl(
        -\alpha_k\,\Pmetric{\bv_k}(\nabla f(\bv_k))
    \bigr),
\end{equation}
where $\Pmetric{\bv}(\nabla f(\bv))$ is the Riemannian gradient and
$\mathrm{Ret}_\bv: \ZD{\bv} \to \Tucker{\fr}{\VD}$ is a retraction.

\begin{definition}[Retraction]
\label{def:app:retraction}
A \emph{retraction} on a Banach manifold $\mathbb{M}$ is a
$\mathcal{C}^1$-map $\mathrm{Ret}: T\mathbb{M} \to \mathbb{M}$
such that, for each $m \in \mathbb{M}$:
\begin{enumerate}
\item[\textup{(i)}] $\mathrm{Ret}_m(0_m) = m$,
\item[\textup{(ii)}] $D(\mathrm{Ret}_m)(0_m) = \Id{T_m\mathbb{M}}$.
\end{enumerate}
\end{definition}

\begin{example}[Tucker retraction via HOSVD truncation]
\label{ex:app:tucker_retraction}
A practical retraction on $\Tucker{\fr}{\VD}$ is:
\begin{equation*}
    \mathrm{Ret}_{\bv}(\dot\bv)
    = \mathrm{HOSVD}_\fr(\bv + \dot\bv),
\end{equation*}
i.e., apply Algorithm~D.1 to $\bv + \dot\bv$ and truncate to rank
$\fr$. This satisfies (i) and (ii) with quasi-optimality constant
$\sqrt{d}$ (Theorem~\ref{thm:app:hosvd}), and remains valid in Banach
spaces as long as the best-approximation step is solvable ---
guaranteed by Theorem~\ref{thm:best_approx_tucker}.
\end{example}

\subsection*{The Banach case: duality-mapping gradient}

In a non-Hilbert Banach tensor space $\VDnorm$, the Riemannian gradient
is replaced by the \emph{duality-mapping gradient}: the unique element
$\nabla_B f(\bv) \in \ZD{\bv}$ satisfying
\begin{equation}
\label{eq:app:banach_gradient}
    \bigl\langle \nabla_B f(\bv),\,
    J\!\bigl(\bv - \bu\bigr)\bigr\rangle
    = Df(\bv)[\bv - \bu]
    \qquad \forall\,\bu \in \Tucker{\fr}{\VD},
\end{equation}
where $J: \VDnorm \to (\VDnorm)^*$ is the normalised duality mapping
(Appendix~\ref{app:functional}, \S\ref{app:functional:index}).
The iteration~\eqref{eq:app:riem_gradient} becomes a generalised
gradient step using the generalised projection $\Pi_{\bv}$ of
Appendix~\ref{app:functional} in place of $\Pmetric{\bv}$. Convergence
analysis in this setting is related to Open Problem~\ref{prob:finsler}.

%------------------------------------------------------------
\section{Historical notes}
\label{app:alg:history}
%------------------------------------------------------------

\paragraph{Tucker decomposition and HOSVD.}
The Tucker decomposition~\cite{Tucker1966} was introduced as a
psychometric tool and has been used in data analysis since the 1970s.
The HOSVD as an algorithm for best Tucker approximations was introduced
by De Lathauwer, De Moor, and Vandewalle~\cite{DeLathauwer2000}; the
quasi-optimality bound~\eqref{eq:app:hosvd_bound} is from the same
paper. A comprehensive survey of tensor decompositions is given in
Kolda and Bader~\cite{KoldaBader2009}.

\paragraph{Tensor Train and DMRG.}
The TT decomposition was introduced independently by Vidal~\cite{Vidal2003}
(matrix product states, MPS) and by Oseledets and
Tyrtyshnikov~\cite{Oseledets2011}. The DMRG algorithm of
White~\cite{White1992} predates both and was recognised
in~\cite{HoltzRohwedderSchneider2012} as ALS on the TT manifold.
The TT-SVD of~\cite{Oseledets2011} provides quasi-optimal TT
approximations analogous to the HOSVD for Tucker.

\paragraph{Projector-splitting integrators.}
The projector-splitting integrator for TT formats was introduced by
Lubich, Oseledets, and Vandereycken~\cite{LubichOseledetsVandereycken2015},
building on the matrix case of Koch and Lubich~\cite{KochLubich2007}
(MCTDH). The extension to general tree-based formats is
in~\cite{CerutiLubichWalach2020}. Convergence in the Banach setting
is Open Problem~\ref{prob:integrator}.

\paragraph{Riemannian optimisation.}
Riemannian gradient methods on Tucker and TT manifolds were
systematically developed by Uschmajew and
Vandereycken~\cite{UschmajewVandereycken2020}, using the Riemannian
structure in a way that parallels
Chapters~\ref{chap:tucker_manifold}--\ref{chap:TB_manifold} but
restricted to the Hilbert case. The Banach generalisation is related
to Open Problem~\ref{prob:finsler}.

%============================================================
%  app_notation_index.tex  —  Index of Notation
%  Tensor Formats in Banach Spaces
%  Antonio Falcó, CEU Universities
%============================================================

\chapter{Index of Notation}
\label{app:notation_index}

All macros are defined in \texttt{notation.tex}. Entries are grouped
thematically and listed in the order they first appear in the text.

%------------------------------------------------------------
\section*{Sets and fields}
%------------------------------------------------------------

\begin{tabular}{@{}lll@{}}
$\KK$ & field ($\RR$ or $\CC$) & \\
$\RR$, $\CC$ & real and complex numbers & \\
$\NN$, $\NNz$, $\ZZ$ & positive integers, non-negative integers, integers & \\
$D = \{1,\ldots,d\}$ & dimension index set & Ch.~\ref{chap:algtensor} \\
$d = |D|$ & number of dimensions & Ch.~\ref{chap:algtensor} \\
$\alpha^c = D \setminus \alpha$ & complement of $\alpha$ in $D$ & Ch.~\ref{chap:algtensor} \\
$2^D$ & power set of $D$ & Ch.~\ref{chap:trees} \\
$\#(S)$ & cardinality of set $S$ & \\
$\overline{S}$ & closure of $S$ & \\
\end{tabular}

%------------------------------------------------------------
\section*{Banach spaces and linear maps}
%------------------------------------------------------------

\begin{tabular}{@{}lll@{}}
$(X, \|\cdot\|)$ & normed (Banach) space & Ch.~\ref{chap:banach} \\
$X^* = \cL(X,\KK)$ & topological dual of $X$ & Ch.~\ref{chap:banach} \\
$X' $ & algebraic dual of $X$ & Ch.~\ref{chap:banach} \\
$\langle x, f \rangle$ & duality pairing, $f \in X^*$, $x \in X$ & \\
$\cL(X,Y)$ & bounded linear maps $X \to Y$ & Ch.~\ref{chap:banach} \\
$\cL(X)$ & $\cL(X,X)$ & \\
$\|A\|_{Y \leftarrow X}$ & operator norm of $A \in \cL(X,Y)$ & Ch.~\ref{chap:banach} \\
$\GL{X}$ & bounded linear automorphisms of $X$ & Ch.~\ref{chap:banach} \\
$\Id{X}$ & identity map on $X$ & \\
$P_{U \oplus W}$ & projection onto $U$ along $W$ & Ch.~\ref{chap:banach} \\
$J: X \to 2^{X^*}$ & normalised duality mapping & Ch.~\ref{chap:DF} \\
$\phi(x,y)$ & Lyapunov functional $\|x\|^2 - 2\langle x,J(y)\rangle + \|y\|^2$ & Ch.~\ref{chap:DF} \\
$\cL_{(U,W)}(X)$ & nilpotent algebra associated to splitting & Ch.~\ref{chap:banach} \\
\end{tabular}

%------------------------------------------------------------
\section*{Grassmannians and subspaces}
%------------------------------------------------------------

\begin{tabular}{@{}lll@{}}
$\GrassB{X}$ & Grassmannian of split subspaces of $X$ & Ch.~\ref{chap:banach} \\
$\Grass{r}{X}$ & $r$-dimensional split subspaces of $X$ & Ch.~\ref{chap:banach} \\
$\GrassB{W,X}$ & split subspaces with complement $W$ & Ch.~\ref{chap:banach} \\
$\Psi_{U \oplus W}$ & Grassmann chart map & Ch.~\ref{chap:banach} \\
$T_U \Grass{r}{X} \cong \cL(U,W)$ & tangent space to Grassmannian & Ch.~\ref{chap:banach} \\
\end{tabular}

%------------------------------------------------------------
\section*{Component spaces and tensor products}
%------------------------------------------------------------

\begin{tabular}{@{}lll@{}}
$V_j$, $V_\alpha$ & component spaces ($j \in D$, $\alpha \subset D$) & Ch.~\ref{chap:algtensor} \\
$\|\cdot\|_j$, $\|\cdot\|_\alpha$ & norms on $V_j$, $V_\alpha$ & \\
$\VD = \bigotimes_{j \in D} V_j$ & algebraic tensor product & Ch.~\ref{chap:algtensor} \\
$\mathbf{V}_\alpha = \bigotimes_{j \in \alpha} V_j$ & partial algebraic tensor product & Ch.~\ref{chap:algtensor} \\
$v \otimes_a w$ & algebraic tensor product of vectors & Ch.~\ref{chap:algtensor} \\
$\bv, \bu, \bw$ & elements of $\VD$ (tensors) & \\
$\VDnorm$ & completion of $\VD$ under $\|\cdot\|_D$ & Ch.~\ref{chap:algtensor} \\
$\|\cdot\|_D$ & norm on $\VD$ (or $\VDnorm$) & \\
$\|\bv\|_\varepsilon$ & injective (spatial) norm & Ch.~\ref{chap:minsubspaces} \\
$\|\bv\|_\pi$ & projective (nuclear) norm & Ch.~\ref{chap:algtensor} \\
\condINJ & condition: $\|\cdot\|_D \geq \|\cdot\|_\varepsilon$ & Ch.~\ref{chap:minsubspaces} \\
\condCROSS & condition: $\|\cdot\|_D$ is a crossnorm & Ch.~\ref{chap:algtensor} \\
$\Mmat{\alpha}(\bv)$ & $\alpha$-matricisation (unfolding) of $\bv$ & Ch.~\ref{chap:minsubspaces} \\
\end{tabular}

%------------------------------------------------------------
\section*{Minimal subspaces and Tucker format}
%------------------------------------------------------------

\begin{tabular}{@{}lll@{}}
$\Umin{\alpha}{\bv}$ & minimal subspace of $\bv$ in $V_\alpha$ & Ch.~\ref{chap:minsubspaces} \\
$r_\alpha(\bv) = \dim\Umin{\alpha}{\bv}$ & $\alpha$-rank of $\bv$ & Ch.~\ref{chap:minsubspaces} \\
$\fr = (r_\alpha)_{\alpha \in D}$ & Tucker rank tuple & Ch.~\ref{chap:minsubspaces} \\
$\AD{\VD}$ & admissible rank tuples for $\VD$ & Ch.~\ref{chap:minsubspaces} \\
$\Tucker{\fr}{\VD}$ & Tucker tensors with fixed rank $\fr$ & Ch.~\ref{chap:minsubspaces} \\
$\Tucker{\leq\fr}{\VD}$ & Tucker tensors with rank $\leq \fr$ & Ch.~\ref{chap:minsubspaces} \\
$\ZD{\bv}$ & tangent subspace at $\bv \in \Tucker{\fr}{\VD}$ & Ch.~\ref{chap:tucker_manifold} \\
$\Emodel{D}$ & model Banach space of Tucker manifold & Ch.~\ref{chap:tucker_manifold} \\
$\xi_{\bv}$ & local chart map at $\bv$ & Ch.~\ref{chap:tucker_manifold} \\
$\mathcal{U}(\bv)$ & chart domain at $\bv$ & Ch.~\ref{chap:tucker_manifold} \\
\end{tabular}

%------------------------------------------------------------
\section*{Trees and tree-based formats}
%------------------------------------------------------------

\begin{tabular}{@{}lll@{}}
$\TD$ & dimension partition tree on $D$ & Ch.~\ref{chap:trees} \\
$S(\alpha)$ & sons of vertex $\alpha$ in $\TD$ & Ch.~\ref{chap:trees} \\
$\mathcal{L}(\TD)$ & leaves of $\TD$ & Ch.~\ref{chap:trees} \\
$\mathrm{level}(\alpha)$ & level of vertex $\alpha$ & Ch.~\ref{chap:trees} \\
$\mathrm{depth}(\TD)$ & depth of the tree $\TD$ & Ch.~\ref{chap:trees} \\
$\mathcal{P}_k(\TD)$ & level-$k$ partition of $D$ & Ch.~\ref{chap:trees} \\
$\TDtuck$ & Tucker tree (depth 1) & Ch.~\ref{chap:trees} \\
$\TDtt$ & Tensor Train tree (linear chain) & Ch.~\ref{chap:trees} \\
$\TDht$ & Hierarchical Tucker tree (binary) & Ch.~\ref{chap:trees} \\
$\fr = (\mathfrak{r}_\alpha)_{\alpha \in \TD \setminus \mathcal{L}(\TD)}$ & tree-based rank tuple & Ch.~\ref{chap:algTBT} \\
$\ADTD$ & admissible tree-based rank tuples & Ch.~\ref{chap:algTBT} \\
$\FTv = \mathcal{FT}_{\fr}(\VD, \TD)$ & tree-based tensors with fixed rank $\fr$ & Ch.~\ref{chap:algTBT} \\
$\mathcal{FT}_{\leq\fr}(\VD,\TD)$ & tree-based tensors with rank $\leq \fr$ & Ch.~\ref{chap:topTBT} \\
$C^{(\alpha)}$ & transfer tensor at vertex $\alpha$ & Ch.~\ref{chap:algTBT} \\
\end{tabular}

%------------------------------------------------------------
\section*{Differential geometry}
%------------------------------------------------------------

\begin{tabular}{@{}lll@{}}
$T_m \mathbb{M}$ & tangent space at $m \in \mathbb{M}$ & Ch.~\ref{chap:banach} \\
$Tf(m)$ & tangent map (derivative) of $f$ at $m$ & Ch.~\ref{chap:banach} \\
$Df(m)$ & Fréchet derivative of $f$ at $m$ & Ch.~\ref{chap:banach} \\
$\mathcal{C}^p$ & $p$ times continuously Fréchet differentiable & \\
$\mathcal{C}^\omega$ & analytic (real or complex) & \\
\end{tabular}

%------------------------------------------------------------
\section*{Dirac--Frenkel variational principle}
%------------------------------------------------------------

\begin{tabular}{@{}lll@{}}
$\bF(t,\bv)$ & vector field (right-hand side of ODE) & Ch.~\ref{chap:DF} \\
$\bv_r(t)$ & solution curve on $\Tucker{\fr}{\VD}$ & Ch.~\ref{chap:DF} \\
$\Pmetric{\bv}$ & metric projection onto $\ZD{\bv}$ & Ch.~\ref{chap:DF} \\
$\Pi_{\bv}$ & generalised projection onto $\ZD{\bv}$ & Ch.~\ref{chap:DF} \\
$\argmin$ & argument of the minimum & Ch.~\ref{chap:DF} \\
\end{tabular}

%------------------------------------------------------------
\section*{Common abbreviations}
%------------------------------------------------------------

\begin{tabular}{@{}ll@{}}
TT & Tensor Train \\
HT & Hierarchical Tucker \\
TBT & Tree-Based Tensor \\
ODE & Ordinary Differential Equation \\
PDE & Partial Differential Equation \\
MCTDH & Multiconfiguration Time-Dependent Hartree \\
\end{tabular}




%------------------------------------------------------------
%  BACK MATTER
%------------------------------------------------------------
\backmatter

\bibliographystyle{spbasic}
\bibliography{bibliography/monograph}

\printindex

\end{document}
